% Lesson 23 — Vocabulary: Words and expressions related to taking part and managing time in community/workplace/school activities — Anglais Tle A — Cours signé OnBuch+
% Programme officiel MINESEC. Compiler avec : tectonic EN23-vocabulary-words-and-expressions-related-to-taking-part.tex
\newcommand{\DOCMATIERE}{Anglais}
\newcommand{\DOCNIVEAU}{Tle A}
\newcommand{\DOCMODULE}{Chapter 2 : Economic Life and Occupations}
\newcommand{\DOCLECON}{EN23}
\newcommand{\DOCTITRE}{Lesson 23 — Vocabulary: Words and expressions related to taking part and managing time in community/workplace/school activities}
% OnBuch+ — préambule « cours signés » ENRICHI (compilé avec tectonic / XeLaTeX)
% Système visuel « Pop » d'OnBuch+ : fond crème (jamais blanc), bordures encre
% épaisses, ombres « dures » décalées, titres Archivo Black en capitales.
% Version MATHÉMATIQUES : graphes (pgfplots/tikz), boîtes pédagogiques
% (définition, propriété, méthode, attention, le savais-tu, exercice résolu,
% exercice, TP, bilan), page de garde et signature OnBuch+ ancrée en bas.
%
% Macros attendues dans main.tex AVANT \input{preamble} :
%   \DOCMATIERE  \DOCNIVEAU  \DOCMODULE  \DOCLECON (numéro)  \DOCTITRE
\documentclass[11pt]{article}

\usepackage{fontspec}
\usepackage[english,french]{babel} % français = langue principale ; l'anglais via \eng{..} / textebox
\usepackage[a4paper,margin=2cm,top=3.4cm,bottom=2.8cm,headheight=1.4cm,footskip=1.3cm]{geometry}
\usepackage{amsmath,amssymb,amsfonts}
\usepackage{mathtools} % \xmapsto, \coloneqq... souvent utilisés par les modèles
\usepackage{cancel} % simplifications barrées (bilans de forces)
% \abs{z} (module, valeur absolue) et \norm{u} (norme), utilisés par les modèles
\providecommand{\abs}{}\let\abs\relax\DeclarePairedDelimiter{\abs}{\lvert}{\rvert}
\providecommand{\norm}{}\let\norm\relax\DeclarePairedDelimiter{\norm}{\lVert}{\rVert}
\usepackage[table]{xcolor} % [table] : fournit aussi \rowcolors (alternance de lignes)
\usepackage{pagecolor}
\usepackage{tikz}
\usetikzlibrary{arrows.meta,positioning,calc,shapes.geometric,shapes.misc,shapes.arrows,decorations.pathmorphing,decorations.pathreplacing,decorations.markings,patterns,shadows,mindmap,trees,fit,backgrounds,matrix,intersections,angles,quotes}
\usepackage{pgfplots}
\usepackage[version=4]{mhchem} % équations nucléaires \ce{^{235}_{92}U}
\usepackage[european,siunitx]{circuitikz} % schémas électriques
\pgfplotsset{compat=1.18}
\usepgfplotslibrary{fillbetween,groupplots} % aires entre courbes (calcul intégral)
\usepackage{enumitem}
\usepackage{fancyhdr}
% [most] charge toutes les bibliothèques tcolorbox (listings, minted, raster,
% documentation...) et gonfle inutilement la mémoire TeX sur un document long
% (~300 boîtes) : on ne charge que ce qui est réellement utilisé.
\usepackage[skins,breakable]{tcolorbox}
\usepackage{graphicx}
\usepackage{colortbl}
\usepackage{booktabs}
\usepackage{tabularx}
\usepackage{diagbox} % case d'en-tête diagonale des tableaux à double entrée
% Colonnes extensibles centrée / gauche / droite (C, L, R), souvent utilisées par les modèles
\newcolumntype{C}{>{\centering\arraybackslash}X}
\newcolumntype{L}{>{\raggedright\arraybackslash}X}
\newcolumntype{R}{>{\raggedleft\arraybackslash}X}
\usepackage{array}
\usepackage{multicol}
\usepackage{multirow}
\usepackage{siunitx}
% parse-numbers=false : accepte aussi \SI{2,5 \times 10^{-3}}{...}
\sisetup{locale=FR,output-decimal-marker={,},group-minimum-digits=4,parse-numbers=false}
\DeclareSIUnit\ohm{\ensuremath{\symup{\Omega}}} % U+2126 absent de la police
\DeclareSIUnit\division{div} % réglages d'oscilloscope (ms/div, V/div)
\DeclareSIUnit\tr{tr} % tours (vitesse de rotation, ex. \SI{1500}{\tr\per\minute})
\DeclareSIUnit\molar{M} % concentration molaire (ex. \SI{3}{\molar}, \SI{1}{\micro\molar})
\DeclareSIUnit\an{an} % année (le modèle confond parfois avec \year, un compteur TeX) — ex. \SI{5}{\kilo\meter\per\an}
\DeclareSIUnit\FCFA{FCFA} % franc CFA (ex. \SI{500}{\FCFA\per\kilo\gram})
\DeclareSIUnit\degreeBrix{\textdegree Brix} % teneur en sucre d'un sirop/jus (réfractométrie)
\DeclareSIUnit\month{mois} % le modèle utilise parfois \month, un compteur TeX, comme unité de durée
\usepackage{qrcode}
% \degree : commande fréquemment utilisée par les modèles pour un angle ou une
% température en dehors de \SI{}{} (non fournie par les paquets chargés).
\providecommand{\degree}{^{\circ}}
% \qed (fin de démonstration), utilisé par les modèles sans amsthm
\providecommand{\qed}{\hfill$\square$}
% \barre : double barre des valeurs interdites dans un tableau de variations (tkz-tab)
\providecommand{\barre}{\ensuremath{\|}}
% environnement proof (amsthm non chargé) utilisé par les modèles
\ifdefined\proof\else\newenvironment{proof}[1][Démonstration]{\par\noindent\textit{#1.}\ }{\qed\par}\fi
\usepackage{lastpage}
\usepackage{hyperref}
\hypersetup{hidelinks}

% ============================================================
% Palette « Pop » OnBuch+ — mêmes hex que la classe P (plus_ui.dart)
% ============================================================
\definecolor{poporange}{HTML}{F59321}
\definecolor{popdark}{HTML}{E07A0C}
\definecolor{popcream}{HTML}{FFF6E8}
\definecolor{popink}{HTML}{1C1714}
\definecolor{poppurple}{HTML}{7A5AE0}
\definecolor{popgreen}{HTML}{2E9E6A}
\definecolor{poppink}{HTML}{E0457B}
\definecolor{popblue}{HTML}{2F7DE1}
\definecolor{popgold}{HTML}{C98A12}
\definecolor{popmuted}{HTML}{8A7F76}
\definecolor{popline}{HTML}{EADFD2}
\definecolor{popgray}{HTML}{8A7F76}
\colorlet{popgrey}{popgray}
% Alias tolérants (le modèle nomme parfois les couleurs différemment)
\colorlet{popwhite}{white}
\colorlet{popblack}{popink}
\colorlet{popred}{poppink}
\colorlet{popyellow}{popgold}
% Teintes claires (fonds de tableaux/encadrés) — le modèle nomme parfois
% les couleurs avec un suffixe L (ex. popblueL) sans les avoir définies.
\colorlet{poporangeL}{poporange!20}
\colorlet{popdarkL}{popdark!20}
\colorlet{poppurpleL}{poppurple!20}
\colorlet{popgreenL}{popgreen!20}
\colorlet{poppinkL}{poppink!20}
\colorlet{popblueL}{popblue!20}
\colorlet{popgoldL}{popgold!20}
\colorlet{popredL}{popred!20}
\colorlet{popgrayL}{popgray!20}
% Teintes claires pour fonds de tableaux / zones
\colorlet{popblueL}{popblue!10!white}
\colorlet{poporangeL}{poporange!14!white}
\colorlet{popgreenL}{popgreen!12!white}
\colorlet{poppurpleL}{poppurple!12!white}
\colorlet{poppinkL}{poppink!12!white}
\colorlet{popgoldL}{popgold!14!white}

\pagecolor{popcream}

% ============================================================
% Typographie
% ============================================================
\defaultfontfeatures{Ligatures=TeX}
\setmainfont{PlusJakartaSans-Medium.ttf}[Path=../../../fonts/,BoldFont=PlusJakartaSans-Bold.ttf]
% Symboles, API et emojis absents de la police principale : repli sur DejaVu Sans (caractères actifs)
\newfontface\symfont{DejaVuSans.ttf}[Path=../../../fonts/]
\makeatletter
\newcommand\symactive[1]{\catcode#1=13 \begingroup\lccode`\~=#1 \lowercase{\endgroup\def~}{{\symfont\char#1}}}
\newcommand\symblank[1]{\catcode#1=13 \begingroup\lccode`\~=#1 \lowercase{\endgroup\def~}{}}
\newcommand\symhyph[1]{\catcode#1=13 \begingroup\lccode`\~=#1 \lowercase{\endgroup\def~}{\mbox{-}}}
\newcount\symc
\newcommand\symrange[2]{\symc=#1 \@whilenum\symc<#2 \do{\expandafter\symactive\expandafter{\the\symc}\advance\symc by 1 }}
\newcommand\symblankrange[2]{\symc=#1 \@whilenum\symc<#2 \do{\expandafter\symblank\expandafter{\the\symc}\advance\symc by 1 }}
\makeatother
\symrange{"0250}{"02FF}\symactive{"2713}\symactive{"2714}\symactive{"2717}\symactive{"2718}\symactive{"2610}\symactive{"2611}\symactive{"2612}
\symhyph{"2011}\symhyph{"2E1A}\symblankrange{"1F300}{"1FAFF}\symblank{"FE0F}
% Maths en sans-serif (Fira Math) avec chiffres et lettres droites de Plus
% Jakarta, pour que les formules se fondent dans le texte.
\usepackage[math-style=ISO,bold-style=ISO]{unicode-math}
\setmathfont{FiraMath-Regular.otf}[Path=../../../fonts/]
\setmathfont{PlusJakartaSans-Medium.ttf}[Path=../../../fonts/,range={up/num,up/{Latin,latin},\mathup}]
\usepackage{icomma} % virgule décimale sans espace en mode math
% Caractères mathématiques Unicode absents des polices de texte (Plus Jakarta,
% Archivo Black) : rendus par leur équivalent mathématique, en texte comme en
% formule — sinon ils s'affichent en carré vide (ex. « Arithmétique dans ℤ »).
\usepackage{newunicodechar}
\newunicodechar{ℝ}{\ensuremath{\mathbb{R}}}
\newunicodechar{ℤ}{\ensuremath{\mathbb{Z}}}
\newunicodechar{ℂ}{\ensuremath{\mathbb{C}}}
\newunicodechar{ℕ}{\ensuremath{\mathbb{N}}}
\newunicodechar{ℚ}{\ensuremath{\mathbb{Q}}}
\newunicodechar{𝔻}{\ensuremath{\mathbb{D}}}
\newunicodechar{α}{\ensuremath{\alpha}}
\newunicodechar{β}{\ensuremath{\beta}}
\newunicodechar{γ}{\ensuremath{\gamma}}
\newunicodechar{δ}{\ensuremath{\delta}}
\newunicodechar{ε}{\ensuremath{\varepsilon}}
\newunicodechar{θ}{\ensuremath{\theta}}
\newunicodechar{λ}{\ensuremath{\lambda}}
\newunicodechar{μ}{\ensuremath{\mu}}
\newunicodechar{σ}{\ensuremath{\sigma}}
\newunicodechar{τ}{\ensuremath{\tau}}
\newunicodechar{φ}{\ensuremath{\varphi}}
\newunicodechar{ψ}{\ensuremath{\psi}}
\newunicodechar{ω}{\ensuremath{\omega}}
\newunicodechar{Σ}{\ensuremath{\Sigma}}
% majuscules produites par \MakeUppercase dans les titres (α-aminés -> Α)
\newunicodechar{Α}{\ensuremath{\alpha}}
\newunicodechar{Β}{\ensuremath{\beta}}
\newunicodechar{Γ}{\ensuremath{\Gamma}}
\newunicodechar{Δ}{\ensuremath{\Delta}}
\newunicodechar{Ε}{\ensuremath{\varepsilon}}
\newunicodechar{Θ}{\ensuremath{\Theta}}
\newunicodechar{Λ}{\ensuremath{\Lambda}}
\newunicodechar{Μ}{\ensuremath{\mu}}
\newunicodechar{Τ}{\ensuremath{\tau}}
\newunicodechar{Φ}{\ensuremath{\Phi}}
\newunicodechar{Ψ}{\ensuremath{\Psi}}
\newunicodechar{Ω}{\ensuremath{\Omega}}
\newunicodechar{↦}{\ensuremath{\mapsto}}
\newunicodechar{∈}{\ensuremath{\in}}
\newunicodechar{∉}{\ensuremath{\notin}}
\newunicodechar{∧}{\ensuremath{\wedge}}
\newunicodechar{⇔}{\ensuremath{\Leftrightarrow}}
\newunicodechar{⟺}{\ensuremath{\Longleftrightarrow}}
\newunicodechar{⇒}{\ensuremath{\Rightarrow}}
\newunicodechar{⟹}{\ensuremath{\Longrightarrow}}
\newunicodechar{∘}{\ensuremath{\circ}}
\newunicodechar{∪}{\ensuremath{\cup}}
\newunicodechar{∩}{\ensuremath{\cap}}
\newunicodechar{≡}{\ensuremath{\equiv}}
\newunicodechar{⊂}{\ensuremath{\subset}}
\newunicodechar{⊥}{\ensuremath{\perp}}
\newunicodechar{∃}{\ensuremath{\exists}}
\newunicodechar{∀}{\ensuremath{\forall}}
\newunicodechar{∛}{\ensuremath{\sqrt[3]{\,}}}
\newunicodechar{∜}{\ensuremath{\sqrt[4]{\,}}}
\newunicodechar{⃗}{\ensuremath{{}^{\to}}}
\newunicodechar{ⁿ}{\ifmmode^{n}\else\textsuperscript{\ensuremath{n}}\fi}
\newunicodechar{ˣ}{\ifmmode^{x}\else\textsuperscript{\ensuremath{x}}\fi}
\newunicodechar{ᵇ}{\ifmmode^{b}\else\textsuperscript{\ensuremath{b}}\fi}
\newunicodechar{ʳ}{\ifmmode^{r}\else\textsuperscript{\ensuremath{r}}\fi}
\newunicodechar{ᵘ}{\ifmmode^{u}\else\textsuperscript{\ensuremath{u}}\fi}
\newunicodechar{ʸ}{\ifmmode^{y}\else\textsuperscript{\ensuremath{y}}\fi}
\newunicodechar{ᵃ}{\ifmmode^{a}\else\textsuperscript{\ensuremath{a}}\fi}
\newunicodechar{ᵉ}{\ifmmode^{e}\else\textsuperscript{\ensuremath{e}}\fi}
\newunicodechar{ᵏ}{\ifmmode^{k}\else\textsuperscript{\ensuremath{k}}\fi}
\newunicodechar{ᵖ}{\ifmmode^{p}\else\textsuperscript{\ensuremath{p}}\fi}
\newunicodechar{ᴺ}{\ifmmode^{N}\else\textsuperscript{\ensuremath{N}}\fi}
\newunicodechar{ᶜ}{\ifmmode^{c}\else\textsuperscript{\ensuremath{c}}\fi}
\newunicodechar{ʰ}{\ifmmode^{h}\else\textsuperscript{\ensuremath{h}}\fi}
\newunicodechar{ᵠ}{\ifmmode^{q}\else\textsuperscript{\ensuremath{q}}\fi}
\newunicodechar{ₙ}{\ifmmode_{n}\else\textsubscript{\ensuremath{n}}\fi}
\newunicodechar{ₐ}{\ifmmode_{a}\else\textsubscript{\ensuremath{a}}\fi}
\newunicodechar{ᵢ}{\ifmmode_{i}\else\textsubscript{\ensuremath{i}}\fi}
\newunicodechar{ᵣ}{\ifmmode_{r}\else\textsubscript{\ensuremath{r}}\fi}
\newunicodechar{ₖ}{\ifmmode_{k}\else\textsubscript{\ensuremath{k}}\fi}
\newunicodechar{ᵦ}{\ifmmode_{\beta}\else\textsubscript{\ensuremath{\beta}}\fi}
\newunicodechar{ₓ}{\ifmmode_{x}\else\textsubscript{\ensuremath{x}}\fi}
\newfontfamily\popdisplay{ArchivoBlack-Regular.ttf}[Path=../../../fonts/]
\newfontfamily\poplogo{SpaceGrotesk-Bold.ttf}[Path=../../../fonts/]
\newfontfamily\popsemi{PlusJakartaSans-SemiBold.ttf}[Path=../../../fonts/]
\newfontfamily\popxbold{PlusJakartaSans-ExtraBold.ttf}[Path=../../../fonts/]
\newfontfamily\popxboldls{PlusJakartaSans-ExtraBold.ttf}[Path=../../../fonts/,LetterSpace=3.0]

\setlist[enumerate]{leftmargin=1.7em,itemsep=5pt,topsep=4pt}
\setlist[itemize]{leftmargin=1.7em,itemsep=4pt,topsep=4pt,label={\color{poporange}\textbullet}}
\setlength{\parindent}{0pt}
\setlength{\parskip}{5pt}
\renewcommand{\arraystretch}{1.25}

% pgfplots : style Pop par défaut
\pgfplotsset{
  every axis/.append style={axis line style={popink,line width=1pt},tick style={popink},
    grid style={popline,line width=0.5pt},label style={font=\small},tick label style={font=\footnotesize}},
  popcurve/.style={poporange,line width=1.8pt,no markers,smooth},
}
% Même style disponible dans un \draw TikZ simple (hors pgfplots)
\tikzset{popcurve/.style={poporange,line width=1.8pt}}
\tikzset{popgrid/.style={popline,very thin}}
\colorlet{popcurve}{poporange} % le nom du style est parfois utilisé comme couleur

% ============================================================
% Pill / squiggle
% ============================================================
\newcommand{\poppill}[3]{%
  \tikz[baseline=-0.55ex]\node[draw=popink,line width=1.2pt,rounded corners=2.6pt,%
    fill=#2,inner xsep=6.5pt,inner ysep=3pt]{%
    \popxboldls\fontsize{7}{7}\selectfont\color{#3}\MakeUppercase{#1}};%
}
\newcommand{\popsquiggle}[1]{%
  \tikz[baseline=-0.4ex]{
    \draw[#1,line width=2.6pt,line cap=round]
      plot [smooth] coordinates {(0,0.09) (0.34,0.01) (0.68,0.21) (1.02,0.01) (1.36,0.10)};
  }%
}
% Mot-clé surligné (marqueur orange sous le texte)
\newcommand{\cle}[1]{{\popxbold\color{popdark}#1}}

% ============================================================
% En-tête / pied
% ============================================================
\pagestyle{fancy}
\fancyhf{}
\renewcommand{\headrulewidth}{0pt}
\renewcommand{\footrulewidth}{0pt}
\lhead{\raisebox{-0.15ex}{\poplogo\fontsize{13}{13}\selectfont\color{popink}OnBuch\color{poporange}+}}
\rhead{\poppill{\DOCMATIERE\ \textbullet\ \DOCNIVEAU\ \textbullet\ Leçon \DOCLECON}{white}{popink}}
\lfoot{\popxbold\fontsize{7.6}{7.6}\selectfont\color{popmuted}\MakeUppercase{Cours signé OnBuch+}\ \textbullet\ {\popsemi\DOCTITRE}}
\rfoot{\tikz[baseline=-0.6ex]\node[rounded corners=8.5pt,fill=popink,minimum height=17pt,minimum width=17pt,inner xsep=5pt,inner sep=0]{\popxbold\fontsize{8}{8}\selectfont\color{popcream}\thepage\,/\,\pageref*{LastPage}};}

% ============================================================
% Titres
% ============================================================
\newcommand{\chaptitre}[2]{%
  \begin{tcolorbox}[enhanced,colback=poporange,colframe=popink,boxrule=2.6pt,arc=11pt,
      drop shadow=popink,left=15pt,right=15pt,top=12pt,bottom=11pt,before skip=3pt,after skip=18pt]
    \poppill{#2}{popink}{popcream}\par\vspace{9pt}
    {\raggedright\hyphenpenalty=10000\exhyphenpenalty=10000\popdisplay\fontsize{22}{24}\selectfont\color{popcream}\MakeUppercase{#1}\par}
  \end{tcolorbox}
}
% Section numérotée : pastille orange avec le numéro + titre Archivo
\newcounter{popsec}
\newcommand{\coursec}[1]{%
  \stepcounter{popsec}\setcounter{popsub}{0}%
  \par\vspace{16pt}\noindent
  \begin{minipage}{\linewidth}
  \tikz[baseline=-0.7ex]\node[draw=popink,line width=1.6pt,fill=poporange,rounded corners=4pt,minimum size=22pt,inner sep=0]{\popdisplay\fontsize{12}{12}\selectfont\color{popcream}\thepopsec};%
  \hspace{8pt}{\popdisplay\fontsize{15}{17}\selectfont\color{popink}\MakeUppercase{#1}}\par
  \vspace{4pt}\noindent{\color{poporange}\rule{36pt}{3.6pt}}
  \end{minipage}\par\nopagebreak\vspace{8pt}
}
\newcounter{popsub}
\newcommand{\courssub}[1]{%
  \stepcounter{popsub}%
  \par\vspace{9pt}\noindent{\popxbold\fontsize{11.5}{13}\selectfont\color{popink}{\color{poporange}\thepopsec.\thepopsub}\hspace{6pt}#1}\par\nopagebreak\vspace{4pt}
}

% ============================================================
% Boîtes pédagogiques — même recette : fond blanc, bordure épaisse colorée,
% ombre dure encre, badge plein. Titre optionnel [#1].
% ============================================================
\tcbset{popbase/.style={breakable,enhanced,colback=white,boxrule=2.3pt,arc=9pt,
  shadow={3pt}{-3pt}{0pt}{popink},left=14pt,right=14pt,top=11pt,bottom=11pt,before skip=11pt,after skip=15pt,
  fontupper=\popsemi\fontsize{10.6}{14.5}\selectfont\color{popink},
  fontlower=\popsemi\fontsize{10.6}{14.5}\selectfont\color{popink}}}
% Coche dessinée (le glyphe ✓ n'existe pas dans les polices du document)
\newcommand{\popcheck}[1]{\tikz[baseline=-0.5ex]\draw[#1,line width=1.6pt,line cap=round,line join=round](-0.13,0.0)--(-0.04,-0.09)--(0.13,0.1);}
\providecommand{\coche}{\popcheck{popgreen}}
\providecommand{\resultat}[1]{\par\smallskip\noindent\fcolorbox{popgreen}{popgreenL}{\parbox{\dimexpr\linewidth-2\fboxsep-2\fboxrule\relax}{\textbf{Résultat.} #1}}\par\smallskip}
\newcommand{\popboxhead}[3]{\poppill{#1}{#2}{white}\ifx\relax#3\relax\else\hspace{6pt}{\popxbold\fontsize{10.6}{12}\selectfont\color{popink}#3}\fi\par\vspace{7pt}}

\newtcolorbox{aretenir}[1][]{popbase,colframe=popgreen,before upper={\popboxhead{À retenir}{popgreen}{#1}}}
% Texte en anglais (document de lecture, dialogue, extrait) : cadre
% d'étude. Un titre optionnel (source, auteur) s'affiche dans l'en-tête.
\newtcolorbox{textebox}[1][]{popbase,colframe=popblue,before upper={\popboxhead{Text}{popblue}{#1}\begin{otherlanguage}{english}},after upper={\end{otherlanguage}}}
% Marques ✓ / ✗ (forme correcte / fautive) souvent utilisées par les modèles
\providecommand{\cross}{\textcolor{poppink}{\ensuremath{\times}}}
\providecommand{\xmark}{\cross}\providecommand{\crossmark}{\cross}\providecommand{\cmark}{\ensuremath{\checkmark}}
% Ligne de réponse / justification à compléter (inventée par les modèles)
\providecommand{\justification}[1]{\par\noindent\textit{Justification~:} \dotfill\par}
% \textipa (package tipa non chargé) : la police ne couvre pas l'API -> simple passage
\providecommand{\textipa}[1]{#1}
% Soulignement (ulem non chargé) : \uline, \ul
\providecommand{\uline}[1]{\underline{#1}}\providecommand{\ul}[1]{\underline{#1}}
% \glossaire{\item ...} inventée par les modèles : liste de glossaire
\providecommand{\glossaire}[1]{\begin{itemize}#1\end{itemize}}
% \alert (beamer) inventée par les modèles : texte mis en évidence
\providecommand{\alert}[1]{\textbf{\textcolor{poppink}{#1}}}
% variantes ASCII d'ordinaux écrites en macro
\providecommand{\iere}{\textsuperscript{re}}\providecommand{\ieme}{\textsuperscript{e}}\providecommand{\ere}{\textsuperscript{re}}\providecommand{\eme}{\textsuperscript{e}}\providecommand{\er}{\textsuperscript{er}}\providecommand{\ie}{\textsuperscript{e}}
% Ordinaux français écrits en macro par les modèles : 1\ière, 3\ième
\newcommand{\ière}{\textsuperscript{re}}\newcommand{\ième}{\textsuperscript{e}}
% \exemple (inventée par les modèles) : étiquette « Exemple : »
\providecommand{\exemple}{\textbf{Exemple~:}\ }
% \square utilisable aussi hors mode math (cases à cocher)
\AtBeginDocument{\let\squareMath\square\renewcommand{\square}{\ensuremath{\squareMath}}}
% Mot ou phrase en anglais à l'intérieur d'un paragraphe français
\newcommand{\eng}[1]{\ifmmode\text{\textit{#1}}\else\begingroup\selectlanguage{english}\itshape #1\endgroup\fi}
\let\esp\eng % alias toléré
% flèches d'intonation inventées par les modèles
\providecommand{\textsearrow}{}\renewcommand{\textsearrow}{\ensuremath{\searrow}}\providecommand{\textnearrow}{}\renewcommand{\textnearrow}{\ensuremath{\nearrow}}
% \fr{..} : mot français (inventé par les modèles)
\providecommand{\fr}[1]{\textit{#1}}
\newtcolorbox{exemplebox}[1][]{popbase,colframe=poporange,before upper={\popboxhead{Exemple}{poporange}{#1}}}
\newtcolorbox{definition}[1][]{popbase,colframe=popblue,before upper={\popboxhead{Définition}{popblue}{#1}}}
\newtcolorbox{propriete}[1][]{popbase,colframe=poppurple,before upper={\popboxhead{Propriété}{poppurple}{#1}}}
\newtcolorbox{methode}[1][]{popbase,colframe=popgold,before upper={\popboxhead{Méthode}{popgold}{#1}}}
\newtcolorbox{attention}[1][]{popbase,colframe=poppink,before upper={\popboxhead{Attention}{poppink}{#1}}}
\newtcolorbox{experience}[1][]{popbase,colframe=popblue,colback=popblueL,before upper={\popboxhead{Expérience}{popblue}{#1}}}
% « Le savais-tu ? » : carte violette pleine (contexte, culture, Cameroun)
\newtcolorbox{savaistu}[1][]{popbase,colback=poppurple,colframe=popink,
  fontupper=\popsemi\fontsize{10.4}{14.2}\selectfont\color{white},
  before upper={\poppill{Le savais-tu\,?}{white}{poppurple}\ifx\relax#1\relax\else\hspace{6pt}{\popxbold\color{white}#1}\fi\par\vspace{7pt}}}
% Exercice résolu : énoncé en haut, \tcblower puis la solution sur fond crème
\newcounter{popexo}\newcounter{popexr}
\newtcolorbox{exoresolu}[1][]{popbase,colframe=poporange,colbacklower=popcream,
  segmentation style={popink,line width=1.2pt,dashed},
  before upper={\stepcounter{popexr}\popboxhead{Exercice résolu \thepopexr}{poporange}{#1}},
  before lower={\poppill{Solution}{popink}{popcream}\par\vspace{6pt}}}
% Exercice (non résolu en place) — la correction est dans la partie Corrigés
\newtcolorbox{exercice}[1][]{popbase,colframe=popink,
  before upper={\stepcounter{popexo}\popboxhead{Exercice \thepopexo}{popink}{#1}}}
% Corrigé
\newtcolorbox{corrige}[1][]{popbase,colframe=popgreen,colback=popgreenL,
  before upper={\popboxhead{Corrigé}{popgreen}{#1}}}
% Objectifs / prérequis / bilan
\newtcolorbox{objectifs}{popbase,colframe=popink,colback=poporangeL,
  before upper={\popboxhead{Objectifs de la leçon}{popink}{}}}
\newtcolorbox{prerequis}{popbase,colframe=popink,colback=popblueL,
  before upper={\popboxhead{Prérequis}{popblue}{}}}
\newtcolorbox{bilan}{popbase,colframe=popink,colback=white,boxrule=2.8pt,
  before upper={\popboxhead{Fiche bilan}{poporange}{L'essentiel à connaître pour l'examen}}}
% Figure encadrée avec légende
\newtcolorbox{popfigure}[1][]{enhanced,breakable=false,colback=white,colframe=popline,boxrule=1.4pt,arc=8pt,
  left=10pt,right=10pt,top=10pt,bottom=8pt,before skip=10pt,after skip=12pt,halign=center,
  lower separated=false,halign lower=center,
  fontlower=\popsemi\fontsize{9}{11.5}\selectfont\color{popmuted},
  #1}
\newcommand{\legende}[1]{\par\vspace{6pt}{\popsemi\fontsize{9}{11.5}\selectfont\color{popmuted}#1}}

% ============================================================
% Signature OnBuch+
% ============================================================
\newcommand{\poplogoinline}[1]{{\poplogo\fontsize{#1}{#1}\selectfont\color{popink}OnBuch\color{poporange}+}}
\newcommand{\signatureonbuch}{%
  \begin{tcolorbox}[enhanced,colback=popink,colframe=popink,boxrule=2.2pt,arc=10pt,
      drop shadow=poporange,left=14pt,right=14pt,top=10pt,bottom=10pt,before skip=0pt,after skip=0pt]
    \tikz[baseline=-0.7ex]\node[circle,fill=popgreen,draw=popcream,line width=1.3pt,minimum size=22pt,inner sep=0]{\popcheck{white}};%
    \hspace{9pt}\begin{minipage}[c]{0.62\linewidth}
      {\popxbold\fontsize{12}{13}\selectfont\color{popcream}Signé\ \textbullet\ {\poplogo OnBuch\color{poporange}+}}\par\vspace{2pt}
      {\raggedright\popsemi\fontsize{8}{10.5}\selectfont\color{popline}\MakeUppercase{Équipe pédagogique OnBuch+}\\\MakeUppercase{Conforme au programme officiel MINESEC}\par}
    \end{minipage}\hfill
    {\poplogo\fontsize{22}{22}\selectfont\color{popcream}OnBuch\color{poporange}+}
  \end{tcolorbox}%
}

% ============================================================
% Page de garde
% #1 titre · #2 matière · #3 niveau · #4 module · #5 n° leçon · #6 durée ·
% #7 description / objectifs (texte ou itemize)
% ============================================================
\newcommand{\pagegarde}[7]{%
  \thispagestyle{empty}
  \newgeometry{a4paper,margin=1.7cm,top=1.6cm,bottom=1.4cm}%
  \begin{tikzpicture}[remember picture,overlay]
    \fill[poporange!18] ([xshift=-3.2cm,yshift=-3.6cm]current page.north east) circle (4.6cm);
    \fill[poppurple!14] ([xshift=2.4cm,yshift=7.6cm]current page.south west) circle (3.8cm);
  \end{tikzpicture}%
  {\poplogo\fontsize{30}{30}\selectfont\color{popink}OnBuch\color{poporange}+}\hfill
  \raisebox{0.6ex}{\poppill{Cours signé \textbullet\ Premium}{popink}{popcream}}\par
  \vspace{1pt}\popsquiggle{poppurple}\par
  \vspace{8pt}
  \begin{tcolorbox}[enhanced,colback=poporange,colframe=popink,boxrule=2.8pt,arc=13pt,
      drop shadow=popink,left=18pt,right=18pt,top=12pt,bottom=13pt,after skip=10pt]
    \poppill{#2 \textbullet\ #3}{popink}{popcream}\hspace{5pt}\poppill{#4}{white}{popink}\par\vspace{8pt}
    {\popdisplay\fontsize{14}{14}\selectfont\color{popink}LEÇON #5}\par\vspace{4pt}
    {\raggedright\hyphenpenalty=10000\exhyphenpenalty=10000\popdisplay\fontsize{25}{27}\selectfont\color{popcream}\MakeUppercase{#1}\par}
    \vspace{8pt}
    {\popxbold\fontsize{9.5}{9.5}\selectfont\color{popink}\MakeUppercase{Durée conseillée : #6}}
  \end{tcolorbox}
  \begin{tcolorbox}[enhanced,colback=white,colframe=popink,boxrule=2.2pt,arc=10pt,
      drop shadow=popink,left=16pt,right=16pt,top=10pt,bottom=10pt,after skip=8pt]
    \poppill{Dans cette leçon}{poppurple}{white}\par\vspace{5pt}
    \popsemi\fontsize{9.3}{12.6}\selectfont\color{popink}#7
  \end{tcolorbox}
  \noindent\begin{minipage}[t]{0.487\linewidth}
    \begin{tcolorbox}[enhanced,colback=popgreen,colframe=popink,boxrule=2.2pt,arc=10pt,
        drop shadow=popink,left=11pt,right=11pt,top=10pt,bottom=10pt,height=3.1cm,valign=center]
      \tcbox[colback=white,colframe=popink,boxrule=1.3pt,arc=5pt,boxsep=0pt,nobeforeafter,
        left=3pt,right=3pt,top=3pt,bottom=3pt,valign=center]{\qrcode[height=1.9cm]{https://play.google.com/store/apps/details?id=com.luvvix.onbuch&hl=fr}}%
      \hspace{8pt}\begin{minipage}[c]{0.5\linewidth}
        {\popdisplay\fontsize{11}{12.5}\selectfont\color{white}\MakeUppercase{Télécharge OnBuch}}\par\vspace{4pt}
        {\raggedright\popsemi\fontsize{8.4}{11}\selectfont\color{white}Gratuit \textbullet\ Google Play\\Tuteur IA, annales, concours, résultats\par}
      \end{minipage}
    \end{tcolorbox}
  \end{minipage}\hfill
  \begin{minipage}[t]{0.487\linewidth}
    \begin{tcolorbox}[enhanced,colback=poppurple,colframe=popink,boxrule=2.2pt,arc=10pt,
        drop shadow=popink,left=11pt,right=11pt,top=10pt,bottom=10pt,height=3.1cm,valign=center]
      \tikz[baseline=-0.7ex]\node[circle,fill=white,draw=popink,line width=1.3pt,minimum size=1.6cm,inner sep=0]{\popdisplay\fontsize{24}{24}\selectfont\color{poppurple}+};%
      \hspace{8pt}\begin{minipage}[c]{0.6\linewidth}
        {\popdisplay\fontsize{11}{12.5}\selectfont\color{white}\MakeUppercase{Passe à OnBuch+}}\par\vspace{4pt}
        {\raggedright\popsemi\fontsize{8.4}{11}\selectfont\color{white}Tous les cours signés, Tuteur IA illimité, fiches PDF — onglet OnBuch+ de l'app\par}
      \end{minipage}
    \end{tcolorbox}
  \end{minipage}
  \vspace{8pt}
  \popcontenu
  \vfill
  \signatureonbuch
  \restoregeometry
  \newpage
}

% Rangée « Ce cours contient » (page de garde)
\newcommand{\poptile}[3]{% #1 couleur #2 gros texte #3 légende
  \begin{tcolorbox}[enhanced,colback=#1,colframe=popink,boxrule=2pt,arc=9pt,shadow={2.5pt}{-2.5pt}{0pt}{popink},
      left=6pt,right=6pt,top=9pt,bottom=9pt,halign=center,valign=center,height=1.75cm,width=\linewidth,nobeforeafter]
    {\popdisplay\fontsize{12.5}{13}\selectfont\color{white}\MakeUppercase{#2}}\par\vspace{3pt}
    {\centering\popsemi\fontsize{7.8}{9.5}\selectfont\color{white}#3\par}
  \end{tcolorbox}}
\newcommand{\popcontenu}{%
  \par\noindent{\popxboldls\fontsize{7.5}{7.5}\selectfont\color{popmuted}CE COURS SIGNÉ CONTIENT}\par\vspace{7pt}
  \noindent\begin{minipage}[t]{0.235\linewidth}\poptile{popblue}{Cours}{complet, illustré, pas à pas}\end{minipage}\hfill
  \begin{minipage}[t]{0.235\linewidth}\poptile{poporange}{Méthodes}{et exercices résolus}\end{minipage}\hfill
  \begin{minipage}[t]{0.235\linewidth}\poptile{poppink}{Exercices}{type examen + corrigés}\end{minipage}\hfill
  \begin{minipage}[t]{0.235\linewidth}\poptile{popgreen}{Fiche bilan}{l'essentiel à retenir}\end{minipage}\par}

% Sommaire
\newtcolorbox{sommaire}{enhanced,colback=white,colframe=popink,boxrule=2.2pt,arc=10pt,
  drop shadow=popink,left=18pt,right=18pt,top=13pt,bottom=13pt,before skip=6pt,after skip=20pt,
  before upper={\poppill{Sommaire}{poppurple}{white}\par\vspace{9pt}\popsemi\fontsize{10.6}{14.5}\selectfont\color{popink}}}

% Signature de fin de cours — poussée tout en bas de la dernière page
\newcommand{\findecours}{%
  \par\vfill\nopagebreak
  \begin{center}\popsquiggle{poporange}\end{center}\vspace{4pt}
  \signatureonbuch
}
\AtBeginDocument{\def\llbracket{\mathopen{[\mkern-3.5mu[}}\def\rrbracket{\mathclose{]\mkern-3.5mu]}}} % intervalles d'entiers (glyphe ⟦ absent de la police)
% Glyphes absents de Fira Math : redéfinis à partir de glyphes disponibles
\AtBeginDocument{%
  \def\setminus{\mathbin{\backslash}}\let\smallsetminus\setminus
  \def\vdots{\mathord{\vbox{\baselineskip3.2pt\lineskiplimit0pt\kern4pt\hbox{.}\hbox{.}\hbox{.}}}}%
  \def\ddots{\mathinner{\mkern1mu\raise6.5pt\hbox{.}\mkern2mu\raise3.6pt\hbox{.}\mkern2mu\raise0.7pt\hbox{.}\mkern1mu}}%
  \def\triangle{\mathord{\scalebox{0.9}{$\symup{\Delta}$}}}%
  \def\nabla{\mathord{\raisebox{\depth}{\scalebox{1}[-1]{$\symup{\Delta}$}}}}%
  \def\checkmark{\popcheck{popdark}}%
  \def\star{\mathbin{\ast}}\def\bigstar{\mathord{\ast}}%
  \def\diamond{\mathbin{\scalebox{0.6}{$\Diamond$}}}%
  \def\bigcup{\mathop{\mathchoice{\raisebox{-0.3em}{\scalebox{1.7}{$\cup$}}}{\scalebox{1.25}{$\cup$}}{\cup}{\cup}}\displaylimits}%
  \def\bigcap{\mathop{\mathchoice{\raisebox{-0.3em}{\scalebox{1.7}{$\cap$}}}{\scalebox{1.25}{$\cap$}}{\cap}{\cap}}\displaylimits}%
  \def\lesseqqgtr{\mathrel{\lessgtr}}\def\bigoplus{\mathop{\oplus}\displaylimits}%
}
\providecommand{\ssi}{si et seulement si} % abréviation fréquente
\AtBeginDocument{\providecommand{\wideparen}{\overparen}} % arc de cercle

%% FIGFIX-BEGIN v4 — correctifs globaux des figures (docs/onbuch-generation/tools/figfix)
\usepackage{adjustbox}\usepackage{etoolbox}
% toute figure TikZ/pgfplots plus large que la ligne est réduite à la largeur disponible
\BeforeBeginEnvironment{tikzpicture}{\begin{adjustbox}{max width=\linewidth}}
\AfterEndEnvironment{tikzpicture}{\end{adjustbox}}
% légendes pgfplots lisibles par-dessus les courbes
\pgfplotsset{every axis/.append style={legend style={fill=white,fill opacity=0.92,text opacity=1,draw=black!15,inner sep=3pt}}}
% étiquettes d'arêtes (syntaxe edge["..."]) avec fond blanc
\tikzset{every edge quotes/.append style={fill=white,inner sep=1.2pt}}
%% FIGFIX-END
\begin{document}

\pagegarde{Lesson 23 — Vocabulary: Words and expressions related to taking part and managing time in community/workplace/school activities}{Anglais}{Tle A}{Chapter 2 : Economic Life and Occupations}{EN23}{2 h}{This vocabulary lesson equips Terminale A students with the words, collocations, idioms and expressions needed to talk about taking part in community, workplace and school activities, and about managing time effectively. Students learn to avoid false friends and to reuse the lexis accurately in speech and writing, in line with the official syllabus (Chapter 2: Economic Life and Occupations).
\vspace{4pt}
\begin{multicols}{2}\small
\begin{itemize}
\item À la fin de cette leçon, je sais identifier et définir le vocabulaire lié à la participation (take part in, join in, volunteer, attend...)
\item Je sais employer le vocabulaire de la gestion du temps (deadline, schedule, postpone, meet a deadline...)
\item Je sais utiliser correctement les collocations courantes de ces deux champs lexicaux
\item Je sais comprendre et employer des expressions idiomatiques liées au temps et à l'engagement
\item Je sais éviter les faux amis et les calques du français (assist, actually, delay, eventually...)
\item Je sais former des mots de la même famille (participate, participation, participant...)
\end{itemize}\end{multicols}}

\begin{sommaire}
\begin{multicols}{2}
\begin{enumerate}
\item Taking part: participation vocabulary
\item Managing time: key vocabulary
\item Collocations and idiomatic expressions
\item False friends and French calques
\item Word families, pronunciation and spelling
\item Using the lexis in context
\item Activité d'intégration
\item Méthodes et exercices résolus
\item Exercices
\item Corrigés des exercices
\item Fiche bilan
\end{enumerate}
\end{multicols}
\end{sommaire}

\begin{objectifs}
\begin{itemize}
\item À la fin de cette leçon, je sais identifier et définir le vocabulaire lié à la participation (take part in, join in, volunteer, attend...)
\item Je sais employer le vocabulaire de la gestion du temps (deadline, schedule, postpone, meet a deadline...)
\item Je sais utiliser correctement les collocations courantes de ces deux champs lexicaux
\item Je sais comprendre et employer des expressions idiomatiques liées au temps et à l'engagement
\item Je sais éviter les faux amis et les calques du français (assist, actually, delay, eventually...)
\item Je sais former des mots de la même famille (participate, participation, participant...)
\item Je sais prononcer et orthographier correctement ce lexique
\item Je sais réemployer ce vocabulaire dans des phrases, des dialogues et un court texte de type Bac
\end{itemize}
\end{objectifs}

\begin{prerequis}
\begin{itemize}
\item Vocabulaire de base de la vie scolaire et des métiers (Première, Chapter: Economic Life)
\item Présent simple et futur (will / be going to) pour parler de projets
\item Prépositions de temps (in, on, at, before, by, until)
\item Notion de registre de langue (formel / informel)
\end{itemize}
\end{prerequis}

\newpage

\coursec{Taking part: participation vocabulary}

\courssub{Introduction : a real-life situation}

\begin{textebox}[Situation : the clean-up campaign in Bamenda]
Imagine: your school in Bamenda is organising a clean-up campaign next Saturday. The principal asks for volunteers, but you must sign up before Friday, manage your revision time, and still attend the youth club meeting on Sunday. How do you say, in correct English, that you will \emph{take part}, \emph{meet the deadline} and \emph{fit everything into your schedule}? Many Cameroonian students say \emph{I will assist to the meeting} — a classic false friend! This lesson gives you the exact words and expressions to participate in activities and manage your time like a native speaker.
\end{textebox}

\textbf{Problématique.} Comment exprimer, en anglais correct, la participation à des activités scolaires, communautaires ou professionnelles ? Quels verbes choisir selon que l'on \emph{adhère} à un club, qu'on \emph{assiste} à une réunion ou qu'on se porte \emph{volontaire} ? Cette première section construit le champ lexical de la participation ; la section suivante traitera de la gestion du temps.

\courssub{Observation : reading the opening text}

Lisons d'abord un court article de presse scolaire. Soulignez mentalement tous les mots et expressions qui expriment la \cle{participation}.

\begin{textebox}[The Buea clean-up campaign — school newspaper article]
Last Saturday, over fifty volunteers took part in a clean-up campaign in Buea. The event was organised by the town council with the support of three secondary schools. Students joined in enthusiastically: some swept the streets around the market, while others collected plastic bottles near the motor park.

The local coordinator, Mrs Eposi, thanked everyone who had signed up before the deadline. “We had more participants than we expected,” she said. “Many students volunteered to come back next month, and two teachers offered to help out with the transport.”

The campaign started at seven in the morning. All the participants attended a short briefing before dividing into teams. Each team had a leader who made sure everyone contributed to the work. By midday, the market area was completely clean.

The mayor of Buea attended the closing ceremony and congratulated the organisers. “Community work teaches young people commitment and team spirit,” he declared. “I encourage every student to get involved in activities like this one.”

The school drama club, which had joined the campaign for the first time, announced that it would run a fund-raising event in December to buy new cleaning equipment. Membership of the club is open to all: interested students can enrol at the secretary's office any day after classes.
\end{textebox}

\textbf{Glossaire.}

\begin{center}
\begin{tabularx}{\linewidth}{|l|l|X|}
\hline
\rowcolor{popblueL} \textbf{English} & \textbf{Nature} & \textbf{Français} \\
\hline
a clean-up campaign & nom & une campagne de nettoyage \\
\hline
the town council & nom & la mairie, le conseil municipal \\
\hline
a briefing & nom & une séance d'information \\
\hline
a fund-raising event & nom & une manifestation de collecte de fonds \\
\hline
equipment & nom (indénombrable) & le matériel, l'équipement \\
\hline
membership & nom & l'adhésion, le fait d'être membre \\
\hline
\end{tabularx}
\end{center}

\textbf{Questions de compréhension.}

\begin{enumerate}
\item How many volunteers took part in the campaign?
\item Who organised the event?
\item What did Mrs Eposi do?
\item What will the drama club do in December?
\item Find in the text three different verbs that express participation.
\end{enumerate}

\courssub{The core verbs of participation}

\begin{definition}[To volunteer]
\eng{To volunteer} means \emph{to offer to do something without being paid or forced} — offrir de faire quelque chose sans être payé ni forcé ; se porter volontaire.\\
\eng{She volunteered to organise the school debate.} « Elle s'est portée volontaire pour organiser le débat scolaire. »
\end{definition}

\begin{propriete}[Constructions de volunteer]
\begin{itemize}
\item \cle{volunteer to do something} : se porter volontaire pour faire quelque chose. \eng{Two students volunteered to carry the chairs.} « Deux élèves se sont portés volontaires pour porter les chaises. »
\item \cle{volunteer for something} : se porter volontaire pour une tâche. \eng{He volunteered for the night shift.} « Il s'est porté volontaire pour l'équipe de nuit. »
\item \cle{a volunteer} (nom) : un bénévole, un volontaire. \eng{The volunteers received certificates.} « Les bénévoles ont reçu des attestations. »
\end{itemize}
Prononciation : « vo-len-TIER » (accent sur la dernière syllabe) pour le verbe et le nom.
\end{propriete}

\begin{propriete}[Les verbes de participation : forme et emploi]
\begin{itemize}
\item \cle{take part in + nom} : participer activement à un événement, une activité. \eng{Over fifty volunteers took part in the campaign.} « Plus de cinquante bénévoles ont participé à la campagne. »
\item \cle{participate in + nom} : synonyme plus soutenu de \eng{take part in}. \eng{All Form Five students participated in the debate.} « Tous les élèves de seconde ont participé au débat. »
\item \cle{join in + nom (ou sans complément)} : se joindre à une activité déjà en cours, sur le moment. \eng{Come on, join in!} « Allez, joins-toi à nous ! »
\item \cle{join + club / group / team} : devenir membre. \eng{She joined the choir in September.} « Elle a rejoint la chorale en septembre. »
\item \cle{attend + nom} (sans préposition) : assister à, être présent à. \eng{He attended the PTA meeting.} « Il a assisté à la réunion des parents d'élèves. »
\item \cle{sign up (for)} : s'inscrire. \eng{You must sign up before Friday.} « Tu dois t'inscrire avant vendredi. »
\item \cle{enrol (in / on)} : s'inscrire officiellement (cours, école). \eng{She enrolled in a computer course.} « Elle s'est inscrite à un cours d'informatique. »
\item \cle{get involved in} : s'impliquer, s'engager dans. \eng{More young people should get involved in community work.} « Plus de jeunes devraient s'impliquer dans le travail communautaire. »
\item \cle{contribute (to)} : contribuer à. \eng{Everyone contributed to the success of the event.} « Chacun a contribué au succès de la manifestation. »
\item \cle{help out} : donner un coup de main. \eng{Can you help out with the decorations?} « Peux-tu donner un coup de main pour les décorations ? »
\item \cle{organise / run an event} : organiser / diriger, mener un événement. \eng{The club runs a quiz night every month.} « Le club organise une soirée quiz chaque mois. »
\end{itemize}
\end{propriete}

\begin{attention}[Faux ami : attend ≠ attendre, assist ≠ assister à]
Ne dites jamais \eng{assist to a meeting} ! Dites \eng{attend a meeting} ou \eng{be present at a meeting}. En anglais, \cle{to assist} signifie \emph{aider} : \eng{She assisted the nurse.} « Elle a aidé l'infirmière. » De même, \eng{to attend to something} signifie \emph{s'occuper de} quelque chose : \eng{Please attend to this customer.} « Occupe-toi de ce client, s'il te plaît. » Enfin, \eng{to attend} ne signifie jamais \emph{attendre} : « attendre » se dit \cle{to wait (for)}.
\end{attention}

\begin{exemplebox}[Exemples traduits]
\begin{itemize}
\item \eng{He joined the drama club last year.} « Il a adhéré au club de théâtre l'an dernier. »
\item \eng{Come on, join in the game!} « Allez, participe au jeu ! »
\item \eng{All students must attend the assembly.} « Tous les élèves doivent assister à la cérémonie. »
\item \eng{She signed up for the clean-up campaign on Tuesday.} « Elle s'est inscrite à la campagne de nettoyage mardi. »
\item \eng{Our class organised a fund-raising event for the school library.} « Notre classe a organisé une collecte de fonds pour la bibliothèque scolaire. »
\item \eng{He always helps out when we run an event.} « Il donne toujours un coup de main quand nous organisons une manifestation. »
\end{itemize}
\end{exemplebox}

\courssub{Join, join in, take part in : the key distinction}

C'est la distinction la plus délicate de cette leçon. Observons :

\begin{itemize}
\item \eng{My brother joined the army in 2020.} → devenir membre d'une institution.
\item \eng{We were singing and she joined in.} → se joindre à une activité en cours.
\item \eng{Fifty students took part in the competition.} → participer activement à un événement organisé.
\end{itemize}

\begin{propriete}[Règle de choix]
\begin{enumerate}
\item \cle{join + organisation} (club, groupe, équipe, parti, armée) = devenir membre. Pas de préposition.
\item \cle{join in + activité} (game, song, discussion, dance) = participer sur le moment, souvent après le début de l'activité. On peut aussi dire \eng{join somebody} : \eng{May I join you?} « Puis-je me joindre à vous ? »
\item \cle{take part in + événement} (competition, campaign, ceremony, meeting, strike) = participer activement, du début à la fin.
\end{enumerate}
\end{propriete}

\begin{center}
\begin{tabularx}{\linewidth}{|X|X|X|}
\hline
\rowcolor{popblueL} \textbf{Français} & \textbf{English} & \textbf{Remarque} \\
\hline
adhérer à un club & join a club & sans préposition \\
\hline
se joindre à une partie & join in the game & activité en cours \\
\hline
participer à une compétition & take part in a competition & participation active \\
\hline
assister à une réunion & attend a meeting & jamais \eng{assist to} \\
\hline
s'inscrire à un cours & sign up for / enrol on a course & \eng{enrol} : plus formel \\
\hline
s'impliquer dans un projet & get involved in a project & engagement durable \\
\hline
\end{tabularx}
\end{center}

\courssub{The people : names of roles}

\begin{definition}[Rôles et participants]
\begin{itemize}
\item \cle{a participant} : un participant. \eng{Each participant received a T-shirt.}
\item \cle{an organiser} : un organisateur. \eng{The organisers planned everything carefully.}
\item \cle{a volunteer} : un bénévole.
\item \cle{a member} : un membre. \eng{She is an active member of the youth club.}
\item \cle{a coordinator} : un coordinateur. \eng{The coordinator divided the volunteers into teams.}
\item \cle{a chairperson} : un(e) président(e) de séance. \eng{The chairperson opened the meeting at nine.}
\item \cle{a helper} : une personne qui aide, un auxiliaire bénévole.
\end{itemize}
Remarque : la plupart de ces noms d'agent se forment avec \cle{-er / -or} : \eng{organise → organiser}, \eng{coordinate → coordinator}, \eng{help → helper}.
\end{definition}

\courssub{Mind map : the vocabulary of taking part}

\begin{popfigure}
\begin{center}
\begin{tikzpicture}[mindmap, grow cyclic, every node/.style={concept}, concept color=popblue!60, text=white,
level 1/.append style={level distance=3.2cm, sibling angle=90},
level 2/.append style={level distance=2.2cm, sibling angle=45, concept color=poporange!70}]
\node[concept, font=\large\bfseries] {TAKING PART}
  child[concept color=popgreen!70] { node {Actions} 
    child { node[align=center] {join\\attend} }
    child { node[align=center] {volunteer\\sign up} }
  }
  child[concept color=poppurple!70] { node {People}
    child { node[align=center] {participant\\member} }
    child { node[align=center] {organiser\\coordinator} }
  }
  child[concept color=poppink!70] { node[align=center] {Places and\\events}
    child { node[align=center] {meeting\\club} }
    child { node[align=center] {campaign\\ceremony} }
  }
  child[concept color=popgold!80] { node {Feelings}
    child { node[align=center] {commitment} }
    child { node[align=center] {team spirit} }
  };
\end{tikzpicture}
\end{center}
\legende{Carte mentale du champ lexical de la participation : quatre familles de mots à mémoriser ensemble.}
\end{popfigure}

\begin{methode}[Comment mémoriser ce vocabulaire]
\begin{enumerate}
\item Classez chaque nouveau mot dans l'une des quatre branches de la carte mentale (action, personne, événement, sentiment).
\item Apprenez le mot avec sa \cle{collocation} : pas seulement \eng{volunteer}, mais \eng{volunteer to organise}.
\item Rédigez une phrase personnelle vraie : \eng{I joined the football team last year.}
\item Vérifiez la préposition : \eng{take part \textbf{in}}, \eng{participate \textbf{in}}, mais \eng{attend} sans préposition.
\end{enumerate}
\end{methode}

\begin{exoresolu}[Choose the correct verb]
Énoncé. Complétez avec \eng{join}, \eng{join in}, \eng{take part in}, \eng{attend} ou \eng{volunteer} (forme correcte) :\\
1. Last year, Amina \dots\ the English club.\\
2. All prefects must \dots\ the Monday assembly.\\
3. We were playing volleyball and the teacher \dots\.\\
4. Over a hundred students \dots\ the inter-school competition.\\
5. Two boys \dots\ to decorate the hall.
\tcblower
Solution détaillée.\\
1. \eng{joined} — devenir membre d'un club, sans préposition. « L'an dernier, Amina a adhéré au club d'anglais. »\\
2. \eng{attend} — assister à une cérémonie, sans préposition. « Tous les préfets doivent assister à la cérémonie du lundi. »\\
3. \eng{joined in} — se joindre à une activité déjà en cours. « Nous jouions au volley et le professeur s'est joint à nous. »\\
4. \eng{took part in} — participation active à un événement (attention au prétérit irrégulier \eng{take → took → taken}). « Plus de cent élèves ont participé à la compétition interscolaire. »\\
5. \eng{volunteered} — \eng{volunteer to do something}. « Deux garçons se sont portés volontaires pour décorer la salle. »
\end{exoresolu}

\begin{exercice}[Your turn]
Traduisez en anglais : 1. « Je me suis inscrite à la campagne de nettoyage. » 2. « Il a assisté à la réunion des parents d'élèves. » 3. « Rejoins-nous dans le jeu ! » 4. « Elle s'est portée volontaire pour organiser la collecte de fonds. »
\end{exercice}

\begin{corrige}[Exercice : Your turn]
1. \eng{I signed up for the clean-up campaign.} 2. \eng{He attended the PTA meeting.} (jamais \eng{assisted to}) 3. \eng{Join in the game!} ou \eng{Join us in the game!} 4. \eng{She volunteered to organise the fund-raising event.}
\end{corrige}

\begin{savaistu}[Volunteering in the English-speaking world]
Au Royaume-Uni et aux États-Unis, le \eng{volunteering} (bénévolat) est très valorisé sur un CV et dans les dossiers de candidature à l'université : les \eng{admissions officers} apprécient les candidats qui ont fait du \eng{community service}. Aux États-Unis, certaines écoles exigent même un nombre d'heures de service communautaire pour obtenir le diplôme. Au Cameroun anglophone, les \eng{community work days} — journées où tout le quartier nettoie les rues et les caniveaux ensemble — sont une tradition forte, héritée du travail collectif villageois. Savoir parler de ces activités en anglais correct est donc utile pour l'examen… et pour votre avenir !
\end{savaistu}

\coursec{Managing time: key vocabulary}

Après avoir exploré le vocabulaire de la participation (\emph{Taking part}) dans la section précédente, nous abordons maintenant une compétence transversale indispensable : la gestion du temps. Que vous soyez élève à Yaoundé, apprenti à Douala ou jeune actif à Buea, savoir \eng{plan}, \eng{organise} et \eng{meet deadlines} conditionne votre réussite. Cette section vous donne les mots exacts pour parler d'horaires, de retards, de reports et d'efficacité.

\courssub{Core vocabulary: schedules, deadlines and appointments}

Commençons par observer un extrait de l'emploi du temps type d'un élève de Terminale A au Lycée de Buea, puis identifions les noms et verbes clés.

\begin{textebox}{Weekly planner extract — Terminale A student}
Monday: Club meeting at 4 p.m. (Debate Club). Prepare arguments.
Tuesday: Free. Revision for Wednesday's test.
Wednesday: \textbf{Deadline} for History essay submission (12 p.m.). Parent-teacher meeting 5 p.m.
Thursday: Football training 4:30 p.m. \textbf{Schedule} dentist appointment.
Friday: \textbf{Sign-up closes} for Saturday clean-up campaign. Pay school fees.
Saturday: Community clean-up campaign 7 a.m. sharp. \textbf{Duration}: 3 hours.
Sunday: Rest. Plan next week.
\end{textebox}

\begin{definition}[Vocabulaire de base : noms et verbes]
Les mots suivants forment le socle lexical pour parler de gestion du temps dans les contextes scolaire, professionnel et communautaire.
\begin{itemize}
    \item \cle{schedule} (n.) /ˈʃɛdjuːl/ : emploi du temps, programme, planning. \eng{My schedule is full this week.} (Mon emploi du temps est chargé cette semaine.)
    \item \cle{schedule} (v.) /ˈʃɛdjuːl/ : programmer, planifier. \eng{We scheduled the meeting for Monday.} (Nous avons programmé la réunion pour lundi.)
    \item \cle{timetable} (n.) /ˈtaɪmteɪbl/ : horaire (transport, cours fixes). Plus rigide que \eng{schedule}. \eng{The school timetable starts at 7:30 a.m.} (L'horaire de l'école commence à 7h30.)
    \item \cle{deadline} (n.) /ˈdɛdlaɪn/ : date limite, échéance ultime.
    \item \cle{appointment} (n.) /əˈpɔɪntmənt/ : rendez-vous (formel, souvent médical ou professionnel). \eng{I have a dentist appointment on Thursday.} (J'ai un rendez-vous chez le dentiste jeudi.)
    \item \cle{meeting} (n.) /ˈmiːtɪŋ/ : réunion, assemblée. \eng{The staff meeting lasts two hours.} (La réunion du personnel dure deux heures.)
    \item \cle{duration} (n.) /djʊəˈreɪʃn/ : durée. \eng{The duration of the exam is three hours.} (La durée de l'examen est de trois heures.)
    \item \cle{delay} (n.) /dɪˈleɪ/ : retard (fait d'être en retard). \eng{The delay was caused by heavy rain.} (Le retard a été causé par la forte pluie.)
    \item \cle{delay} (v.) /dɪˈleɪ/ : retarder, faire attendre. \eng{The rain delayed the match.} (La pluie a retardé le match.)
    \item \cle{postponement} (n.) /pəʊstˈpəʊnmənt/ : report (action de reporter).
\end{itemize}
\end{definition}

\begin{center}
\begin{tabularx}{\linewidth}{|X|X|X|}
\hline
\rowcolor{popblueL} \textbf{English} & \textbf{Français} & \textbf{Contexte / Collocation typique} \\
\hline
\cle{schedule} (n.) & emploi du temps, planning & \eng{a busy schedule}, \eng{weekly schedule} \\
\cle{timetable} & horaire (cours, bus, train) & \eng{school timetable}, \eng{bus timetable} \\
\cle{deadline} & date limite, échéance & \cle{meet a deadline}, \cle{miss a deadline}, \eng{set a deadline} \\
\cle{appointment} & rendez-vous (formel) & \cle{make an appointment}, \cle{keep an appointment}, \cle{cancel an appointment} \\
\cle{meeting} & réunion & \cle{hold a meeting}, \cle{attend a meeting}, \cle{chair a meeting} \\
\cle{duration} & durée & \eng{short duration}, \eng{maximum duration} \\
\cle{delay} (n.) & retard & \eng{flight delay}, \eng{without delay} (sans délai) \\
\cle{postponement} & report & \eng{announce a postponement} \\
\hline
\end{tabularx}
\end{center}

\begin{attention}[Faux ami majeur : \eng{delay} vs \eng{délai}]
En français, « un délai » signifie un temps imparti pour faire quelque chose (ex. \emph{un délai de deux semaines}). En anglais, \cle{delay} signifie \textbf{retard} (le fait d'arriver après l'heure).
\begin{itemize}
    \item \eng{The deadline is Friday.} = La date limite (le délai) est vendredi.
    \item \eng{There was a delay of two hours.} = Il y a eu un retard de deux heures.
    \item \eng{Without delay.} = Sans délai / immédiatement (ici \eng{delay} = attente inutile).
\end{itemize}
Ne dites jamais : \eng{*I have a delay of two weeks to finish.} Dites : \eng{I have a \textbf{deadline} / \textbf{time limit} of two weeks.}
\end{attention}

\courssub{Action verbs: planning, changing and cancelling}

La gestion du temps est une activité dynamique. Voici les verbes d'action pour modifier le cours du temps.

\begin{propriete}[Verbes d'action temporelle : formes et constructions]
\begin{itemize}
    \item \cle{plan} / \cle{organise} / \cle{arrange} : préparer à l'avance. \eng{We planned the event carefully.} / \eng{She organised a study group.} / \eng{He arranged a meeting with the principal.}
    \item \cle{schedule} / \cle{fix} / \cle{set a date} : fixer une date/heure précise. \eng{The exam is scheduled for June.} / \eng{Let's fix a date for the party.} / \eng{They set the date for the wedding.}
    \item \cle{postpone} / \cle{put off} : reporter à plus tard. \textbf{Suivis du gérondif (V-ing) ou de \eng{until} + date.}
    \item \cle{bring forward} : avancer (mettre à une date/heure plus proche).
    \item \cle{cancel} : annuler.
    \item \cle{meet a deadline} : respecter une échéance. \cle{miss a deadline} : manquer une échéance.
\end{itemize}
\end{propriete}

\begin{exemplebox}[Postpone, Put off, Bring forward, Cancel]
\begin{itemize}
    \item \eng{We \textbf{postponed holding} the meeting until next week.} (Nous avons reporté la tenue de la réunion à la semaine prochaine.) $\rightarrow$ \eng{postpone + V-ing}.
    \item \eng{Don't \textbf{put off doing} your homework.} (Ne remets pas tes devoirs à plus tard.) $\rightarrow$ \eng{put off + V-ing}.
    \item \eng{The match was \textbf{postponed until} Saturday.} (Le match a été reporté à samedi.) $\rightarrow$ Passif + \eng{until}.
    \item \eng{The conference has been \textbf{brought forward} to 2 p.m.} (La conférence a été avancée à 14h.)
    \item \eng{They \textbf{cancelled} the trip due to the strike.} (Ils ont annulé le voyage à cause de la grève.)
\end{itemize}
\end{exemplebox}

\begin{aretenir}[Postpone vs Put off vs Bring forward vs Cancel]
\begin{center}
\begin{tabular}{|l|l|l|}
\hline
\rowcolor{poporangeL} \textbf{Verbe} & \textbf{Sens} & \textbf{Construction clé} \\
\hline
\cle{postpone} & Reporter (formel) & \eng{postpone + V-ing} / \eng{postpone sth until + date} \\
\cle{put off} & Reporter (courant, phrasal verb) & \eng{put off + V-ing} / \eng{put sth off until + date} \\
\cle{bring forward} & Avancer (date/heure) & \eng{bring sth forward to + date/heure} \\
\cle{cancel} & Annuler & \eng{cancel sth} (pas de gérondif) \\
\hline
\end{tabular}
\end{center}
\end{aretenir}

\courssub{Punctuality and time status: on time vs in time}

C'est l'une des distinctions les plus subtiles et les plus testées. Lisez attentivement les deux situations.

\begin{exemplebox}[On time vs In time — Observation]
\begin{enumerate}
    \item \eng{The train arrived \textbf{on time}.} (Le train est arrivé \textbf{à l'heure exacte}, selon l'horaire.)
    \item \eng{I arrived \textbf{in time for} the interview.} (Je suis arrivé \textbf{à temps pour} / \textbf{assez tôt pour} l'entretien — peut-être 10 min avant, l'heure exacte n'importe pas, l'essentiel est de ne pas le rater.)
\end{enumerate}
\end{exemplebox}

\begin{propriete}[Distinction \eng{on time} / \eng{in time (for)}]
\begin{itemize}
    \item \cle{on time} = \textbf{ponctuel(le)}, à l'heure prévue, ni en avance ni en retard par rapport à un horaire fixe. Synonyme : \cle{punctual} (adj.).
    \item \cle{in time (for)} = \textbf{assez tôt pour} faire quelque chose, avant qu'il ne soit trop tard. Implique une marge (même mince) avant un événement.
    \item \cle{be late for} : être en retard pour (événement).
    \item \cle{be ahead of schedule} : être en avance sur le planning.
    \item \cle{be behind schedule} : avoir du retard sur le planning.
\end{itemize}
\end{propriete}

\begin{methode}[Choisir entre \eng{on time} et \eng{in time}]
Posez-vous la question : \textbf{« Y a-t-il un horaire précis à respecter ? »}
\begin{enumerate}
    \item \textbf{OUI} (horaire de train, début de cours, heure de rendez-vous fixe) $\rightarrow$ \eng{on time}. \eng{The class starts at 7:30. Be on time!}
    \item \textbf{NON / Événement à ne pas rater} (attraper un bus, rendre un devoir, voir quelqu'un avant son départ) $\rightarrow$ \eng{in time (for)}. \eng{Run! We need to get there in time for the kick-off.}
\end{enumerate}
\end{methode}

\begin{center}
\begin{tabularx}{\linewidth}{|X|X|X|}
\hline
\rowcolor{popgreenL} \textbf{Expression} & \textbf{Sens français} & \textbf{Exemple} \\
\hline
\cle{on time} & à l'heure (ponctuel) & \eng{She is always on time for work.} \\
\cle{in time (for)} & à temps pour / assez tôt pour & \eng{We got home in time for dinner.} \\
\cle{punctual} (adj.) & ponctuel(le) & \eng{Be punctual for the exam.} \\
\cle{be late for} & être en retard pour & \eng{Don't be late for the meeting.} \\
\cle{be ahead of schedule} & en avance sur le planning & \eng{The project is ahead of schedule.} \\
\cle{be behind schedule} & en retard sur le planning & \eng{Construction is behind schedule.} \\
\hline
\end{tabularx}
\end{center}

\courssub{Personal time management: save, waste, spare, run out of}

Au-delà des rendez-vous extérieurs, il faut gérer son temps personnel. Ces expressions sont très idiomatiques.

\begin{definition}[Expressions de gestion personnelle du temps]
\begin{itemize}
    \item \cle{manage one's time} : gérer son temps. \eng{Students must learn to manage their time.}
    \item \cle{save time} : gagner du temps (économiser du temps). \eng{Typing saves time compared to handwriting.}
    \item \cle{waste time} : perdre du temps (gaspiller). \eng{Stop wasting time on social media!}
    \item \cle{spare time} / \cle{free time} : temps libre. \eng{In my spare time, I read novels.} (Souvent interchangeables, \eng{spare time} insiste sur le temps « restant » après obligations.)
    \item \cle{make time for} : se ménager du temps pour, s'arranger pour faire. \eng{You must make time for revision.} (Effort conscient.)
    \item \cle{run out of time} : manquer de temps, ne plus avoir de temps. \eng{We ran out of time and didn't finish the test.}
\end{itemize}
\end{definition}

\begin{attention}[Calques français à éviter]
\begin{itemize}
    \item \eng{*I lost time.} $\rightarrow$ \eng{I \textbf{wasted} time.} (\eng{lose time} existe mais signifie « être ralenti », ex. \eng{We lost time in the traffic}.)
    \item \eng{*I won time.} $\rightarrow$ \eng{I \textbf{saved} time.}
    \item \eng{*I have few time.} $\rightarrow$ \eng{I have \textbf{little} time.} (\eng{time} est indénombrable.)
    \item \eng{*Make a pause.} $\rightarrow$ \eng{Take a break.} / \eng{Have a break.}
\end{itemize}
\end{attention}

\courssub{Pronunciation, spelling and word families}

\begin{propriete}[Familles de mots et orthographe]
\begin{center}
\begin{tabular}{|l|l|l|l|}
\hline
\rowcolor{poppurpleL} \textbf{Verbe} & \textbf{Nom (action/état)} & \textbf{Nom (agent)} & \textbf{Adjectif} \\
\hline
\cle{organise} & \cle{organisation} & \cle{organiser} & \cle{organised} / \cle{organisational} \\
\cle{plan} & \cle{plan} / \cle{planning} & \cle{planner} & \cle{planned} \\
\cle{schedule} & \cle{schedule} / \cle{scheduling} & \cle{scheduler} & \cle{scheduled} \\
\cle{postpone} & \cle{postponement} & — & \cle{postponed} \\
\cle{cancel} & \cle{cancellation} & — & \cle{cancelled} \\
\cle{—} & \cle{punctuality} & — & \cle{punctual} \\
\cle{delay} & \cle{delay} & — & \cle{delayed} \\
\hline
\end{tabular}
\end{center}
\textbf{Orthographe :} \eng{cancel} $\rightarrow$ \eng{cancelled}, \eng{cancellation} (double L en anglais britannique/camerounais). \eng{organise} (UK) vs \eng{organize} (US) — au Cameroun, norme britannique : \eng{organise}.
\end{propriete}

\begin{center}
\begin{tabular}{|l|l|}
\hline
\rowcolor{poppinkL} \textbf{Mot} & \textbf{Prononciation (approx. fr.)} \\
\hline
\cle{schedule} & « chèd-ioul » (UK) / « skèd-ioul » (US) — Au Cameroun : « chèd-ioul » \\
\cle{deadline} & « déd-laïne » \\
\cle{appointment} & « a-poïnt-mente » \\
\cle{postpone} & « pouste-poune » \\
\cle{cancellation} & « canssè-la-ï-chène » \\
\cle{punctuality} & « ponk-tiou-a-li-ti » \\
\cle{duration} & « diou-ré-ï-chène » \\
\hline
\end{tabular}
\end{center}

\courssub{Illustration : A typical week in Terminale A}

La figure ci-dessous synthétise le vocabulaire vu dans une frise hebdomadaire. Observez l'articulation entre \eng{meeting}, \eng{deadline}, \eng{sign-up closes} et \eng{campaign}.

\begin{popfigure}
\begin{tikzpicture}[
    node distance=0.8cm and 1.5cm,
    day/.style={rectangle, draw=popblue, thick, fill=popblueL, minimum width=2.2cm, minimum height=1.2cm, font=\bfseries\small, align=center},
    event/.style={rectangle, draw=poporange, fill=poporangeL, rounded corners=3pt, minimum width=2.2cm, minimum height=0.9cm, font=\small, align=center},
    arrow/.style={-Stealth, popdark, thick}
]
% Nodes for days
\node[day] (mon) {Monday};
\node[day, below of=mon] (tue) {Tuesday};
\node[day, below of=tue] (wed) {Wednesday};
\node[day, below of=wed] (thu) {Thursday};
\node[day, below of=thu] (fri) {Friday};
\node[day, below of=fri] (sat) {Saturday};
\node[day, below of=sat] (sun) {Sunday};

% Events
\node[event, right=of mon, align=center] (mon1){Club meeting\\ 4 p.m.};
\node[event, right=of wed, align=center] (wed1){Deadline:\\ History essay};
\node[event, right=of wed, yshift=-1.2cm, align=center] (wed2){Parent-teacher\\ meeting 5 p.m.};
\node[event, right=of fri, align=center] (fri1){Sign-up closes:\\ Clean-up};
\node[event, right=of sat, align=center] (sat1){Clean-up campaign\\ 7 a.m. sharp\\ Duration: 3h};
\node[event, right=of sun, align=center] (sun1){Rest\\ Plan next week};

% Arrows
\draw[arrow] (mon.east) -- (mon1.west);
\draw[arrow] (wed.east) -- (wed1.west);
\draw[arrow] (wed.east) -- (wed2.west);
\draw[arrow] (fri.east) -- (fri1.west);
\draw[arrow] (sat.east) -- (sat1.west);
\draw[arrow] (sun.east) -- (sun1.west);

% Legend box
\node[draw=popmuted, fill=popcream, rounded corners, font=\footnotesize, align=left, anchor=north east] at (current bounding box.north east) {
    \textbf{Key vocab in context:}\\
    \textbullet\ \eng{meeting} (club)\\
    \textbullet\ \eng{deadline} (essay)\\
    \textbullet\ \eng{appointment} (parent-teacher)\\
    \textbullet\ \eng{sign-up closes} (deadline)\\
    \textbullet\ \eng{campaign} (community)\\
    \textbullet\ \eng{duration} (3h)\\
    \textbullet\ \eng{plan} (next week)
};
\end{tikzpicture}
\legende{Frise hebdomadaire : gestion du temps d'un élève de Terminale A (vocabulaire clé en contexte).}
\end{popfigure}

\courssub{Using the lexis in context: a dialogue at the workplace}

Voici un dialogue authentique entre un stagiaire (\eng{intern}) et son superviseur (\eng{supervisor}) dans une entreprise de transformation de cacao à Douala. Il regroupe le vocabulaire de la section.

\begin{textebox}{Dialogue: Internship review — Managing priorities}
\textbf{Supervisor:} Good morning, Achille. Let's review your \textbf{schedule} for this week. The \textbf{deadline} for the export report is Wednesday noon. Are you \textbf{on track}?

\textbf{Achille:} Good morning, Sir. Yes, I've \textbf{planned} my time carefully. I've \textbf{scheduled} the data collection for Monday and Tuesday. I should \textbf{meet the deadline} without any problem.

\textbf{Supervisor:} Good. But the management meeting has been \textbf{brought forward} to Tuesday 10 a.m. You'll need to \textbf{attend}. It might \textbf{delay} your data entry.

\textbf{Achille:} Understood. I'll \textbf{rearrange} my \textbf{timetable}. I can \textbf{put off} the formatting until Wednesday morning. I'll \textbf{make time for} the meeting.

\textbf{Supervisor:} Excellent. Don't \textbf{waste time} hesitating. If you \textbf{run out of time}, let me know immediately. We cannot \textbf{postpone} the submission to the Ministry.

\textbf{Achille:} I won't, Sir. I'll be \textbf{punctual} for the meeting and \textbf{submit} the report \textbf{in time}.
\end{textebox}

\begin{center}
\begin{tabularx}{\linewidth}{|l|X|}
\hline
\rowcolor{popgoldL} \textbf{Mot / Expression} & \textbf{Traduction / Glossaire} \\
\hline
\eng{on track} & sur la bonne voie, dans les temps \\
\eng{data collection} & collecte de données \\
\eng{brought forward} & avancé (reporté à une date plus proche) \\
\eng{data entry} & saisie de données \\
\eng{rearrange} & réorganiser, modifier l'ordre \\
\eng{formatting} & mise en forme \\
\eng{hesitating} & hésiter \\
\eng{submit} & soumettre, remettre (un document) \\
\eng{the Ministry} & le Ministère \\
\hline
\end{tabularx}
\end{center}

\courssub{Exercice résolu : choisir la bonne expression}

\begin{exoresolu}[Complétez avec l'expression appropriée : \eng{on time / in time / deadline / postponed / cancelled / behind schedule / spare time / run out of time}]
\textbf{1.} The registration \_\_\_\_\_\_\_\_\_ for the competition is next Friday. \\
\textbf{2.} Hurry up! We must arrive \_\_\_\_\_\_\_\_\_ \_\_\_\_\_\_\_\_\_ the start of the ceremony. \\
\textbf{3.} My watch says 7:30. The train leaves at 7:30. It left \_\_\_\_\_\_\_\_\_. \\
\textbf{4.} The football match was \_\_\_\_\_\_\_\_\_ until next Saturday because of the heavy rain. \\
\textbf{5.} The manager \_\_\_\_\_\_\_\_\_ the meeting because he was sick. \\
\textbf{6.} The construction project is \_\_\_\_\_\_\_\_\_ \_\_\_\_\_\_\_\_\_; it will finish two months late. \\
\textbf{7.} In my \_\_\_\_\_\_\_\_\_ \_\_\_\_\_\_\_\_\_, I like playing chess. \\
\textbf{8.} We worked fast because we were \_\_\_\_\_\_\_\_\_ \_\_\_\_\_\_\_\_\_ \_\_\_\_\_\_\_\_\_.

\tcblower
\textbf{Solution détaillée :}
\begin{enumerate}
    \item \textbf{deadline} (nom : date limite). \eng{The registration deadline...}
    \item \textbf{in time for} (assez tôt pour un événement). \eng{...arrive in time for the start...}
    \item \textbf{on time} (ponctuel, à l'heure exacte prévue). \eng{...left on time.}
    \item \textbf{postponed} (reporté — passif). \eng{...was postponed until...} (pas \eng{put off} car passif formel souvent préféré, mais \eng{put off} possible).
    \item \textbf{cancelled} (annulé). \eng{The manager cancelled...}
    \item \textbf{behind schedule} (en retard sur le planning). \eng{...is behind schedule;}
    \item \textbf{spare time} (ou \eng{free time}) : temps libre.
    \item \textbf{running out of time} (manquer de temps — forme continue pour action en cours). \eng{...we were running out of time.}
\end{enumerate}
\end{exoresolu}

\courssub{Exercices d'application}

\begin{exercice}[Entraînement : Vocabulaire et structures]
\textbf{A. Traduisez en anglais en utilisant le vocabulaire de la leçon.}
\begin{enumerate}
    \item La date limite pour le paiement des frais de scolarité est vendredi.
    \item Nous avons reporté la réunion du club de débat à mercredi prochain.
    \item Le train est arrivé à l'heure.
    \item J'ai manqué de temps pour finir la rédaction.
    \item Elle a annulé son rendez-vous chez le dentiste.
    \item Le projet a du retard sur le planning.
    \item Dans mon temps libre, j'aide ma mère au marché.
    \item Il faut que tu te ménages du temps pour réviser tes leçons.
\end{enumerate}

\textbf{B. Choisissez la bonne réponse (a, b ou c).}
\begin{enumerate}
    \item "The bus arrived \_\_\_\_\_\_\_\_\_." (It arrived at 8:00, scheduled for 8:00).
    a) in time \quad b) on time \quad c) at time
    \item "We \_\_\_\_\_\_\_\_\_ the conference until July."
    a) cancelled \quad b) brought forward \quad c) postponed
    \item "Don't \_\_\_\_\_\_\_\_\_ your homework until the last minute."
    a) put off \quad b) put on \quad c) put up
    \item "The rain \_\_\_\_\_\_\_\_\_ our departure."
    a) delayed \quad b) postponed \quad c) cancelled
    \item "I have very \_\_\_\_\_\_\_\_\_ time left."
    a) few \quad b) little \quad c) small
\end{enumerate}
\end{exercice}

\begin{corrige}[Exercice d'application]
\textbf{A. Traduction}
\begin{enumerate}
    \item The \textbf{deadline} for school fees payment is Friday. / \textbf{The deadline to pay} school fees is Friday.
    \item We \textbf{postponed} the debate club meeting \textbf{until} next Wednesday. / We \textbf{put off} the meeting \textbf{until}...
    \item The train arrived \textbf{on time}.
    \item I \textbf{ran out of time} to finish the essay. / I \textbf{was running out of time}...
    \item She \textbf{cancelled} her dentist \textbf{appointment}.
    \item The project is \textbf{behind schedule}.
    \item In my \textbf{spare time} / \textbf{free time}, I help my mother at the market.
    \item You must \textbf{make time for} revising your lessons. / \textbf{Make time to revise}...
\end{enumerate}

\textbf{B. Choix multiple}
\begin{enumerate}
    \item \textbf{b) on time} (Respect de l'horaire précis).
    \item \textbf{c) postponed} (Reporter à plus tard). \eng{Brought forward} = avancer.
    \item \textbf{a) put off} (Phrasal verb courant pour reporter une action : \eng{put off doing sth}).
    \item \textbf{a) delayed} (Verbe transitif : la pluie a retardé NOTRE départ). \eng{Postpone} implique une décision humaine de reporter ; \eng{cancel} = annuler.
    \item \textbf{b) little} (\eng{Time} est indénombrable $\rightarrow$ \eng{little} / \eng{much}. \eng{Few} pour dénombrables).
\end{enumerate}
\end{corrige}

\begin{savaistu}[Culture : "African time" vs "Western time"]
Dans de nombreux contextes anglophones (Royaume-Uni, USA, monde des affaires international), la ponctualité (\eng{punctuality}) est une marque de respect et de professionnalisme. \eng{"Time is money."} Au Cameroun, comme dans beaucoup de cultures africaines, le temps peut être perçu de façon plus fluide (\emph{"African time"}), où les relations priment sur l'horloge stricte. Cependant, dans les examens officiels (GCE, Baccalauréat), les concours administratifs, les entretiens d'embauche ou le travail en entreprise formelle, la norme attendue est \textbf{\eng{on time}}. Savoir naviguer entre ces deux registres est une compétence interculturelle clé : \eng{be on time for the exam}, mais \eng{arrive in time for the celebration} (où l'heure de début est souvent indicative).
\end{savaistu}

\coursec{Collocations and idiomatic expressions}

Dans la section précédente, nous avons isolé les mots-clés du thème « Managing time ». Cependant, en anglais naturel, ces mots ne vivent pas seuls : ils s'associent à d'autres pour former des \cle{collocations} (associations de mots figées, fréquentes et naturelles) ou des \cle{expressions idiomatiques} dont le sens global ne se déduit pas de la somme des sens littéraux. Maîtriser ces blocs lexicaux est indispensable pour passer d'un anglais « scolaire » à un anglais fluide, que ce soit pour \eng{chair a meeting}, \eng{meet a deadline} ou simplement \eng{lend a hand} lors d'une activité communautaire à Buea ou Yaoundé. Observons d'abord ces associations dans leur contexte naturel.

\courssub{Collocations with 'time' : spend, save, waste, kill, take}

\begin{textebox}{Extrait d'un article de blog : « Managing my Week »}
Last Monday, I decided to \textbf{spend time} planning my semester instead of \textbf{wasting time} scrolling through social media. It helped me \textbf{save time} later in the week. On Wednesday, I had to \textbf{kill time} at the bus station in Douala because the bendskin was late, so I reviewed my vocabulary cards. My grandmother always says : “ \textbf{Take your time} to do things well, because \textbf{time flies}. ”
\tcblower
\textbf{Glossaire :} \textit{scroll through} (faire défiler), \textit{semester} (semestre), \textit{bus station} (gare routière), \textit{bendskin} (moto-taxi), \textit{review} (réviser).
\end{textebox}

\begin{propriete}{Les collocations majeures avec \textbf{time}}
En anglais, \eng{time} (temps) ne s'utilise pas avec n'importe quel verbe. Les associations suivantes sont figées et obligatoires pour exprimer ces idées :
\begin{center}
\begin{tabularx}{\linewidth}{|X|X|X|}
\hline
\rowcolor{popblueL}
\textbf{Collocation} & \textbf{Sens français} & \textbf{Exemple traduit} \\
\hline
\eng{spend time (+ -ing)} & passer du temps (à faire qqch) & \eng{We spent three hours cleaning the classroom.} (Nous avons passé trois heures à nettoyer la salle.) \\
\hline
\eng{waste time (+ -ing)} & perdre du temps (à faire qqch) & \eng{Don't waste time arguing.} (Ne perds pas ton temps à te disputer.) \\
\hline
\eng{save time} & gagner du temps / économiser du temps & \eng{Using a template saves time.} (Utiliser un modèle fait gagner du temps.) \\
\hline
\eng{kill time} & tuer le temps / occuper son temps (en attendant) & \eng{I read a newspaper to kill time.} (J'ai lu un journal pour tuer le temps.) \\
\hline
\eng{take one's time} & prendre son temps (ne pas se presser) & \eng{Take your time, the deadline is next week.} (Prends ton temps, l'échéance est la semaine prochaine.) \\
\hline
\eng{have a good time} & passer un bon moment / s'amuser & \eng{The students had a good time at the cultural festival.} (Les élèves se sont bien amusés au festival culturel.) \\
\hline
\eng{time flies} & le temps passe vite / le temps file & \eng{Time flies when you're studying hard.} (Le temps passe vite quand on étudie dur.) \\
\hline
\end{tabularx}
\end{center}
\end{propriete}

\begin{attention}{Piège majeur : \eng{Time flies} $\neq$ « Le temps vole »}
Ne traduisez \textbf{jamais} les idiomes mot à mot. \eng{Time flies} (le temps file / passe vite) n'a rien à voir avec l'oiseau (\eng{a fly} / \eng{to fly}). De même, \eng{kill time} ne signifie pas « tuer le temps » au sens violent, mais « occuper un temps mort ». \eng{Save time} ne prend pas de préposition (pas de *\eng{save *of* time}).
\end{attention}

\courssub{Collocations with 'part' : take part in, play a part, do one's part}

Le nom \eng{part} (partie, rôle) génère des expressions figées essentielles pour exprimer la participation.

\begin{propriete}{Collocations clés avec \textbf{part}}
\begin{center}
\begin{tabularx}{\linewidth}{|X|X|X|}
\hline
\rowcolor{poporangeL}
\textbf{Expression} & \textbf{Sens / Contexte} & \textbf{Exemple traduit} \\
\hline
\eng{take part in} (+ nom / -ing) & participer à (activités, événements, concours) & \eng{All Form 5 students must take part in the mock exams.} (Tous les élèves de 1ère doivent participer aux examens blancs.) \\
\hline
\eng{play a part / role in} & jouer un rôle dans (souvent positif, influent) & \eng{Teachers play a vital part in shaping society.} (Les enseignants jouent un rôle vital dans la formation de la société.) \\
\hline
\eng{do one's part} & faire sa part / faire son devoir (effort collectif) & \eng{If everyone does their part, the school will be clean.} (Si chacun fait sa part, l'école sera propre.) \\
\hline
\eng{for my part} & pour ma part / quant à moi (marque d'opinion personnelle) & \eng{For my part, I prefer morning classes.} (Pour ma part, je préfère les cours du matin.) \\
\hline
\end{tabularx}
\end{center}
\end{propriete}

\begin{exemplebox}{Nuance : \eng{Take part in} vs \eng{Participate in}}
\eng{Take part in} est plus courant à l'oral et dans l'écrit semi-formel (journal scolaire, discours d'élèves). \eng{Participate in} est plus formel (rapports administratifs, lettres officielles).
\begin{itemize}
\item \eng{Students \textbf{took part in} the cleaning campaign.} (Registre standard)
\item \eng{The delegation \textbf{participated in} the international summit.} (Registre formel)
\end{itemize}
\end{exemplebox}

\courssub{Verb + Noun collocations : Meetings, Deadlines and Schedules}

Dans la vie professionnelle, associative ou scolaire (Conseil d'établissement, club de presse, association des parents d'élèves), certaines associations \textbf{verbe + nom} sont obligatoires. Les utiliser prouve votre maîtrise de l'anglais de communication.

\begin{popfigure}
\begin{tikzpicture}[
    box/.style={draw=popblue, thick, rounded corners, fill=popblue!5, minimum width=7cm, minimum height=5.5cm, inner sep=5mm},
    title/.style={fill=popblue, text=white, font=\bfseries\large, rounded corners=3pt, inner sep=2pt},
    item/.style={font=\ttfamily\small, text=popdark, anchor=west}
]
% Colonne 1 : TIME
\node[box] (timebox) at (0,0) {};
\node[title, anchor=north] at (timebox.north) {TIME COLLOCATIONS};
\begin{scope}[xshift=-3.2cm, yshift=2.2cm]
\node[item] (t1) {\textbullet\ \eng{spend time}};
\node[item, below=4mm of t1] (t2) {\textbullet\ \eng{save time}};
\node[item, below=4mm of t2] (t3) {\textbullet\ \eng{waste time}};
\node[item, below=4mm of t3] (t4) {\textbullet\ \eng{kill time}};
\node[item, below=4mm of t4] (t5) {\textbullet\ \eng{take one's time}};
\node[item, below=4mm of t5] (t6) {\textbullet\ \eng{have a good time}};
\node[item, below=4mm of t6, text=popgreen!80!black] (t7) {\textbullet\ \eng{time flies} (idiom)};
\end{scope}

% Colonne 2 : MEETING
\node[box] (meetbox) at (8.5,0) {};
\node[title, anchor=north] at (meetbox.north) {MEETING COLLOCATIONS};
\begin{scope}[xshift=5.3cm, yshift=2.2cm]
\node[item] (m1) {\textbullet\ \eng{hold a meeting}};
\node[item, below=4mm of m1] (m2) {\textbullet\ \eng{attend a meeting}};
\node[item, below=4mm of m2] (m3) {\textbullet\ \eng{chair a meeting}};
\node[item, below=4mm of m3] (m4) {\textbullet\ \eng{call a meeting}};
\node[item, below=4mm of m4] (m5) {\textbullet\ \eng{cancel a meeting}};
\node[item, below=4mm of m5] (m6) {\textbullet\ \eng{meet a deadline}};
\node[item, below=4mm of m6] (m7) {\textbullet\ \eng{set a deadline}};
\node[item, below=4mm of m7] (m8) {\textbullet\ \eng{keep to a schedule}};
\node[item, below=4mm of m8] (m9) {\textbullet\ \eng{draw up a timetable}};
\end{scope}
\end{tikzpicture}
\legende{Carte mentale des collocations « Time » et « Meeting/Deadline » à mémoriser par blocs.}
\end{popfigure}

\begin{propriete}{Collocations professionnelles et scolaires (Verbe + Nom)}
\begin{center}
\begin{tabularx}{\linewidth}{|l|X|X|}
\hline
\rowcolor{popgreenL}
\textbf{Collocation} & \textbf{Sens précis} & \textbf{Contexte d'usage / Exemple} \\
\hline
\eng{hold a meeting} & tenir/ organiser une réunion (organisateur) & \eng{The principal will \textbf{hold a meeting} with prefects tomorrow.} \\
\hline
\eng{attend a meeting} & assister à / être présent à une réunion (participant) & \eng{All delegates must \textbf{attend the meeting} punctually.} \\
\hline
\eng{chair a meeting} & présider une réunion (rôle de leader) & \eng{The class monitor \textbf{chaired the meeting} efficiently.} \\
\hline
\eng{call a meeting} & convoquer / appeler à une réunion & \eng{The manager \textbf{called an emergency meeting}.} \\
\hline
\eng{cancel a meeting} & annuler une réunion & \eng{Due to the rain, we had to \textbf{cancel the meeting}.} \\
\hline
\eng{meet a deadline} & respecter une échéance (terminer à temps) & \eng{We worked hard to \textbf{meet the deadline}.} \\
\hline
\eng{set a deadline} & fixer une date limite & \eng{The teacher \textbf{set a deadline} for Friday noon.} \\
\hline
\eng{keep to a schedule} & respecter un emploi du temps / un planning & \eng{It is crucial to \textbf{keep to the schedule} during exams.} \\
\hline
\eng{draw up a timetable} & établir / rédiger un emploi du temps & \eng{The committee \textbf{drew up a timetable} for the cultural week.} \\
\hline
\end{tabularx}
\end{center}
\end{propriete}

\begin{attention}{Faux amis et calques fréquents}
\begin{itemize}
\item \textbf{*Assist a meeting} $\rightarrow$ \textbf{FAUX}. \eng{Assist} = aider. On dit \eng{\textbf{attend} a meeting}.
\item \textbf{*Make a meeting} $\rightarrow$ \textbf{FAUX} (sauf « fabriquer »). On dit \eng{\textbf{hold} / \textbf{call} a meeting}.
\item \textbf{*Respect a deadline} $\rightarrow$ \textbf{FAUX} (calque du français). On dit \eng{\textbf{meet} a deadline} ou \eng{\textbf{hit} a deadline}.
\item \textbf{*Do a timetable} $\rightarrow$ \textbf{PEU NATUREL}. On préfère \eng{\textbf{draw up} / \textbf{prepare} / \textbf{make} a timetable}.
\end{itemize}
\end{attention}

\courssub{Idiomatic expressions related to time and pressure}

Les idiomes colorent la langue. Ils sont très présents dans les médias anglophones (BBC, CNN, journaux nigérians, sud-africains) et dans les conversations de travail.

\begin{definition}[In the nick of time]
\eng{In the nick of time} : \textbf{au tout dernier moment possible, juste à temps} (souvent avec soulagement).
\eng{The firefighters arrived \textbf{in the nick of time} to save the building.} (Les pompiers sont arrivés juste à temps pour sauver le bâtiment.)
\end{definition}

\begin{center}
\begin{tabularx}{\linewidth}{|X|X|X|}
\hline
\rowcolor{poppurpleL}
\textbf{Idiom} & \textbf{Sens français} & \textbf{Exemple contextualisé} \\
\hline
\eng{around the clock} & 24h/24, jour et nuit, sans arrêt & \eng{The medical team worked \textbf{around the clock} during the epidemic.} \\
\hline
\eng{at the eleventh hour} & à la dernière minute, au dernier moment & \eng{He submitted his application \textbf{at the eleventh hour}.} \\
\hline
\eng{once in a blue moon} & très rarement, une fois tous les 36 du mois & \eng{Our principal visits the dormitories \textbf{once in a blue moon}.} \\
\hline
\eng{better late than never} & mieux vaut tard que jamais & \eng{Here is your report. \textbf{Better late than never}!} \\
\hline
\eng{time is money} & le temps, c'est de l'argent (efficacité économique) & \eng{Let's start now; \textbf{time is money}.} \\
\hline
\eng{beat the clock} & faire qqch avant la fin du temps imparti, gagner du temps sur l'horloge & \eng{We must \textbf{beat the clock} to finish the project today.} \\
\hline
\end{tabularx}
\end{center}

\begin{savaistu}{Time is money : Benjamin Franklin}
L'expression \eng{« Time is money »} a été popularisée par \textbf{Benjamin Franklin} dans son essai \textit{Advice to a Young Tradesman} (1748). Il écrivait : \eng{« Remember that time is money. »} Aujourd'hui, c'est une maxime centrale dans la culture du travail anglo-saxonne (USA, UK, Nigeria, Ghana, Afrique du Sud) : le temps est une ressource qu'on \eng{spend}, \eng{save}, \eng{invest} ou \eng{waste}, exactement comme l'argent.
\end{savaistu}

\begin{exemplebox}{Exemples traduits (Idiomes en contexte)}
\begin{itemize}
\item \eng{We worked \textbf{around the clock} to finish the report on time.} (Nous avons travaillé 24h/24 pour finir le rapport à temps.)
\item \eng{The agreement was signed \textbf{at the eleventh hour}.} (L'accord a été signé à la dernière minute / au dernier moment.)
\item \eng{She visits her village \textbf{once in a blue moon}.} (Elle rend visite à son village très rarement.)
\item \eng{The team \textbf{beat the clock} and delivered the project before 5 PM.} (L'équipe a gagné la course contre la montre et a livré le projet avant 17h.)
\end{itemize}
\end{exemplebox}

\courssub{Expressions of commitment : Register (Informal vs Formal)}

S'engager dans une activité collective (nettoyage du campus, collecte de fonds, projet de groupe) demande un vocabulaire d'adhésion. Le choix du registre (familier / standard / formel) dépend de votre interlocuteur.

\begin{textebox}{Dialogue : Organising a Fund-raising Concert}
\textbf{A:} We're organising a fund-raising concert for the school library. \textbf{Are you on board?}
\textbf{B:} \textbf{Count me in!} I'll \textbf{lend a hand} with the tickets.
\textbf{A:} Great! But we must \textbf{beat the clock} — the deadline is next Friday. Everyone needs to \textbf{pull their weight}.
\textbf{B:} Don't worry. \textbf{For my part}, I'll make sure the posters are ready by Wednesday.
\tcblower
\textbf{Glossaire :} \textit{fund-raising} (collecte de fonds), \textit{on board} (d'accord, partant, dans l'équipe), \textit{count me in} (compte sur moi, j'y suis), \textit{lend a hand} (donner un coup de main), \textit{beat the clock} (gagner la course contre la montre), \textit{deadline} (date limite), \textit{pull one's weight} (faire sa part, assumer sa charge), \textit{posters} (affiches).
\end{textebox}

\begin{propriete}{Engagement et registre : du familier au formel}
\begin{center}
\begin{tabularx}{\linewidth}{|X|X|X|}
\hline
\rowcolor{poppinkL}
\textbf{Expression (Informel / Standard)} & \textbf{Sens} & \textbf{Équivalent Formel (Écrit / Officiel)} \\
\hline
\eng{Count me in!} & Compte sur moi / J'accepte & \eng{I would be delighted to participate.} / \eng{I accept the invitation.} \\
\hline
\eng{Count me out.} & Compte sans moi / Je refuse & \eng{I regret that I cannot participate.} / \eng{I must decline.} \\
\hline
\eng{Are you on board?} & Es-tu partant / d'accord ? & \eng{Do you agree to take part?} / \eng{Will you join the committee?} \\
\hline
\eng{Lend a hand} & Donner un coup de main / aider & \eng{Provide assistance} / \eng{Contribute support} \\
\hline
\eng{Pull one's weight} & Faire sa part (effort équitable) & \eng{Fulfill one's obligations} / \eng{Contribute one's fair share} \\
\hline
\eng{Do one's bit / Do one's part} & Faire sa part (général) & \eng{Fulfill one's role} / \eng{Discharge one's duties} \\
\hline
\end{tabularx}
\end{center}
\end{propriete}

\begin{attention}{Registre : Attention au contexte !}
Ne dites pas \eng{« Count me in »} dans une lettre de motivation pour un stage ou un emploi, ni face au Proviseur. Écrivez : \eng{« I would be delighted to contribute to this project. »}. Inversement, dans un SMS à un camarade pour le nettoyage de la classe, \eng{« I would be delighted to participate »} sonne faux et guindé. Adaptez votre \cle{registre de langue}.
\end{attention}

\courssub{Word families and pronunciation focus (Preview)}

Avant de passer aux exercices, notez les familles de mots dérivées des racines vues ici. La section 5 y reviendra en détail, mais anticipez l'orthographe et l'accent tonique :

\begin{center}
\begin{tabular}{|l|l|l|l|}
\hline
\rowcolor{popgoldL}
\textbf{Verbe / Base} & \textbf{Nom (Action/Agent)} & \textbf{Adjectif / Participe} & \textbf{Accent tonique (approx.)} \\
\hline
\eng{participate} & \eng{participation, participant} & \eng{participatory} & par-\textbf{TI}-ci-pate / par-ti-ci-\textbf{PA}-tion \\
\hline
\eng{organise} & \eng{organisation, organiser} & \eng{organised, organisational} & \textbf{OR}-gan-ise / or-gan-i-\textbf{SA}-tion \\
\hline
\eng{manage} & \eng{management, manager} & \eng{manageable, managerial} & \textbf{MA}-nage / \textbf{MA}-nage-ment \\
\hline
\eng{contribute} & \eng{contribution, contributor} & \eng{contributory} & con-\textbf{TRI}-bute / con-tri-\textbf{BU}-tion \\
\hline
\eng{deadline} (nom composé) & -- & \eng{deadline-driven} & \textbf{DEAD}-line \\
\hline
\end{tabular}
\end{center}

\begin{aretenir}{Synthèse : Le lexique indispensable à maîtriser}
\begin{itemize}
\item \textbf{Temps :} \eng{spend/waste/save/kill/take time}, \eng{time flies}, \eng{meet/set a deadline}, \eng{keep to a schedule}, \eng{draw up a timetable}.
\item \textbf{Participation :} \eng{take part in}, \eng{play a part in}, \eng{do one's part}, \eng{for my part}.
\item \textbf{Réunions :} \eng{hold/attend/chair/call/cancel a meeting}.
\item \textbf{Idiomes pression/temps :} \eng{in the nick of time}, \eng{around the clock}, \eng{at the eleventh hour}, \eng{once in a blue moon}, \eng{beat the clock}, \eng{time is money}.
\item \textbf{Registre engagement :} \eng{Count me in/out} (informel) $\leftrightarrow$ \eng{I would be delighted to participate / I must decline} (formel) ; \eng{lend a hand} $\leftrightarrow$ \eng{provide assistance} ; \eng{pull one's weight} $\leftrightarrow$ \eng{fulfill one's obligations}.
\item \textbf{Orthographe UK :} \eng{organise}, \eng{organisation}, \eng{programme} (mais \eng{program} en informatique).
\end{itemize}
\end{aretenir}

\begin{exoresolu}[Transformation de registre (Formal $\leftrightarrow$ Informal)]
\textbf{Consigne :} Réécrivez les phrases suivantes en changeant de registre (Informel $\rightarrow$ Formel ou Formel $\rightarrow$ Informel) en utilisant les expressions de la leçon.

\begin{enumerate}
\item \textbf{Informal :} \eng{Hey guys, are you on board for the cleaning Saturday? Count me in!}
\item \textbf{Formal :} \eng{The committee requests that all members fulfill their obligations regarding the upcoming deadline.}
\item \textbf{Informal :} \eng{We have to beat the clock if we want to finish the banner.}
\item \textbf{Formal :} \eng{I would be delighted to provide assistance with the organisation of the conference.}
\end{enumerate}

\tcblower
\textbf{Solution détaillée :}
\begin{enumerate}
\item \eng{\textbf{Formal :} « Dear colleagues, do you agree to participate in the cleaning day this Saturday? I would be delighted to take part. »} (\textit{on board} $\rightarrow$ \textit{agree to participate} ; \textit{count me in} $\rightarrow$ \textit{I would be delighted to take part}).
\item \eng{\textbf{Informal :} « Everyone needs to pull their weight before the deadline. »} (\textit{fulfill their obligations} $\rightarrow$ \textit{pull their weight} ; \textit{regarding} $\rightarrow$ contexte implicite).
\item \eng{\textbf{Formal :} « We must work against the clock / meet the deadline to complete the banner. »} (\textit{beat the clock} $\rightarrow$ \textit{work against the clock} ou \textit{meet the deadline}).
\item \eng{\textbf{Informal :} « Count me in! I'll lend a hand with organising the conference. »} (\textit{delighted to provide assistance} $\rightarrow$ \textit{count me in / lend a hand} ; \textit{organisation} $\rightarrow$ \textit{organising}).
\end{enumerate}
\end{exoresolu}

\begin{exercice}[Entraînement : Collocations, Idiomes et Registre]
\begin{enumerate}
\item \textbf{Complétez} avec la collocation correcte (\eng{hold, attend, chair, call, cancel, meet, set, keep, draw up}) : \\
The Principal decided to \_\_\_\_\_\_ a meeting for Monday. The Vice-Principal will \_\_\_\_\_\_ it. All teachers must \_\_\_\_\_\_. We hope nobody will \_\_\_\_\_\_ it. The secretary will \_\_\_\_\_\_ a timetable for the week. Students must \_\_\_\_\_\_ the deadline for registration.
\item \textbf{Associez} l'idiome à sa définition :
\begin{center}
\begin{tabular}{ll}
1. Around the clock & a. Juste à temps \\
2. In the nick of time & b. 24h/24 \\
3. At the eleventh hour & c. Très rarement \\
4. Once in a blue moon & d. À la dernière minute \\
5. Beat the clock & e. Gagner la course contre la montre \\
\end{tabular}
\end{center}
\item \textbf{Traduisez} en anglais (choisissez le registre approprié) : \\
a) (À un ami) « Allez, participe au tournoi de foot ! » \\
b) (Dans un rapport officiel) « Tous les membres ont rempli leurs obligations. » \\
c) (Dialogue) « On a fini juste à temps ! » \\
d) (Conseil) « Ne perds pas ton temps sur les réseaux sociaux. »
\end{enumerate}
\end{exercice}

\begin{corrige}[Exercice n]
\begin{enumerate}
\item \eng{call / chair / attend / cancel / draw up / meet}.
\item 1-b, 2-a, 3-d, 4-c, 5-e.
\item a) \eng{Come on, take part in the football tournament!} (ou \eng{Join in!}) \\
b) \eng{All members have fulfilled their obligations.} \\
c) \eng{We finished in the nick of time!} (ou \eng{at the eleventh hour} / \eng{just in time}). \\
d) \eng{Don't waste time on social media.}
\end{enumerate}
\end{corrige}

\coursec{False friends and French calques}

After exploring the collocations and idiomatic expressions that give fluency to your participation and time-management discourse, we must now address a major source of errors for Francophone learners: \textbf{false friends (faux amis)} and \textbf{French calques (calques du français)}. These “transparent” words look familiar but betray you in exams and real communication. Mastering them is the hallmark of a candidate aiming for excellence at the Baccalauréat.

\courssub{Observation : when familiar words deceive}

\begin{textebox}[Authentic workplace dialogue]
— \eng{Are you going to \textbf{attend} the opening ceremony of the new cocoa processing plant in Douala next Monday?}
— \eng{No, I cannot \textbf{attend}. I have a prior \textbf{engagement}. Actually, I am \textbf{currently} negotiating a contract in Yaoundé.}
— \eng{I thought you were \textbf{eventually} moving to the capital.}
— \eng{That was the plan, but the project was \textbf{delayed}. We need to \textbf{manage our time} better. Let's not \textbf{participate in} the confusion; I will \textbf{participate in} the video conference instead.}
— \eng{Good idea. I will update the \textbf{agenda} and send it to your \textbf{planner}. Don't \textbf{demand} a reply immediately, though; just \textbf{ask for} confirmation.}
\end{textebox}

\begin{definition}[Glossaire du dialogue]
\begin{itemize}
\item \eng{attend} (v.) : assister à, être présent à.
\item \eng{engagement} (n.) : engagement, rendez-vous, obligation.
\item \eng{currently} (adv.) : actuellement, en ce moment.
\item \eng{actually} (adv.) : en fait, en réalité.
\item \eng{eventually} (adv.) : finalement, à la fin (après un processus).
\item \eng{delay} (v./n.) : retarder / retard ; \eng{deadline} (n.) : délai, date limite.
\item \eng{manage time} (expr.) : gérer son temps (pas \eng{programme time}).
\item \eng{participate in} (v. + prép.) : participer à.
\item \eng{agenda} (n.) : ordre du jour (réunion).
\item \eng{planner} (n.) / \eng{diary} (n. UK) : agenda personnel.
\item \eng{demand} (v.) : exiger (fort) ; \eng{ask for} : demander.
\item \eng{assist} (v.) : aider, porter assistance (ne signifie \textbf{jamais} “assister à”).
\end{itemize}
\end{definition}

\courssub{The “Big Five” false friends: rigorous analysis}

The following five pairs account for the vast majority of lexical errors in the \emph{Vocabulary} and \emph{Writing} papers. Study the \textbf{form}, the \textbf{meaning}, and the \textbf{contextual usage} carefully.

\begin{propriete}[False Friend 1 : \eng{Assist} vs \eng{Attend}]
\textbf{Forme :} \eng{Assist} (verbe transitif direct) / \eng{Attend} (verbe transitif direct).
\textbf{Emploi :}
\begin{itemize}
    \item \eng{Assist} = \textbf{Aider}, porter secours à quelqu'un. \eng{She assisted the surgeon during the operation.} « Elle a aidé le chirurgien pendant l'opération. »
    \item \eng{Attend} = \textbf{Assister à}, être présent à un événement. \eng{All students must attend the orientation session.} « Tous les élèves doivent assister à la séance d'orientation. »
\end{itemize}
\textbf{Piège :} En français, “assister” signifie “être présent”. En anglais, \eng{assist} ne prend \textbf{jamais} la préposition \eng{to} dans ce sens. \eng{Assist to} n'existe pas.
\end{propriete}

\begin{propriete}[False Friend 2 : \eng{Actually} vs \eng{Currently}]
\textbf{Forme :} Deux adverbes en \eng{-ly}.
\textbf{Emploi :}
\begin{itemize}
    \item \eng{Currently} / \eng{At present} = \textbf{Actuellement}, en ce moment (marque le présent). \eng{The Minister is currently visiting the Far North region.} « Le ministre visite actuellement la région de l'Extrême-Nord. »
    \item \eng{Actually} / \eng{In fact} = \textbf{En fait}, en réalité (marque la correction, la surprise, l'opposition à une attente). \eng{People think Buea is cold; actually, it is quite pleasant.} « On pense que Buea est froid ; en fait, c'est assez agréable. »
\end{itemize}
\textbf{Astuce mnémotechnique :} \eng{Current} $\rightarrow$ \eng{Current affairs} (l'actualité) $\rightarrow$ \eng{Currently} = \textbf{Actuellement}. \eng{Actual} $\rightarrow$ \eng{Actual facts} (les faits réels) $\rightarrow$ \eng{Actually} = \textbf{En fait}.
\end{propriete}

\begin{propriete}[False Friend 3 : \eng{Eventually} vs \eng{Possibly} / \eng{If necessary}]
\textbf{Forme :} Adverbes.
\textbf{Emploi :}
\begin{itemize}
    \item \eng{Eventually} = \textbf{Finalement}, après un certain temps, après des efforts. Implique une \textbf{chronologie} et une \textbf{conclusion}. \eng{After hours of debate, the committee eventually reached a consensus.} « Après des heures de débat, le comité est finalement parvenu à un consensus. »
    \item \eng{Possibly} / \eng{If necessary} / \eng{If need be} = \textbf{Éventuellement}, si nécessaire. Implique une \textbf{conditionnalité}, une option.
\end{itemize}
\textbf{Distinction cruciale :} \eng{“I will eventually call you”} = Je t'appellerai \textbf{finalement} (un jour, après avoir hésité). \eng{“I will call you if necessary”} = Je t'appellerai \textbf{éventuellement} (si c'est nécessaire).
\end{propriete}

\begin{propriete}[False Friend 4 : \eng{Delay} vs \eng{Deadline} / \eng{Time limit}]
\textbf{Forme :} Noms (et verbes pour \eng{delay}).
\textbf{Emploi :}
\begin{itemize}
    \item \eng{A delay} / \eng{To be delayed} = \textbf{Un retard} / \textbf{Être en retard}. \eng{The train was delayed by 30 minutes.} « Le train a eu 30 minutes de retard. »
    \item \eng{A deadline} / \eng{A time limit} = \textbf{Un délai}, une \textbf{date limite} (butoir). \eng{The deadline for registration is Friday.} « Le délai d'inscription expire vendredi. »
\end{itemize}
\textbf{Attention :} \eng{Delay} a une connotation négative (non-respect de l'heure). \eng{Deadline} est neutre/administratif (date butoir).
\end{propriete}

\begin{propriete}[False Friend 5 : \eng{Demand} vs \eng{Ask for}]
\textbf{Forme :} Verbes transitifs.
\textbf{Emploi :}
\begin{itemize}
    \item \eng{Demand} = \textbf{Exiger}, réclamer avec autorité (ton fort, impératif). \eng{The workers demanded a pay rise.} « Les travailleurs ont exigé une augmentation de salaire. »
    \item \eng{Ask for} / \eng{Request} (formel) = \textbf{Demander} (poliment, standard). \eng{She asked for an extension of the deadline.} « Elle a demandé une prolongation du délai. »
\end{itemize}
\textbf{Conséquence :} Écrire \eng{“I demand your help”} dans un mail formel est perçu comme \textbf{agressif et impoli}. Utilisez \eng{“I would like to request your assistance”} ou \eng{“Could I ask for your help?”}.
\end{propriete}

\courssub{French calques: structural traps}

Beyond single words, entire syntactic structures are calqued from French. These are penalised in \emph{Writing} (coherence/accuracy) and \emph{Speaking} (fluency).

\begin{propriete}[Calque 1 : \eng{Participate IN} (never \eng{TO})]
En français : “Participer \textbf{à}”. La préposition \eng{à} tente l'élève d'écrire \eng{to}.
\textbf{Règle :} \eng{Participate} + \textbf{IN} + \eng{gerund} / \eng{noun}.
\begin{center}
\begin{tabularx}{\linewidth}{|X|X|}\hline
\rowcolor{popblueL} \textbf{Correct} & \textbf{Incorrect (Calque FR)} \\ \hline
\eng{He participated in the workshop.} & \eng{He participated to the workshop.} \\ \hline
\eng{We look forward to participating in the fair.} & \eng{We look forward to participate to the fair.} \\ \hline
\end{tabularx}
\end{center}
\end{propriete}

\begin{propriete}[Calque 2 : \eng{Be BUSY} vs \eng{Be OCCUPIED}]
En français : “Je suis \textbf{occupé}”.
\textbf{Nuance :}
\begin{itemize}
    \item \eng{I am busy} = J'ai du travail, je suis pris (standard, courant). \eng{“Sorry, I'm busy right now.”} « Désolé, je suis occupé en ce moment. »
    \item \eng{I am occupied} = Je suis \textbf{physiquement rempli} (ex. toilettes, siège) ou sonne formel/vieilli pour une personne. \eng{“The seat is occupied.”} « Le siège est pris. »
\end{itemize}
\textbf{Conseil :} Pour dire “j'ai beaucoup de travail”, utilisez toujours \eng{busy}, \eng{tied up} (informel) ou \eng{snowed under} (idiomatique : submergé).
\end{propriete}

\begin{propriete}[Calque 3 : \eng{Manage / Plan / Organise ONE'S TIME}]
En français : “Programmer son temps / son emploi du temps”.
\textbf{Erreur :} \eng{*I need to programme my time.} (\eng{Programme} = programmer un ordinateur, une machine, un spectacle).
\textbf{Correct :}
\begin{itemize}
    \item \eng{Manage my time} (gérer).
    \item \eng{Plan my schedule / week} (planifier).
    \item \eng{Organise my timetable} (organiser son emploi du temps).
\end{itemize}
\end{propriete}

\begin{propriete}[Calque 4 : \eng{Agenda} vs \eng{Diary / Planner}]
C'est le piège “fournitures scolaires / bureau” par excellence.
\begin{center}
\begin{tabularx}{\linewidth}{|X|X|X|}\hline
\rowcolor{popblueL} \textbf{Mot anglais} & \textbf{Sens anglais} & \textbf{Équivalent français} \\ \hline
\eng{Agenda} & List of items to be discussed at a meeting & \textbf{Ordre du jour} \\ \hline
\eng{Diary} (UK) / \eng{Planner} (US) & Book with dates for appointments & \textbf{Agenda (personnel)} \\ \hline
\eng{Calendar} & System of dates / Wall chart & \textbf{Calendrier} \\ \hline
\end{tabularx}
\end{center}
\eng{“Please check the \textbf{agenda} for today's meeting.”} (Ordre du jour). \eng{“I'll write it in my \textbf{diary}.”} (Mon agenda perso).
\end{propriete}

\courssub{Synthesis table: the “FAUX AMIS” wall}

The following TikZ figure summarises the most dangerous false friends for this chapter. Print it mentally before the exam.

\begin{popfigure}
\begin{tikzpicture}[
    node distance=0pt,
    header/.style={fill=popblue, text=white, font=\bfseries\small, minimum height=1.1cm, align=center},
    cell/.style={fill=popblueL, font=\small, minimum height=0.9cm, align=center, text width=5.5cm},
    cellalt/.style={fill=popcream, font=\small, minimum height=0.9cm, align=center, text width=5.5cm},
    header2/.style={fill=popdark, text=white, font=\bfseries\small, minimum height=1.1cm, align=center, text width=5.5cm},
    line/.style={draw=popline, thin}
]
% En-têtes
\node[header, text width=5.5cm] (h1) at (0,0) {English Word};
\node[header, text width=5.5cm] (h2) at (5.5,0) {True Meaning (Vrai sens)};
\node[header2, text width=5.5cm] (h3) at (11,0) {False Friend to Avoid (Piège FR)};

% Ligne 1
\node[cell, anchor=north west] (r1c1) at (0,-1.1) {\eng{Assist}};
\node[cell, anchor=north west] (r1c2) at (5.5,-1.1) {\eng{To help / aid (Aider)}};
\node[cellalt, anchor=north west] (r1c3) at (11,-1.1) {\eng{Assister à (Attend)}};

% Ligne 2
\node[cellalt, anchor=north west] (r2c1) at (0,-2.0) {\eng{Attend}};
\node[cellalt, anchor=north west] (r2c2) at (5.5,-2.0) {\eng{To be present at (Assister à)}};
\node[cell, anchor=north west] (r2c3) at (11,-2.0) {\eng{S'occuper de (Take care of)}};

% Ligne 3
\node[cell, anchor=north west] (r3c1) at (0,-2.9) {\eng{Actually}};
\node[cell, anchor=north west] (r3c2) at (5.5,-2.9) {\eng{In fact / Really (En fait)}};
\node[cellalt, anchor=north west] (r3c3) at (11,-2.9) {\eng{Actuellement (Currently)}};

% Ligne 4
\node[cellalt, anchor=north west] (r4c1) at (0,-3.8) {\eng{Currently}};
\node[cellalt, anchor=north west] (r4c2) at (5.5,-3.8) {\eng{At present (Actuellement)}};
\node[cell, anchor=north west] (r4c3) at (11,-3.8) {\eng{— (Mot correct pour « actuellement »)}};

% Ligne 5
\node[cell, anchor=north west] (r5c1) at (0,-4.7) {\eng{Eventually}};
\node[cell, anchor=north west] (r5c2) at (5.5,-4.7) {\eng{Finally / In the end (Finalement)}};
\node[cellalt, anchor=north west] (r5c3) at (11,-4.7) {\eng{Éventuellement (Possibly / If necessary)}};

% Ligne 6
\node[cellalt, anchor=north west] (r6c1) at (0,-5.6) {\eng{Delay (n/v)}};
\node[cellalt, anchor=north west] (r6c2) at (5.5,-5.6) {\eng{Lateness / To be late (Retard / Retarder)}};
\node[cell, anchor=north west] (r6c3) at (11,-5.6) {\eng{Délai / Date limite (Deadline / Time limit)}};

% Ligne 7
\node[cell, anchor=north west] (r7c1) at (0,-6.5) {\eng{Demand}};
\node[cell, anchor=north west] (r7c2) at (5.5,-6.5) {\eng{To require forcefully (Exiger)}};
\node[cellalt, anchor=north west] (r7c3) at (11,-6.5) {\eng{Demander poliment (Ask for / Request)}};

% Ligne 8
\node[cellalt, anchor=north west] (r8c1) at (0,-7.4) {\eng{Agenda}};
\node[cellalt, anchor=north west] (r8c2) at (5.5,-7.4) {\eng{Meeting schedule (Ordre du jour)}};
\node[cell, anchor=north west] (r8c3) at (11,-7.4) {\eng{Agenda personnel (Diary / Planner)}};

% Lignes de séparation
\draw[line] (0,-1.1) -- (16.5,-1.1);
\draw[line] (0,-2.0) -- (16.5,-2.0);
\draw[line] (0,-2.9) -- (16.5,-2.9);
\draw[line] (0,-3.8) -- (16.5,-3.8);
\draw[line] (0,-4.7) -- (16.5,-4.7);
\draw[line] (0,-5.6) -- (16.5,-5.6);
\draw[line] (0,-6.5) -- (16.5,-6.5);
\draw[line] (0,-7.4) -- (16.5,-7.4);
\draw[line] (0,-8.3) -- (16.5,-8.3);
\draw[line] (5.5,0) -- (5.5,-8.3);
\draw[line] (11,0) -- (11,-8.3);
\draw[line] (0,0) -- (0,-8.3);
\draw[line] (16.5,0) -- (16.5,-8.3);
\end{tikzpicture}
\legende{Tableau récapitulatif des faux amis majeurs du thème « Participation \& Gestion du temps ».}
\end{popfigure}

\courssub{Methodology: spotting a false friend at the Baccalauréat}

\begin{methode}[Comment déjouer un faux ami en examen (Bac Strategy)]
\textbf{Étape 1 : Repérage visuel.} Vous lisez un mot anglais qui ressemble \textbf{parfaitement} à un mot français (\eng{actually, demand, agenda, assist, eventually, library, sensible, large, etc.}). \textbf{Alerte rouge !}

\textbf{Étape 2 : Vérification contextuelle.} Ne traduisez pas par le mot français qui vient à l'esprit. Regardez la phrase entière :
\begin{itemize}
    \item Temps du verbe ? Sujet ? Compléments ?
    \item \eng{“He \textbf{actually} lives in Bamenda.”} $\rightarrow$ Contraste implicite $\rightarrow$ \textbf{En fait}.
    \item \eng{“The \textbf{agenda} is on the table.”} $\rightarrow$ Réunion $\rightarrow$ \textbf{Ordre du jour}.
\end{itemize}

\textbf{Étape 3 : Substitution mentale.} Remplacez le mot suspect par son \textbf{vrai synonyme anglais} appris en cours.
\begin{itemize}
    \item \eng{Actually} $\rightarrow$ testez \eng{In fact}. Si la phrase a du sens, c'est gagné.
    \item \eng{Demand} $\rightarrow$ testez \eng{Ask for}. Si le ton est poli, \eng{demand} est faux.
\end{itemize}

\textbf{Étape 4 : Production écrite/orale.} \textbf{Interdisez-vous} d'utiliser un mot “transparent” si vous n'êtes pas certain à 100\% de son sens anglais. Utilisez le synonyme sûr (\eng{currently} au lieu de \eng{actually} pour “actuellement” ; \eng{ask for} au lieu de \eng{demand} pour “demander”).
\end{methode}

\courssub{Pronunciation and spelling traps}

\begin{aretenir}[Prononciation : pièges sonores]
\begin{itemize}
    \item \eng{Actually} : « ak-tchou-eu-li » (4 syllabes) — \textbf{pas} « ak-tu-el-le-ment ».
    \item \eng{Eventually} : « i-ven-tchou-eu-li » (5 syllabes) — accent sur \textbf{ven}.
    \item \eng{Agenda} : « a-djen-da » (3 syllabes) — \textbf{pas} « a-gen-da » à la française.
    \item \eng{Delay} (n/v) : « dé-lèï » — diphtongue \eng{/eɪ/}, pas « dé-lè ».
    \item \eng{Busy} : « \textbf{bi-zi} » (u = son \eng{/ɪ/}), pas « bu-zi ».
\end{itemize}
\end{aretenir}

\begin{aretenir}[Orthographe : doublements et suffixes]
\begin{itemize}
    \item \eng{Occupied} (double \eng{c}, double \eng{p}, \eng{ie} avant \eng{d}) $\leftarrow$ \eng{Occupy}.
    \item \eng{Participation} (\eng{ti} + \eng{on}) — attention au \eng{t} (pas \eng{ss} comme en fr. \emph{participation}).
    \item \eng{Management} (garde \eng{e} muet avant \eng{ment} : \eng{manage} + \eng{ment}) — règle des verbes en \eng{-ge}/\eng{-ce} (\eng{judgment} US / \eng{judgement} UK, mais \eng{management} toujours avec \eng{e}).
\end{itemize}
\end{aretenir}

\courssub{Solved exercise: translation and correction}

\begin{exoresolu}[Correction de calques et choix lexicaux — Type Bac]
\textbf{Énoncé :} Traduisez les phrases suivantes en anglais naturel. Évitez les calques et les faux amis.

1. Le directeur a assisté à la réunion hier. (Il était présent.)
2. Actuellement, nous préparons le budget pour l'année prochaine.
3. Les élèves ont finalement compris la leçon après l'explication.
4. Le train a eu deux heures de retard à cause de la pluie.
5. Je dois programmer mon emploi du temps pour la semaine.
6. Elle a exigé de voir le responsable immédiatement. (Contexte : plainte formelle, ton fort.)
7. Notez ce rendez-vous dans votre agenda.
8. Je participe à ce projet depuis janvier.

\textbf{Solution détaillée :}

1. \eng{The Director \textbf{attended} the meeting yesterday.}
   \textit{(Piège : “assisté” $\neq$ \eng{assisted}. \eng{Assisted} = a aidé. Ici, présence physique = \eng{attend}.)}

2. \eng{\textbf{Currently} / \textbf{At present}, we are preparing the budget for next year.}
   \textit{(Piège : “Actuellement” $\neq$ \eng{Actually}. \eng{Actually} = “En fait”.)}

3. \eng{The students \textbf{eventually} understood the lesson after the explanation.}
   \textit{(Piège : “Finalement” (après un processus) = \eng{Eventually}. Pas \eng{Finally} seul qui marque juste la fin d'une liste, ni \eng{Eventually} pour “éventuellement”.)}

4. \eng{The train \textbf{was delayed} by two hours because of the rain.} \quad ou \quad \eng{The train had a two-hour \textbf{delay}.}
   \textit{(Piège : “Retard” (incident) = \eng{Delay}. Pas \eng{deadline} qui est la date limite prévue.)}

5. \eng{I need to \textbf{plan my schedule} / \textbf{organise my timetable} for the week.}
   \textit{(Piège : “Programmer” $\neq$ \eng{programme}. \eng{Programme} = coder, ou établir le programme d'un spectacle. Pour son temps : \eng{plan / manage / organise}.)}

6. \eng{She \textbf{demanded} to see the manager immediately.}
   \textit{(Ici, “exigé” (ton fort) = \eng{Demand}. Si c'était une demande polie : \eng{She asked to see...} ou \eng{She requested to see...}.)}

7. \eng{Write this appointment in your \textbf{diary} (UK) / \textbf{planner} (US).}
   \textit{(Piège : “Agenda” (objet personnel) $\neq$ \eng{Agenda}. \eng{Agenda} = ordre du jour d'une réunion.)}

8. \eng{I \textbf{have been participating in} this project since January.}
   \textit{(Pièges : “Participe à” = \eng{Participate IN} (pas \eng{to}). Durée depuis le passé = \eng{Present Perfect Continuous}.)}
\end{exoresolu}

\courssub{Pedagogical alerts: the examiner's red pen}

\begin{attention}[Erreurs fatales au Bac — “Zero Tolerance” Zone]
\textbf{1. “I will assist to the ceremony.”} $\rightarrow$ \textbf{FAUX.}
\begin{itemize}
    \item Correction : \eng{“I will \textbf{attend} the ceremony.”}
    \item Explication : \eng{Assist} = aider. \eng{Assist to} n'existe pas.
\end{itemize}

\textbf{2. “Actually, I am in Douala.”} (pour dire “Actuellement, je suis à Douala”) $\rightarrow$ \textbf{FAUX.}
\begin{itemize}
    \item Correction : \eng{“\textbf{Currently} / \textbf{At present}, I am in Douala.”}
    \item \eng{Actually, I am in Douala} signifie : “En fait / Contrairement à ce que tu penses, je suis à Douala.”
\end{itemize}

\textbf{3. “We will eventually do it if we have time.”} $\rightarrow$ \textbf{FAUX (sens inversé).}
\begin{itemize}
    \item Correction : \eng{“We will do it \textbf{if necessary} / \textbf{possibly}.”}
    \item \eng{Eventually} implique la certitude finale après attente, pas l'option.
\end{itemize}

\textbf{4. “The deadline of the train is 10 minutes.”} $\rightarrow$ \textbf{FAUX.}
\begin{itemize}
    \item Correction : \eng{“The train has a 10-minute \textbf{delay}.”} ou \eng{“The train is 10 minutes \textbf{late}.”}
    \item \eng{Deadline} = date limite de dépôt (inscription, projet), pas retard de transport.
\end{itemize}

\textbf{5. “I participated to the cleaning campaign.”} $\rightarrow$ \textbf{FAUX.}
\begin{itemize}
    \item Correction : \eng{“I \textbf{participated in} the cleaning campaign.”}
    \item Règle immuable : \eng{Participate IN} + V-ing / Nom.
\end{itemize}

\textbf{6. “I am very occupied this week.”} $\rightarrow$ \textbf{MALADROIT / PEU NATUREL.}
\begin{itemize}
    \item Correction : \eng{“I am very \textbf{busy} this week.”} ou \eng{“I have a \textbf{busy} week ahead.”}
\end{itemize}
\end{attention}

\courssub{Cultural insight: the diary vs the agenda}

\begin{savaistu}[Agenda, Diary, Planner : une histoire de réunions]
Le mot \eng{agenda} vient du latin \emph{agenda} (“choses à faire”, gérondif de \emph{agere}). Au XVIIe siècle, en anglais, il désignait une liste de choses à faire \textbf{collectivement} lors d'une assemblée (ordre du jour).
Au XIXe siècle, avec l'essor du commerce et de la bureaucratie, les Britanniques ont gardé \eng{agenda} pour la réunion, mais ont adopté \eng{diary} (de \emph{diarium}, journal quotidien) pour le carnet personnel de rendez-vous.
Aux États-Unis, on a préféré \eng{planner} (celui qui planifie), plus “actif” et professionnel.
\textbf{Au Cameroun (contexte bilingue) :} Vous entendez souvent “Donne-moi ton agenda” (franglais) pour “Give me your diary”. À l'écrit officiel (anglais), séparez strictement :
\begin{itemize}
    \item \eng{The \textbf{agenda} of the PTA meeting} (L'ordre du jour de la réunion APE).
    \item \eng{My school \textbf{diary} / \textbf{planner}} (Mon agenda scolaire / carnet de correspondance).
\end{itemize}
\end{savaistu}

\courssub{Consolidation exercise}

\begin{exercice}[Entraînement autonome : False Friends \& Calques]
\textbf{A. Choisissez le bon mot (False Friends).}
1. The meeting was \_\_\_\_\_\_ cancelled due to lack of quorum. (\eng{actually / eventually})
2. The Minister is \_\_\_\_\_\_ visiting the region. (\eng{actually / currently})
3. We must respect the \_\_\_\_\_\_ for submission. (\eng{delay / deadline})
4. Please \_\_\_\_\_\_ the form to the secretary. (\eng{demand / submit / hand in})
5. Check the \_\_\_\_\_\_ for today's meeting. (\eng{agenda / diary})

\textbf{B. Corrigez les erreurs (French Calques).}
1. I will participate to the workshop next week.
2. She is occupied with her new baby.
3. We need to programme our time better.
4. He assisted to the conference in Yaoundé.
5. I demand you to help me with this exercise. (Contexte : demande polie à un camarade).

\textbf{C. Traduisez en anglais.}
1. En fait, je n'aime pas le café.
2. Finalement, il a réussi son examen.
3. Le bus a 20 minutes de retard.
4. Note ce devoir dans ton agenda.
5. Elle a aidé le vieux monsieur à traverser la rue.
\end{exercice}

\begin{corrige}[Exercice : False Friends \& Calques]
\textbf{A.}
1. \eng{eventually} (Finalement, après processus).
2. \eng{currently} (Actuellement).
3. \eng{deadline} (Délai/Date limite).
4. \eng{submit} ou \eng{hand in} (\eng{Demand} = exiger ; \eng{fill in} = remplir).
5. \eng{agenda} (Ordre du jour).

\textbf{B.}
1. \eng{I will participate \textbf{in} the workshop...}
2. \eng{She is \textbf{busy} with her new baby.} (ou \eng{tied up}).
3. \eng{We need to \textbf{manage} / \textbf{plan} our time better.}
4. \eng{He \textbf{attended} the conference in Yaoundé.}
5. \eng{Could you \textbf{help me} / \textbf{assist me} with this exercise?} (\eng{Demand} trop fort ; \eng{Ask you to help} possible mais \eng{help/assist} direct plus naturel).

\textbf{C.}
1. \eng{\textbf{Actually}, I don't like coffee.}
2. \eng{He \textbf{eventually} passed his exam.}
3. \eng{The bus is 20 minutes \textbf{late} / has a 20-minute \textbf{delay}.}
4. \eng{Write this homework in your \textbf{diary} / \textbf{planner}.}
5. \eng{She \textbf{assisted} the old man across the street.} (ou \eng{helped}).
\end{corrige}

\coursec{Word families, pronunciation and spelling}

Dans la section précédente, nous avons débusqué les \cle{faux amis} et les calques du français qui menacent votre expression. Nous allons maintenant franchir une étape décisive : apprendre à \textbf{fabriquer des mots}. En anglais comme en français, un même mot-racine peut donner naissance à toute une famille : verbe, nom, adjectif, adverbe. Maîtriser ces familles, c'est multiplier votre vocabulaire par quatre sans effort supplémentaire — et surtout, c'est pouvoir \textbf{varier le style} de vos copies de Bac, ce que les correcteurs récompensent toujours.

\courssub{Observation : one root, many words}

Lisons d'abord ce court texte, puis observons les mots en gras.

\begin{textebox}[A school clean-up campaign in Bamenda]
Last term, the students of Government High School Bamenda decided to \textbf{organise} a clean-up campaign around the market. The \textbf{organisation} of the event was not easy: the \textbf{organisers} had to contact the council, borrow tools and fix a date. About sixty \textbf{participants} took part, most of them \textbf{volunteers} from Form Five and Lower Sixth. Their \textbf{participation} was entirely \textbf{voluntary}: nobody was paid, and everyone worked \textbf{voluntarily} on Saturday morning. The project \textbf{manager}, a disciplined prefect, insisted on \textbf{punctuality}: everybody had to arrive \textbf{punctually} at seven o'clock, because the \textbf{deadline} had been fixed by the council. Thanks to good time \textbf{management}, the work was finished before noon. The mayor congratulated the school and said the campaign was a model of \textbf{participatory} citizenship.
\end{textebox}

\textbf{Glossaire :} \emph{clean-up campaign} (nom, campagne de nettoyage) ; \emph{council} (nom, conseil municipal) ; \emph{to borrow} (verbe, emprunter) ; \emph{prefect} (nom, élève responsable, « préfet de classe ») ; \emph{deadline} (nom, date limite) ; \emph{mayor} (nom, maire).

\textbf{Questions de compréhension :}
\begin{enumerate}
\item Who took part in the campaign?
\item Was the participation compulsory or voluntary?
\item Why was punctuality important?
\end{enumerate}

\textbf{Observation.} Soulignons les mots de la même famille : \eng{organise} (verbe), \eng{organisation} (nom de l'action), \eng{organisers} (nom des personnes). De même : \eng{participants} (personnes) et \eng{participation} (action) viennent du verbe \eng{participate}. L'anglais, comme le français, construit des familles de mots à l'aide de \cle{suffixes}. Repérer le suffixe d'un mot inconnu permet souvent de deviner sa nature grammaticale — et donc sa place dans la phrase.

\courssub{The great word families of this lesson}

\begin{propriete}[Le suffixe -tion : du verbe au nom d'action]
De nombreux verbes anglais se transforment en \textbf{noms d'action} en ajoutant le suffixe \cle{-tion} (souvent avec une petite modification orthographique) :
\begin{itemize}
\item \eng{participate} $\rightarrow$ \eng{participation} (le -e final tombe)
\item \eng{organise} $\rightarrow$ \eng{organisation} (le -e final tombe)
\item \eng{contribute} $\rightarrow$ \eng{contribution} ; \eng{elect} $\rightarrow$ \eng{election}
\end{itemize}
Autre suffixe fréquent pour le nom d'action : \cle{-ment}, qui s'ajoute sans rien enlever : \eng{postpone} $\rightarrow$ \eng{postponement} (le report), \eng{manage} $\rightarrow$ \eng{management}, \eng{appoint} $\rightarrow$ \eng{appointment}.
\end{propriete}

\begin{center}
\begin{tabularx}{\linewidth}{|X|X|X|X|}\hline
\rowcolor{popblueL} \textbf{Verbe} & \textbf{Nom (action)} & \textbf{Nom (personne)} & \textbf{Adjectif / adverbe} \\ \hline
participate (participer) & participation & participant (participant) & participatory (participatif) \\ \hline
organise (organiser) & organisation & organiser (organisateur) & organised / disorganised (organisé / désorganisé) \\ \hline
volunteer (se porter volontaire) & volunteering (volontariat) & volunteer (volontaire) & voluntary (volontaire) / voluntarily (volontairement) \\ \hline
manage (gérer) & management (gestion) & manager (responsable) & manageable (gérable) \\ \hline
--- & punctuality (ponctualité) & --- & punctual (ponctuel) / punctually (ponctuellement) \\ \hline
\end{tabularx}
\end{center}

\begin{exemplebox}[Exemples traduits]
\eng{Her participation was remarkable.} « Sa participation était remarquable. »\par
\eng{The organisers postponed the event.} « Les organisateurs ont reporté l'événement. »\par
\eng{Punctuality is essential at work.} « La ponctualité est essentielle au travail. »\par
\eng{He works as a volunteer at the hospital.} « Il travaille comme volontaire à l'hôpital. »\par
\eng{The workload is manageable if you plan your time.} « La charge de travail est gérable si tu planifies ton temps. »
\end{exemplebox}

\begin{popfigure}
\begin{tikzpicture}[level distance=1.9cm,
  level 1/.style={sibling distance=3.6cm},
  every node/.style={align=center, font=\small},
  edge from parent/.style={draw=popblue, thick}]
\node[draw=popdark, fill=popblueL, rounded corners, font=\bfseries] {ORGANISE (v.)}
  child { node[draw=popblue, fill=popblue!15, rounded corners, align=center]{organisation (n.)\\ \emph{The organisation was perfect.}} }
  child { node[draw=popgreen, fill=popgreenL, rounded corners, align=center]{organiser (n.)\\ \emph{She is the main organiser.}} }
  child { node[draw=poporange, fill=poporangeL, rounded corners, align=center]{organised (adj.)\\ \emph{He is very organised.}} }
  child { node[draw=poppink, fill=poppinkL, rounded corners, align=center]{disorganised (adj.)\\ \emph{A disorganised meeting.}} };
\end{tikzpicture}
\legende{L'arbre de famille du verbe \eng{organise} : une racine, quatre branches, un exemple par branche.}
\end{popfigure}

\begin{methode}[Enrichir une copie de Bac : remplacer le verbe par le nom]
Pour obtenir un style plus soutenu, transformez un verbe simple en nom de la même famille :
\begin{enumerate}
\item Repérez un verbe courant dans votre phrase : \eng{We organised a campaign.}
\item Transformez-le en nom d'action : \eng{organisation}.
\item Reconstruisez la phrase autour du nom : \eng{The organisation of the campaign required careful planning.}
\end{enumerate}
Autres modèles : \eng{Students participated actively.} $\rightarrow$ \eng{The participation of students was remarkable.} ; \eng{They managed their time well.} $\rightarrow$ \eng{Their time management was excellent.} Ce procédé, appelé \cle{nominalisation}, est très apprécié des correcteurs.
\end{methode}

\begin{savaistu}[Le suffixe -ee : celui qui subit l'action]
En anglais, le suffixe \cle{-ee} désigne la personne qui \emph{subit} l'action, tandis que \cle{-er}/\cle{-or} désigne celui qui la \emph{fait} : \eng{employer} (employeur) / \eng{employee} (employé) ; \eng{trainer} (formateur) / \eng{trainee} (stagiaire) ; \eng{interviewer} / \eng{interviewee} (la personne interviewée). Très utile pour le vocabulaire des métiers ! Remarquez la prononciation : l'accent tombe sur la dernière syllabe — \eng{employee} se dit « èm-plo-YI ».
\end{savaistu}

\courssub{Prononciation : les mots pièges de la leçon}

\begin{center}
\begin{tabularx}{\linewidth}{|X|X|X|}\hline
\rowcolor{popblueL} \textbf{Mot} & \textbf{Prononciation (lettres françaises)} & \textbf{Traduction} \\ \hline
schedule (UK) & « ché-dioul » & emploi du temps, programme \\ \hline
schedule (US) & « ské-dioul » & idem \\ \hline
volunteer & « vo-leun-TIER » & volontaire ; se porter volontaire \\ \hline
punctual & « PEUNK-tchouel » & ponctuel \\ \hline
deadline & « DÈDE-laïne » & date limite \\ \hline
participate & « par-TI-ci-peite » & participer \\ \hline
\end{tabularx}
\end{center}

\begin{attention}[Schedule : deux prononciations correctes]
\eng{Schedule} se prononce « chédioul » au Royaume-Uni et « skédioul » aux États-Unis. \textbf{Les deux sont correctes} : choisissez-en une et soyez cohérent dans tout votre exposé. Autres pièges : \eng{volunteer} porte l'accent sur la dernière syllabe (pas « VO-lon-teur » à la française), et \eng{deadline} se dit « dède-laïne » — ne prononcez pas le -a- comme dans « dead » isolé.
\end{attention}

\courssub{Spelling : double letters and one-word traps}

\begin{propriete}[Orthographe : les pièges à mémoriser]
\begin{itemize}
\item \eng{committee} : \textbf{deux m, deux t, deux e} (comité) — le champion des fautes d'orthographe.
\item \eng{occurrence} : \textbf{deux c, deux r} (événement, fait).
\item \eng{appointment} : \textbf{deux p} (rendez-vous, nomination).
\item \eng{timetable} : \textbf{un seul mot}, sans trait d'union (emploi du temps).
\item \eng{deadline} : \textbf{un seul mot} (date limite).
\end{itemize}
\end{propriete}

\begin{attention}[Erreurs fréquentes des francophones]
N'écrivez jamais \emph{comittee}, \emph{apointment} ou \emph{occurence}. Ne séparez jamais \eng{time table} ni \eng{dead line} en deux mots. Et attention au calque : un \eng{appointment} est un rendez-vous, pas une « pointe » ; on dit \eng{to make an appointment} (prendre rendez-vous), jamais \emph{to take a rendez-vous}. Enfin, dans ce cours nous adoptons l'orthographe britannique : \eng{organise}, \eng{organised} (l'orthographe américaine \eng{organize} existe, mais ne mélangez pas les deux dans une même copie).
\end{attention}

\begin{exoresolu}[Word families in action]
\textbf{Énoncé.} Complétez chaque phrase avec la forme correcte du mot entre parenthèses.\par
1. The \dots\ of the festival took six months. (organise)\par
2. All \dots\ must register before Friday. (participate)\par
3. She works \dots\ for a charity in Douala. (voluntary)\par
4. \dots\ is the first quality of a good employee. (punctual)\par
5. The \dots\ of the project meets the staff every Monday. (manage)
\tcblower
\textbf{Solution détaillée.}\par
1. \eng{organisation} — il faut un nom d'action après l'article \eng{the} : « L'organisation du festival a pris six mois. »\par
2. \eng{participants} — il faut un nom de personne au pluriel (sujet de \eng{must register}) : « Tous les participants doivent s'inscrire avant vendredi. »\par
3. \eng{voluntarily} — il faut un adverbe qui modifie le verbe \eng{works} : « Elle travaille bénévolement pour une œuvre caritative à Douala. »\par
4. \eng{Punctuality} — il faut un nom abstrait en position de sujet : « La ponctualité est la première qualité d'un bon employé. »\par
5. \eng{manager} — il faut le nom de la personne (celle qui rencontre l'équipe) : « Le responsable du projet rencontre l'équipe chaque lundi. »
\end{exoresolu}

\begin{exercice}[À vous de jouer]
1. Construisez l'arbre de famille complet de \eng{participate} et de \eng{manage}, avec un exemple par branche.\par
2. Réécrivez en style soutenu : \eng{They organised the meeting well.} / \eng{Many students participated.}\par
3. Dictez-vous mentalement et écrivez sans faute : committee, occurrence, appointment, timetable, deadline.
\end{exercice}

\begin{corrige}[Exercice — éléments de réponse]
1. PARTICIPATE $\rightarrow$ participation (n.), participant (n.), participatory (adj.) ; MANAGE $\rightarrow$ management (n.), manager (n.), manageable (adj.). Exemples libres, vérifier la nature et la place du mot dans la phrase.\par
2. \eng{The organisation of the meeting was excellent.} / \eng{The participation of many students was remarkable.}\par
3. Vérifiez : deux m et deux t pour \eng{committee} ; deux c et deux r pour \eng{occurrence} ; deux p pour \eng{appointment} ; un seul mot pour \eng{timetable} et \eng{deadline}.
\end{corrige}

\begin{aretenir}[L'essentiel de la section]
Une racine = une famille : verbe, nom d'action (\cle{-tion}/\cle{-ment}), nom de personne (\cle{-er}/\cle{-or}/\cle{-ee}), adjectif, adverbe. Prononcez \eng{volunteer} « vo-leun-TIER », \eng{deadline} « dède-laïne », et choisissez votre prononciation de \eng{schedule}. Orthographe blindée : \eng{committee}, \eng{occurrence}, \eng{appointment}, \eng{timetable}, \eng{deadline}. Enfin, pour briller au Bac, nominalisez : transformez vos verbes en noms d'action.
\end{aretenir}

\coursec{Using the lexis in context}

Après avoir étudié les familles de mots, la prononciation et l'orthographe du vocabulaire de la participation et de la gestion du temps, nous passons maintenant à l'étape cruciale : \textbf{la mise en contexte}. C'est dans des phrases authentiques, des dialogues réalistes ou des textes cohérents que le lexique prend tout son sens. Cette section vous propose un texte original ancré dans le contexte camerounais, des activités de compréhension guidées, une méthode de lecture efficace et une production écrite pour ancrer durablement les acquis.

\courssub{Reading comprehension: A Busy Week at GHS Limbe}

Le texte suivant mobilise l'essentiel du vocabulaire vu dans les sections précédentes (participation, gestion du temps, collocations, idiomes). Lisez-le attentivement en appliquant la méthode qui sera présentée par la suite.

\begin{textebox}[A Busy Week at GHS Limbe — Article original]
On Monday morning, the principal called a special meeting in the main hall to organise the annual fund-raising day scheduled for Saturday. Thirty enthusiastic students signed up immediately to join the various committees. The deadline for the poster competition was set for Wednesday afternoon, but because of heavy rain, the submission date was put off until Friday. By Thursday, the organising committee realised they were seriously behind schedule with the decoration of the school grounds. Consequently, the head boy made an urgent appeal over the intercom asking every available student to lend a hand. Thanks to those who pulled their weight and stayed late after classes, the event took place on time on Saturday morning and was a huge success.
\end{textebox}

\begin{definition}[Glossaire du texte — Mots et expressions clés]
\begin{center}
\begin{tabularx}{\linewidth}{|X|X|X|}
\hline
\rowcolor{popblueL} \textbf{English} & \textbf{Français} & \textbf{Catégorie / Note} \\
\hline
to sign up & s'inscrire, se porter volontaire & Verbe à particule (phrasal verb) — \eng{sign up for} \\
\hline
fund-raising (n) & collecte de fonds, activité caritative & Nom composé (trait d'union) — \eng{a fund-raising event} \\
\hline
deadline (n) & date limite, échéance & \eng{meet a deadline} (respecter), \eng{miss a deadline} (manquer) \\
\hline
to put off & reporter, remettre à plus tard & Séparable : \eng{put it off} / \eng{put off the meeting} \\
\hline
behind schedule (adj) & en retard (sur le planning) & Antonyme : \eng{ahead of schedule} (en avance), \eng{on schedule} (dans les temps) \\
\hline
to lend a hand & donner un coup de main, aider & Idiome — \eng{lend a hand with sth / to sb} \\
\hline
to pull one's weight & faire sa part du travail, assumer ses responsabilités & Idiome (souvent au passé \eng{pulled their weight}) \\
\hline
head boy / head girl (n) & délégué principal / déléguée principale & Voir définition détaillée ci-dessous \\
\hline
\end{tabularx}
\end{center}
\end{definition}

\begin{definition}[Head boy / Head girl — Contexte culturel anglophone]
Dans le système éducatif anglophone (Cameroun, Royaume-Uni, Nigeria, Ghana, etc.), le \textbf{head boy} (garçon) ou la \textbf{head girl} (fille) est un élève de la classe terminale (Upper Sixth), élu par ses pairs et/ou désigné par l'administration, pour représenter l'ensemble du corps étudiant. Il/Elle fait le lien entre l'administration (principal, discipline masters) et les élèves, coordonne les préfets (prefects) et prend la parole lors des cérémonies officielles. C'est l'équivalent du \emph{délégué principal} ou du \emph{président du bureau des élèves} dans le système francophone, mais avec souvent plus de responsabilités protocolaires et disciplinaires.
\end{definition}

\courssub{Comprehension activities}

\begin{exercice}[Compréhension du texte — Bac Blanc]
\textbf{Instructions :} Répondez aux questions en vous basant \textbf{uniquement} sur le texte. Pour les questions Vrai/Faux, citez la phrase ou le segment de texte qui justifie votre réponse (numéro de ligne ou citation entre guillemets anglais “ ”).

\textbf{A. True or False? Justify by quoting the text.}
\begin{enumerate}
\item The meeting organised by the principal took place on Tuesday morning.
\item The deadline for the poster competition was postponed because of bad weather.
\item The head boy asked students to help on Wednesday.
\end{enumerate}

\textbf{B. Multiple Choice Questions. Choose the correct answer (a, b or c).}
\begin{enumerate}
\setcounter{enumi}{3}
\item In the text, “signed up” means:
\begin{itemize}
\item[a)] registered as volunteers
\item[b)] signed a contract
\item[c)] arrived on time
\end{itemize}
\item The expression “behind schedule” means:
\begin{itemize}
\item[a)] ahead of time
\item[b)] late compared to the plan
\item[c)] exactly on time
\end{itemize}
\end{enumerate}

\textbf{C. Short Answer Questions. Answer in complete sentences.}
\begin{enumerate}
\setcounter{enumi}{5}
\item Why was the submission date for the poster competition put off?
\item Who made an appeal for help over the intercom?
\item When did the fund-raising day finally take place?
\end{enumerate}
\end{exercice}

\begin{exoresolu}[Corrigé détaillé de l'exercice de compréhension]
\tcblower
\textbf{A. True or False?}
\begin{enumerate}
\item \textbf{False.} Justification : “On \textbf{Monday} morning, the principal called a special meeting…”
\item \textbf{True.} Justification : “The deadline for the poster competition was set for Wednesday afternoon, but because of \textbf{heavy rain}, the submission date was put off until Friday.”
\item \textbf{False.} Justification : “By \textbf{Thursday}, the organising committee realised they were seriously behind schedule… Consequently, the head boy made an urgent appeal…”
\end{enumerate}

\textbf{B. Multiple Choice Questions.}
\begin{enumerate}
\setcounter{enumi}{3}
\item \textbf{a) registered as volunteers.} Contexte : “Thirty enthusiastic students signed up immediately to join the various committees.”
\item \textbf{b) late compared to the plan.} Contexte : “By Thursday… they were seriously behind schedule with the decoration…”
\end{enumerate}

\textbf{C. Short Answer Questions.}
\begin{enumerate}
\setcounter{enumi}{5}
\item The submission date was put off \textbf{because of heavy rain}.
\item \textbf{The head boy} made an urgent appeal over the intercom asking every available student to lend a hand.
\item The event took place \textbf{on Saturday morning}.
\end{enumerate}
\end{exoresolu}

\courssub{Reading strategy: Spotting time markers and action verbs}

Pour réussir une compréhension de texte (reading comprehension) au Baccalauréat ou en classe, il ne suffit pas de lire : il faut \textbf{naviguer} dans le texte. La méthode ci-dessous, testée et approuvée, vous permet de reconstruire la chronologie et d'identifier les informations clés en un temps record.

\begin{methode}[Méthode de lecture active — 5 étapes pour maîtriser un texte narratif]
\begin{enumerate}
\item \textbf{Lisez le titre et la source.} Ils donnent le thème (fund-raising), le lieu (GHS Limbe) et le genre (récit chronologique).
\item \textbf{Repérez les marqueurs de temps (Time markers).} Surlignez ou entourez : \eng{On Monday morning}, \eng{scheduled for Saturday}, \eng{Wednesday afternoon}, \eng{until Friday}, \eng{By Thursday}, \eng{on Saturday morning}. Ils structurent la ligne du temps.
\item \textbf{Entourez les verbes d'action principaux (Action verbs).} Identifiez ce qui se passe : \eng{called}, \eng{signed up}, \eng{was set}, \eng{was put off}, \eng{realised}, \eng{made an appeal}, \eng{asked}, \eng{lend a hand}, \eng{pulled their weight}, \eng{took place}. Suivez le fil de l'histoire.
\item \textbf{Repérez les connecteurs logiques.} \eng{But} (contraste), \eng{Consequently} (conséquence), \eng{Thanks to} (cause/résultat positif). Ils expliquent les liens entre les événements.
\item \textbf{Répondez aux questions en citant le texte.} Ne répondez jamais de mémoire. Retournez au paragraphe concerné, copiez la preuve exacte entre guillemets anglais “ ”, puis formulez votre réponse en anglais correct.
\end{enumerate}
\end{methode}

\begin{popfigure}
\begin{tikzpicture}[
    node distance=1.2cm and 0.5cm,
    every node/.style={font=\small, align=center},
    stepS/.style={rectangle, rounded corners, draw=popblue, thick, fill=popblueL, minimum width=2.8cm, minimum height=1.2cm},
    arrow/.style={-Stealth, thick, popdark}
]
\node[stepS, align=center] (s1){1. Lire le titre\\ \& la source};
\node[stepS, right=of s1, align=center] (s2){2. Repérer les\\ marqueurs de temps};
\node[stepS, right=of s2, align=center] (s3){3. Entourer les\\ verbes d'action};
\node[stepS, right=of s3, align=center] (s4){4. Identifier les\\ connecteurs logiques};
\node[stepS, right=of s4, align=center] (s5){5. Répondre en\\ citant le texte};

\draw[arrow] (s1) -- (s2);
\draw[arrow] (s2) -- (s3);
\draw[arrow] (s3) -- (s4);
\draw[arrow] (s4) -- (s5);

% Annotations sous les boîtes
\node[below=0.3cm of s1, text=popmuted, font=\tiny] {Theme, lieu, genre};
\node[below=0.3cm of s2, text=popmuted, font=\tiny, align=center]{Monday, Wednesday,\\ Friday, Saturday};
\node[below=0.3cm of s3, text=popmuted, font=\tiny, align=center]{called, signed up,\\ put off, took place};
\node[below=0.3cm of s4, text=popmuted, font=\tiny, align=center]{But, Consequently,\\ Thanks to};
\node[below=0.3cm of s5, text=popmuted, font=\tiny, align=center]{Preuve textuelle \\ + phrase complète};
\end{tikzpicture}
\legende{Méthode de lecture active : les 5 étapes pour une compréhension efficace}
\end{popfigure}

\courssub{Classroom practice: Lexical fields in the text}

En classe, l'enseignant guide les élèves pour opérationnaliser le vocabulaire. Voici un exemple type d'activité menée collectivement au tableau.

\begin{exemplebox}[Activité guidée : Identifier les champs lexicaux dans le texte]
\textbf{Consigne :} Retrouvez dans le texte \textbf{5 mots ou expressions du champ du temps / de la gestion du temps} et \textbf{5 mots ou expressions du champ de la participation / de l'action collective}.

\textbf{Résultat attendu au tableau :}

\begin{center}
\begin{tabularx}{\linewidth}{|X|X|}
\hline
\rowcolor{poporangeL} \textbf{Champ lexical : TIME MANAGEMENT (Gestion du temps)} & \textbf{Champ lexical : PARTICIPATION / ACTION (Participation)} \\
\hline
\begin{itemize}
\item \eng{scheduled for Saturday} (prévu pour)
\item \eng{deadline} (date limite)
\item \eng{Wednesday afternoon / Friday} (repères calendaires)
\item \eng{put off until Friday} (reporter)
\item \eng{behind schedule} (en retard)
\item \eng{By Thursday} (échéance)
\item \eng{on time} (à l'heure / dans les temps)
\item \eng{stayed late} (rester tard)
\end{itemize}
&
\begin{itemize}
\item \eng{signed up} (s'inscrire / se porter volontaire)
\item \eng{join the various committees} (rejoindre les comités)
\item \eng{organising committee} (comité d'organisation)
\item \eng{lend a hand} (donner un coup de main)
\item \eng{pulled their weight} (faire sa part)
\item \eng{urgent appeal} (appel urgent)
\item \eng{every available student} (tout élève disponible)
\item \eng{huge success} (grand succès)
\end{itemize}
\\
\hline
\end{tabularx}
\end{center}
\textbf{Observation pédagogique :} Remarquez que \eng{deadline}, \eng{schedule}, \eng{on time} sont des termes techniques de gestion de projet. \eng{Lend a hand} et \eng{pull one's weight} sont des idiomes figés (ne pas traduire mot à mot).
\end{exemplebox}

\begin{methode}[Comment répondre aux questions de compréhension (Type Bac) — Règle d'or]
\begin{enumerate}
\item \textbf{Lisez la question} et identifiez le mot-clé (Who, When, Why, Which expression…).
\item \textbf{Scannez le texte} pour localiser le mot-clé ou son synonyme (étape 2 et 3 de la méthode ci-dessus).
\item \textbf{Copiez la justification exacte} entre guillemets anglais : “ … ”.
\item \textbf{Rédigez la réponse} en phrase complète, en reprenant les mots du texte si possible, ou en reformulant simplement.
\item \textbf{Vérifiez} : la réponse vient-elle \textbf{uniquement} du texte ? (Pas de connaissance extérieure).
\end{enumerate}
\emph{Exemple :} \eng{Question: Why was the date put off?} $\rightarrow$ \eng{Justification: “because of heavy rain”.} $\rightarrow$ \eng{Answer: It was put off because of heavy rain.}
\end{methode}

\begin{exemplebox}[Exemples traduits — Phrases pivots du texte]
\begin{enumerate}
\item \eng{The deadline was put off until Friday.} \\
« La date limite a été reportée à vendredi. » (Passif : \eng{was put off} — attention à l'ordre : particule \eng{off} après le verbe, avant \eng{until}).
\item \eng{The committee was behind schedule.} \\
« Le comité était en retard (sur le planning). » (\eng{be} + \eng{behind schedule} = état).
\item \eng{The head boy asked everyone to lend a hand.} \\
« Le délégué principal a demandé à tout le monde de donner un coup de main. » (\eng{ask sb to do sth} + infinitif).
\item \eng{Thanks to those who pulled their weight, the event took place on time.} \\
« Grâce à ceux qui ont fait leur part du travail, l'événement a eu lieu à l'heure prévue. » (\eng{Thanks to} + nom/pronom/\eng{-ing} ; \eng{take place} = avoir lieu, intransitif).
\end{enumerate}
\end{exemplebox}

\begin{attention}[Piège du Bac : True / False — Justify = Points obligatoires]
\textbf{Au Baccalauréat camerounais (série A), les questions “True or False? Justify” sont notées sur 2 points : 1 point pour V/F, 1 point pour la justification.}
\begin{itemize}
\item \textbf{Zéro point} si vous écrivez seulement “True” ou “False” sans citation.
\item \textbf{Zéro point} si la justification est inventée, vague ou en français (“Parce qu'il a plu” au lieu de “because of heavy rain”).
\item \textbf{La citation doit être en anglais}, copiée \textbf{fidèlement} du texte, entre guillemets anglais “ ”.
\item \textbf{Astuce :} Si la phrase est longue, citez le segment pertinent seulement (ex. : “because of heavy rain” au lieu de toute la phrase).
\end{itemize}
\end{attention}

\courssub{Guided production: My weekly schedule}

Maintenant à vous de jouer. L'objectif est de \textbf{réemployer activement} le lexique pour décrire votre propre réalité (école, maison, quartier, église, association de jeunesse).

\begin{exercice}[Production écrite guidée — My Busy Week]
\textbf{Consigne :} Rédigez un court paragraphe (6 à 8 phrases) intitulé “My Busy Week”. Décrivez une semaine type ou un événement récent (préparation de la fête de la jeunesse, interclasse, examen, marché, travail aux champs, répétition chorale…).

\textbf{Contraintes obligatoires :}
\begin{enumerate}
\item Utilisez \textbf{au moins 4 mots ou expressions cibles} de la leçon (ex. : \eng{sign up}, \eng{deadline}, \eng{put off}, \eng{behind schedule}, \eng{lend a hand}, \eng{pull one's weight}, \eng{on time}, \eng{schedule}, \eng{meeting}, \eng{committee}).
\item Intégrez \textbf{au moins 1 idiome} (\eng{lend a hand} ou \eng{pull one's weight} ou \eng{take place}).
\item Respectez la chronologie (marqueurs de temps : \eng{On Monday}, \eng{By Tuesday}, \eng{Finally on Saturday}…).
\item Écrivez en \textbf{anglais correct} (majuscules, ponctuation, accord SV).
\end{enumerate}
\end{exercice}

\begin{exoresolu}[Modèle de production attendue — Élève de Tle A]
\tcblower
\textbf{My Busy Week}

Last week was intense at Government Bilingual High School Yaoundé. On Monday, our form master announced the upcoming \eng{inter-class sports competition} \eng{scheduled for} Friday. I immediately \eng{signed up} for the 400 metres race and the relay team. The \eng{deadline} to submit the team list was Wednesday, but our class representative \eng{put it off} until Thursday morning because some students were absent. By Thursday afternoon, we were \eng{behind schedule} with the training sessions. Our sports prefect asked the senior students to \eng{lend a hand} with the warm-up exercises. Fortunately, everyone \eng{pulled their weight} and we trained hard. The competition \eng{took place on time} on Friday and our class won the trophy!

\textbf{Analyse du modèle :}
\begin{itemize}
\item \textbf{Cibles utilisées (6) :} scheduled for, signed up, deadline, put it off, behind schedule, on time.
\item \textbf{Idiomes (2) :} lend a hand, pulled their weight, took place.
\item \textbf{Marqueurs de temps :} Last week, On Monday, Wednesday, Thursday, By Thursday afternoon, Friday.
\item \textbf{Grammaire :} Prétérit simple (announced, signed, put, was, asked, trained, took, won) — cohérent pour un récit passé.
\end{itemize}
\end{exoresolu}

\courssub{Speaking practice: Pair work — “My Commitments”}

\begin{experience}[Jeu de rôle / Pair Work — Planning a School Event]
\textbf{Situation :} Vous êtes deux membres du \eng{Students' Representative Council} (Conseil des élèves). Vous devez organiser la \eng{End-of-Year Party} (Fête de fin d'année) dans deux semaines.

\textbf{Rôles :}
\begin{itemize}
\item \textbf{Student A (Chairperson / Président) :} Vous menez la réunion. Vous donnez les \eng{deadlines}, vous \eng{assign tasks} (répartissez les tâches), vous vérifiez si l'équipe est \eng{on schedule} ou \eng{behind schedule}.
\item \textbf{Student B (Secretary / Secrétaire) :} Vous prenez des notes. Vous \eng{sign up} des volontaires. Vous proposez de \eng{lend a hand} pour la déco. Vous rappelez à tous de \eng{pull their weight}.
\end{itemize}

\textbf{Script de départ (à compléter oralement) :}
\begin{quote}
\textbf{A:} “Good morning everyone. The party is \eng{scheduled for} Saturday 15th June. The \eng{deadline} for the budget is tomorrow. Are we \eng{on schedule}?”
\textbf{B:} “Not really, Chairperson. We are \eng{behind schedule} with the decorations. Can we \eng{put off} the budget meeting to Thursday?”
\textbf{A:} “No, we can't \eng{put it off}. I need volunteers to \eng{lend a hand} today. Who can \eng{sign up}?”
\textbf{B:} “I'll \eng{sign up}. And I'll ask the Form 5 students to \eng{pull their weight} with the sound system.”
\end{quote}

\textbf{Consigne :} Jouez la scène 3 minutes. Changez de rôle. L'enseignant circule et note l'usage correct du vocabulaire cible (\eng{deadline, behind schedule, put off, sign up, lend a hand, pull one's weight, on schedule, scheduled for}).
\end{experience}

\begin{aretenir}[Synthèse lexicale — Taking Part \& Managing Time]
\begin{center}
\begin{tabularx}{\linewidth}{|X|X|X|}
\hline
\rowcolor{popgreenL} \textbf{Notion} & \textbf{Expressions clés (English)} & \textbf{Équivalent français / Contexte} \\
\hline
\textbf{S'inscrire / Participer} & \eng{sign up (for)}, \eng{join}, \eng{volunteer (for)}, \eng{take part in}, \eng{enroll in} & S'inscrire à une activité, rejoindre un groupe \\
\hline
\textbf{Aider / Contribuer} & \eng{lend a hand (to sb / with sth)}, \eng{pull one's weight}, \eng{chip in}, \eng{contribute to} & Donner un coup de main, faire sa part, participer financièrement/effort \\
\hline
\textbf{Responsabilité / Rôle} & \eng{be in charge of}, \eng{be responsible for}, \eng{head boy/girl}, \eng{prefect}, \eng{committee member}, \eng{chair a meeting} & Être responsable de, délégué, préfet, membre du comité, présider \\
\hline
\textbf{Planifier / Échéances} & \eng{schedule (v/n)}, \eng{deadline}, \eng{set a date}, \eng{meet a deadline}, \eng{miss a deadline}, \eng{timetable}, \eng{agenda} & Planifier, date limite, respecter/manquer l'échéance, emploi du temps, ordre du jour \\
\hline
\textbf{Gérer les retards / changements} & \eng{put off}, \eng{postpone}, \eng{delay (v/n)}, \eng{behind schedule}, \eng{ahead of schedule}, \eng{on schedule / on time}, \eng{run late} & Reporter, retard, en retard/en avance/dans les temps, avoir du retard \\
\hline
\textbf{Réalisation} & \eng{take place}, \eng{hold (an event)}, \eng{organise}, \eng{run smoothly}, \eng{be a success} & Avoir lieu, organiser, se dérouler sans accroc, être un succès \\
\hline
\end{tabularx}
\end{center}

\end{aretenir}

\begin{savaistu}[Le temps, c'est de l'argent ? — Culture anglophone \& Cameroun]
L'expression \eng{“Time is money”} (Benjamin Franklin, 1748) résume la vision \textbf{anglo-saxonne} du temps : une ressource rare, quantifiable, qu'on \eng{spend} (dépense), \eng{save} (économise), \eng{waste} (gaspille), \eng{invest} (investit). Au Cameroun, dans le monde professionnel (banques, administrations, ONG, entreprises du bois/cacao à Douala), la \textbf{ponctualité} (\eng{punctuality}) et le respect des \eng{deadlines} sont des critères d'évaluation majeurs. Cependant, dans la vie communautaire (village, église, \eng{njangi} / tontine, cérémonies traditionnelles), le temps est plus \textbf{élastique} (\eng{flexible time} / \eng{African time}) : l'événement \eng{takes place} quand les personnes importantes sont arrivées. Un bon communicant en anglais doit savoir \textbf{naviguer entre ces deux registres} : \eng{formal/deadline-driven} au travail, \eng{relational/event-driven} dans la communauté.
\end{savaistu}

\newpage

\coursec{Activité d'intégration}

\coursec{Lesson 23 — Vocabulary: Words and expressions related to taking part and managing time in community/workplace/school activities}

\courssub{Objectifs de la leçon}
\noindent Cette leçon d'intégration (type \eng{Task-Based Learning}) vise à mobiliser l'ensemble du vocabulaire et des expressions du thème « \cle{Participation} » et « \cle{Gestion du temps} » dans une situation de communication authentique, contextualisée au Cameroun. À l'issue de cette séquence de 2 h, l'élève doit être capable de :
\begin{itemize}
    \item \textbf{Identifier et réemployer} activement le lexique de la participation (\eng{volunteer, sign up, take part in, contribute to, lend a hand}) et de la gestion du temps (\eng{deadline, schedule, timetable, postpone, on time, in time, beat the clock, time slot, prioritise, allocate, clash, feasible, realistic}) à l'écrit et à l'oral.
    \item \textbf{Produire} un document écrit structuré (\eng{schedule} / programme horaire) et une présentation orale cohérente (2 min) en utilisant \textbf{au minimum huit mots-cibles} de la leçon et \textbf{deux expressions idiomatiques}.
    \item \textbf{Argumenter} un choix en justifiant son vote avec le vocabulaire précis acquis (\eng{realistic, feasible, well-organised, tight schedule, clash, buffer zone}).
    \item \textbf{Éviter} les faux amis et calques du français identifiés (\eng{assist/attend, participate in/participate to, planning/schedule, eventually/finalement, actual/actuel}).
    \item \textbf{Maîtriser} la prononciation des mots-clés (accent tonique, sons difficiles) et l'orthographe britannique (\eng{prioritise, programme}).
\end{itemize}

\courssub{Situation déclenchante \& Document support}
\noindent Le contexte se déroule au \textbf{Lycée Bilingue de Buea}, région du Sud-Ouest. Le Proviseur lance un appel à projets pour organiser la \eng{School Community Day} annuelle (journée de service communautaire, cohésion sociale, collecte de fonds pour la bibliothèque). Vous êtes membres du \eng{Student Council} (Conseil des élèves). Mission : concevoir le programme officiel (7 h 00 – 17 h 00) et le présenter devant l'administration et les délégués pour obtenir approbation et budget.

\begin{textebox}[Memo officiel du Proviseur — Lycée Bilingue de Buea]
\textbf{TO:} All Student Council Members \\
\textbf{FROM:} Principal's Office \\
\textbf{DATE:} 15 October 2025 \\
\textbf{SUBJECT:} Call for Proposals — School Community Day (20 November)

Dear Students,

As part of our civic engagement \textbf{programme}, we will hold our annual \textbf{School Community Day} on \textbf{Thursday, 20 November}. The objective is threefold: \textbf{clean-up} of the campus and neighbouring streets, a \textbf{blood donation drive} in partnership with the Buea Regional Hospital, and a \textbf{fund-raising cultural show} to equip the school library.

We invite each committee to submit a detailed \textbf{schedule (timeline)} for the day (7:00 AM – 5:00 PM). Your proposal must include:
\begin{itemize}
    \item A list of \textbf{volunteers} needed per activity and a \textbf{sign-up} strategy.
    \item Clear \textbf{time slots} for each event, including set-up and clean-up.
    \item A \textbf{deadline} for confirmation of external partners (Red Cross, Sponsors).
    \item A contingency plan if rain forces us to \textbf{postpone} outdoor activities.
\end{itemize}

The best proposal will be adopted officially. The \textbf{deadline for submission} is \textbf{Friday, 24 October, at 3:00 PM sharp}. Late proposals will not be considered. Let's \textbf{beat the clock} and make this day a success!

Yours sincerely,
\textit{Mr. N. Ekosso, Principal}
\end{textebox}

\courssub{Démarche pédagogique \& Consignes}
\noindent Travaillez en groupes de 4 élèves. Répartissez les rôles : \eng{Project Manager} (chef de projet/chronométreur), \eng{Secretary} (rédacteur du programme), \eng{Logistics Officer} (gestion des créneaux/ressources), \eng{Spokesperson} (présentateur oral). Suivez les étapes ci-dessous :

\begin{enumerate}
    \item \textbf{Analyse du document support (10 min) :} Lisez le memo. Surlignez ou listez tous les mots-cibles de la leçon 23 qui y figurent (\eng{deadline, volunteers, sign up, postpone, beat the clock, schedule, time slot, clean-up, fund-raising, drive, partnership, programme}). Vérifiez leur sens dans le contexte à l'aide d'un dictionnaire monolingue si nécessaire.
    \item \textbf{Planification écrite — Rédaction du \eng{Schedule} (25 min) :} Rédigez un mini-programme officiel sur feuille A4 (format tableau horaire). Votre document doit impérativement contenir \textbf{au moins 8 mots-cibles} de la leçon (ex. : \eng{deadline, sign up, volunteers, postpone, on time, time slot, prioritise, allocate, clash, feasible, contingency, buffer zone}). Utilisez des phrases complètes pour les descriptions d'activités. Respectez l'horaire 7 h 00 – 17 h 00.
    \item \textbf{Préparation de la présentation orale (15 min) :} Le \eng{Spokesperson} prépare une intervention de 2 minutes maximum. Il/Elle doit intégrer \textbf{deux expressions idiomatiques} vues en classe (ex. : \eng{count me in, lend a hand, beat the clock, in the nick of time, on the dot}). Les autres membres préparent des arguments pour la phase de questions/réponses (anticiper les objections : \eng{What if it rains? Is the timeline feasible?}).
    \item \textbf{Présentations, évaluation par les pairs et vote (40 min) :} Chaque groupe présente son programme. L'auditoire note sur la grille d'évaluation (respect du vocabulaire, réalisme du planning, clarté orale, utilisation des idiomes, gestion du temps de parole). À la fin, chaque élève vote pour le meilleur projet \textbf{en justifiant son choix en anglais} avec le vocabulaire de la leçon (ex. : \eng{Their schedule is realistic because they allocated enough time for set-up; they didn't forget a contingency plan if we have to postpone.}).
\end{enumerate}

\courssub{Ressources linguistiques mobilisables}
\noindent Rappel synthétique des outils lexicaux, grammaticaux et phonologiques nécessaires pour réussir la tâche.

\begin{definition}[Lexique clé — Participation \eng{(Taking Part)}]
\begin{center}
\begin{tabularx}{\linewidth}{|X|X|X|}\hline
\rowcolor{popblueL} \textbf{Mot / Expression} & \textbf{Classe} & \textbf{Contexte / Exemple (traduction)} \\ \hline
\eng{volunteer} & n. / v. & \eng{We need 20 volunteers for the clean-up.} (Il nous faut 20 bénévoles.) / \eng{I volunteer to lead the team.} (Je me porte volontaire.) \\ \hline
\eng{sign up (for)} & phr. v. & \eng{Please sign up for the blood drive by Friday.} (Inscrivez-vous pour la collecte de sang d'ici vendredi.) \\ \hline
\eng{take part in} & phr. v. & \eng{Every student must take part in at least one activity.} (Chaque élève doit participer à au moins une activité.) \\ \hline
\eng{contribute to} & v. + prep. & \eng{Local businesses contributed to the fund-raising.} (Les commerces locaux ont contribué à la collecte de fonds.) \\ \hline
\eng{lend a hand} & idiom & \eng{Can you lend a hand with the stage set-up?} (Peux-tu donner un coup de main pour l'installation de la scène ?) \\ \hline
\eng{count me in} & idiom & \eng{“Count me in for the cultural show!”} (Compte sur moi pour le spectacle !) \\ \hline
\eng{participate in} & v. + prep. & \textbf{Attention :} pas \eng{*participate to}. \eng{They participated in the drive.} (Ils ont participé à la collecte.) \\ \hline
\eng{attend} & v. & \eng{All delegates must attend the meeting.} (Tous les délégués doivent assister à la réunion.) \\ \hline
\end{tabularx}
\end{center}
\end{definition}

\begin{definition}[Lexique clé — Gestion du temps \eng{(Managing Time)}]
\begin{center}
\begin{tabularx}{\linewidth}{|X|X|X|}\hline
\rowcolor{poporangeL} \textbf{Mot / Expression} & \textbf{Classe} & \textbf{Contexte / Exemple (traduction)} \\ \hline
\eng{deadline} & n. & \eng{The deadline for proposals is 24 October.} (La date limite pour les propositions est le 24 octobre.) \\ \hline
\eng{schedule / timetable} & n. & \eng{The schedule is tight; we must start on time.} (L'emploi du temps est serré ; nous devons commencer à l'heure.) \\ \hline
\eng{time slot} & n. & \eng{We allocated a 30-minute time slot for the speech.} (Nous avons alloué un créneau de 30 min pour le discours.) \\ \hline
\eng{postpone} & v. & \eng{If it rains, we will postpone the clean-up to Saturday.} (S'il pleut, nous reporterons le nettoyage à samedi.) \\ \hline
\eng{on time} & adv. & \eng{The event started on time.} (L'événement a commencé à l'heure / ponctuellement.) \\ \hline
\eng{in time} & adv. & \eng{We arrived in time for the opening.} (Nous sommes arrivés à temps pour l'ouverture / avant le début.) \\ \hline
\eng{beat the clock} & idiom & \eng{We worked hard to beat the clock and finish the report.} (Nous avons travaillé dur pour gagner du temps / finir avant l'échéance.) \\ \hline
\eng{prioritise} & v. & \eng{We must prioritise safety during the blood drive.} (Nous devons donner la priorité à la sécurité pendant la collecte de sang.) \\ \hline
\eng{allocate} & v. & \eng{We allocated two hours for the cultural show.} (Nous avons alloué deux heures au spectacle culturel.) \\ \hline
\eng{clash (with)} & v. & \eng{The football match clashes with the fund-raising.} (Le match de football coïncide / entre en conflit avec la collecte de fonds.) \\ \hline
\eng{feasible / realistic} & adj. & \eng{Is this timeline feasible with only 10 volunteers?} (Ce calendrier est-il réalisable avec seulement 10 bénévoles ?) \\ \hline
\eng{contingency plan} & n. & \eng{We have a contingency plan in case of rain.} (Nous avons un plan de secours en cas de pluie.) \\ \hline
\eng{buffer zone / buffer time} & n. & \eng{We added a 15-minute buffer zone between events.} (Nous avons ajouté une marge de 15 min entre les événements.) \\ \hline
\eng{on standby} & adj. & \eng{Volunteers are on standby for the contingency plan.} (Des bénévoles sont en attente / de réserve pour le plan de secours.) \\ \hline
\end{tabularx}
\end{center}
\end{definition}

\begin{attention}[Faux amis et pièges fréquents — Rappel strict]
\begin{itemize}
    \item \textbf{\eng{Assist}} = aider (\eng{The nurse assisted the doctor.}). Pour « assister à » (être présent) : \textbf{\eng{attend}} ou \textbf{\eng{take part in}}. \eng{*I assisted the meeting} $\rightarrow$ \eng{I attended the meeting.}
    \item \textbf{\eng{Planning}} (n.m. français) $\rightarrow$ en anglais : \textbf{\eng{schedule}}, \textbf{\eng{timetable}} ou \textbf{\eng{plan}}. \eng{Planning} existe mais est rare et technique (\eng{urban planning} = aménagement du territoire). \eng{*Our planning is ready} $\rightarrow$ \eng{Our schedule is ready.}
    \item \textbf{\eng{Eventually}} = finalement / à la fin (\eng{He eventually agreed.}). Pour « éventuellement » : \textbf{\eng{possibly}}, \textbf{\eng{potentially}}, \textbf{\eng{if necessary}}.
    \item \textbf{\eng{Actual / Actually}} = réel / en fait (\eng{The actual cost was higher. / Actually, I disagree.}). Pour « actuel / actuellement » : \textbf{\eng{current}}, \textbf{\eng{present}}, \textbf{\eng{at the moment}}, \textbf{\eng{currently}}.
    \item \textbf{\eng{Participate \textbf{in}}} (préposition obligatoire). \eng{*Participate to} est une erreur majeure (calque du français).
    \item \textbf{\eng{Prioritise}} (orthographe britannique \textbf{-ise}) vs US \eng{prioritize}. Au Cameroun (norme UK) : \eng{prioritise}.
    \item \textbf{\eng{Programme}} (UK) vs \eng{Program} (US). Contexte scolaire/administratif : \eng{programme}.
\end{itemize}
\end{attention}

\begin{methode}[Structurer une présentation orale formelle (2 min) — \eng{Pitch}]
\begin{enumerate}
    \item \textbf{Introduction (15 s) :} Saluer, nommer le groupe, annoncer la thèse centrale (\eng{feasibility, inclusivity}). \eng{Good morning everyone. We are Group 3. Our proposal for the Community Day focuses on efficiency and inclusivity.}
    \item \textbf{Corps — Présentation du \eng{schedule} (1 min 15 s) :} Découper par grands blocs horaires (\eng{First, Then, Finally}). Citer les \eng{time slots} clés. Justifier les choix lexicaux (\eng{We allocated... because... ; We prioritised...}). \textbf{Insérer les 2 idiomes} naturellement.
    \item \textbf{Plan de secours \eng{(Contingency Plan)} (15 s) :} Structure conditionnelle Type 1 : \eng{If it rains, we will postpone the clean-up to the 2:00 PM slot. Volunteers are on standby.}
    \item \textbf{Conclusion — Appel à l'action (15 s) :} Réaffirmer la faisabilité. \eng{We are confident this schedule is feasible and realistic. Count us in! Thank you for your attention.}
\end{enumerate}
\end{methode}

\begin{savaistu}[Culture anglophone — \eng{Volunteering \& Community Service}]
\noindent Au Royaume-Uni et aux États-Unis, le \eng{community service} (service civique/bénévolat) est une composante majeure de l'éducation et de la vie citoyenne.
\begin{itemize}
    \item \textbf{UK :} \eng{National Citizen Service (NCS)} pour les 16–17 ans (résidentiel + projet social). \eng{Duke of Edinburgh's Award} (compétences, bénévolat, expédition).
    \item \textbf{USA :} \eng{Service Learning} intégré au cursus (lycée/université). \eng{Volunteer hours} souvent requis pour le diplôme (\eng{graduation}) ou les bourses.
    \item \textbf{Cameroun :} Mouvement \eng{Scout}, \eng{Jeunesse Étudiante Chrétienne (JEC)}, \eng{Clubs UNESCO}, nettoyage du \eng{campus} (\"clean-up\"), journées \eng{Community Work} (travaux d'intérêt général). Le \eng{volunteering} valorise le \eng{CV} et l'esprit d'initiative (\eng{leadership}).
\end{itemize}
\eng{“Volunteers don't get paid, not because they're worthless, but because they're priceless.”} — Sherry Anderson
\end{savaistu}

\courssub{Production modèle et analyse didactique}
\noindent Voici une production modèle (écrite + script oral) pour le \textbf{Groupe 3}, suivie d'une analyse commentée des choix lexicaux, grammaticaux et phonologiques.

\begin{exoresolu}[Production modèle — Groupe 3 : \eng{“Unity in Action” Schedule \& Oral Pitch}]
\textbf{A. Mini-programme écrit (Schedule) — \eng{School Community Day, 20 Nov}}

\begin{center}
\begin{tabularx}{\linewidth}{|l|X|}\hline
\rowcolor{popblueL} \textbf{Time Slot} & \textbf{Activity \& Details} \\ \hline
07:00 – 07:30 & \eng{Volunteers arrive and sign in at the main gate. Project Manager allocates tasks.} \\ \hline
07:30 – 09:30 & \eng{Clean-up Drive: Campus \& Molyko Street. \textbf{Prioritise} waste sorting. \textbf{Lend a hand} where needed.} \\ \hline
09:30 – 10:00 & \eng{Break / Hydration. \textbf{Contingency}: If heavy rain, \textbf{postpone} outdoor clean-up to 14:00.} \\ \hline
10:00 – 12:30 & \eng{Blood Donation Drive (School Hall). Red Cross partners \textbf{on time}. 50 donors targeted.} \\ \hline
12:30 – 13:30 & \eng{Lunch Break. Cultural team \textbf{set up} stage. No \textbf{clash} with donation.} \\ \hline
13:30 – 15:30 & \eng{Fund-raising Cultural Show. \textbf{Count me in} for the dance troupe! Entry fee: 500 FCFA.} \\ \hline
15:30 – 16:00 & \eng{Closing Ceremony \& Prize Giving. Principal's speech. \textbf{Beat the clock} to finish \textbf{on time}.} \\ \hline
16:00 – 17:00 & \eng{Final Clean-up \& Debrief. \textbf{Deadline} for expense reports: 22 Nov.} \\ \hline
\end{tabularx}
\end{center}

\textbf{B. Script de la présentation orale (Spokesperson)}

\eng{Good morning, Principal Ekosso, teachers, fellow students.}
\eng{We are Group 3: “Unity in Action”. Our schedule is \textbf{feasible} and \textbf{realistic} because we \textbf{allocated} buffer zones between major events.}
\eng{\textbf{First}, the clean-up starts early at 7:30 AM while energy is high. We \textbf{prioritise} safety with gloves and bags. \textbf{Count me in} — I’ll be leading the Molyko street team!}
\eng{\textbf{Then}, the blood drive runs from 10:00 to 12:30 in the hall. No \textbf{clash} with the stage set-up outside. The Red Cross confirmed they will be \textbf{on the dot}.}
\eng{\textbf{If rain forces us to \textbf{postpone} the clean-up}, our contingency plan \textbf{will move} it to the 2:00 PM slot, pushing the cultural show slightly later. We have \textbf{volunteers} on standby.}
\eng{\textbf{Finally}, the fund-raising show. Entry fees go straight to the library fund. We aim to \textbf{beat the clock} and wrap up by 4:00 PM sharp.}
\eng{To join, students simply \textbf{sign up} at the Student Council office by Friday’s \textbf{deadline}. \textbf{Lend a hand}, make a difference. Thank you.}

\tcblower
\textbf{Commentaire didactique des choix de langue}

\begin{enumerate}
    \item \textbf{Cible lexicale atteinte (11/8 minimum) :} Le texte écrit et oral intègre \textbf{11 mots-cibles} distincts : \eng{volunteers, sign up/sign in, allocate, prioritise, lend a hand, contingency, postpone, on time/on the dot, clash, beat the clock, deadline, feasible, realistic}. L'exigence quantitative est largement satisfaite.
    \item \textbf{Expressions idiomatiques (3/2) :} \eng{“Count me in”} (adhésion enthousiaste), \eng{“beat the clock”} (gagner du temps/terminer à temps) et \eng{“lend a hand”} (appel final + instruction écrite) sont utilisées naturellement. Cohésion lexicale renforcée.
    \item \textbf{Gestion du temps (prépositions/particules) :} \eng{at 7:30 AM} (heure précise), \eng{from 10:00 to 12:30} (durée), \eng{by Friday’s deadline} (échéance), \eng{on standby} (disponible/réserve). Aucune erreur de préposition temporelle.
    \item \textbf{Faux amis évités :} Utilisation correcte de \eng{schedule} (pas \eng{planning}), \eng{participate in} / \eng{take part in} (pas \eng{*participate to}), \eng{attend} implicite via \eng{arrive/sign in}. \eng{Actually} non utilisé à tort. Orthographe UK : \eng{prioritise, programme}.
    \item \textbf{Structures grammaticales ancrées dans la fonction :}
    \begin{itemize}
        \item \textbf{Conditionnel réel (Type 1)} pour le plan de secours : \eng{If rain forces us to postpone..., our plan will move...}. Correct (présent + will).
        \item \textbf{Voix passive / tournures impersonnelles} pour l'objectivité administrative : \eng{Entry fees go straight to...}, \eng{Red Cross confirmed they would be...} (discours rapporté), \eng{Volunteers are needed}.
        \item \textbf{Impératif / Infinitif de but} dans le programme écrit : \eng{Prioritise waste sorting}, \eng{To join, students simply sign up...}.
        \item \textbf{Future continuous} pour action planifiée : \eng{I’ll be leading...}.
    \end{itemize}
    \item \textbf{Registre, cohésion et prononciation :} Ton formel mais engageant (\eng{fellow students, We are confident...}). Connecteurs logiques (\eng{First, Then, Finally, If}). Vocabulaire de liaison expert (\eng{buffer zones, contingency plan, on standby, wrap up}). Mots-clés à prononcer soigneusement : \eng{deadline} (« déd-laïne »), \eng{volunteer} (« vo-lon-ti-eu »), \eng{feasible} (« fi-zi-beul »), \eng{prioritise} (« pra-io-ri-taïze »), \eng{allocate} (« a-lo-keute »).
\end{enumerate}
\end{exoresolu}

\courssub{Bilan de l'activité \& Auto-évaluation}
\noindent Cette séquence de type \cle{Task-Based Learning} a permis de valider les compétences visées par la leçon 23 dans une situation intégrée authentique (contexte camerounais, enjeu réel).

\begin{itemize}
    \item \textbf{Compréhension (Reading/Listening) :} Le memo du Proviseur a servi de déclencheur authentique ; extraction d'informations précises (\eng{deadline, partners, threefold objective, sharp}).
    \item \textbf{Lexique (Vocabulary) :} Réemploi contraint (8 mots-cibles + 2 idiomes) forçant l'activation de la mémoire à long terme et la vérification des collocations (\eng{sign up for, participate in, on time vs in time, clash with, allocate to}).
    \item \textbf{Grammaire en contexte :} Usage naturel du \eng{First Conditional} pour la contingence (\eng{If rain forces us to postpone..., we will move...}) et de l'impératif/infinitif pour les consignes du \eng{schedule}.
    \item \textbf{Expression écrite (Writing) :} Rédaction d'un \eng{schedule} structuré (tableau horaire) — compétence « produire un document professionnel/administratif » (programme, rapport), attendue au Baccalauréat.
    \item \textbf{Expression orale (Speaking) :} Présentation chronométrée (2 min) avec rôles répartis : fluidité, prononciation des mots-clés, interaction (phase de vote/justification).
    \item \textbf{Compétence transversale :} Ancrage camerounais (Buea, Molyko, FCFA, Red Cross, Principal Ekosso) donnant du sens à l'apprentissage et valorisant le contexte socioculturel.
\end{itemize}

\begin{aretenir}[Auto-évaluation de l'élève — Grille de validation]
L'élève coche les items réussis :
\begin{itemize}
    \item[ ] J'ai utilisé \textbf{8+ mots-cibles} correctement dans mon \eng{schedule} écrit (orthographe UK).
    \item[ ] J'ai intégré \textbf{2 idiomes} naturellement à l'oral (\eng{count me in, beat the clock, lend a hand, on the dot, in the nick of time}).
    \item[ ] Je n'ai fait \textbf{aucun faux ami} (\eng{participate in, schedule, actually, eventually, assist/attend}).
    \item[ ] Mon \eng{contingency plan} utilise le \textbf{Type 1} : \eng{If + present, will + V} (\eng{If it rains, we will postpone...}).
    \item[ ] J'ai justifié mon vote en anglais avec du vocabulaire de la leçon (\eng{feasible, clash, allocated, on time, buffer zone, prioritise}).
    \item[ ] Je prononce correctement : \eng{deadline /'dedlaɪn/ (« déd-laïne »), volunteer /vɒlən'tɪə/ (« vo-lon-ti-eu »), feasible /'fi:zəbl/ (« fi-zi-beul »), prioritise /praɪ'ɒrɪtaɪz/ (« pra-io-ri-taïze »)}.
\end{itemize}
\textit{Objectif atteint ? $\rightarrow$ Poursuivre avec la fiche de révision Chapter 2. Difficultés persistantes ? $\rightarrow$ Revoir les boîtes \textbf{Attention} (faux amis), \textbf{Propriété} (collocations) et \textbf{Méthode} (structure orale) de la leçon 23.}
\end{aretenir}

\newpage

\coursec{Méthodes et exercices résolus}

\courssub{Méthodes et exercices résolus}

\begin{methode}[Identifier et employer le vocabulaire de la participation (community, workplace, school)]
\textbf{Objectif :} Maîtriser les verbes, noms et adjectifs clés pour décrire un engagement actif dans une collectivité, une entreprise ou un établissement scolaire.
\begin{enumerate}
    \item \textbf{Repérer le champ lexical :} Dans un texte ou un script d'écoute, surlignez les mots qui expriment l'action de \cle{joining}, \cle{contributing}, \cle{leading} ou \cle{supporting} une activité collective.
    \item \textbf{Classifier par fonction :}
    \begin{itemize}
        \item \textit{Verbes d'action :} \eng{to volunteer, to participate, to contribute, to engage, to collaborate, to spearhead, to facilitate, to attend}.
        \item \textit{Noms d'acteurs/rôles :} \eng{a volunteer, a participant, a contributor, a stakeholder, a chairperson, a facilitator, a member, an attendee}.
        \item \textit{Noms d'actions/structures :} \eng{involvement, participation, contribution, engagement, committee, task force, board, assembly, outreach, initiative}.
        \item \textit{Adjectifs/Qualificatifs :} \eng{active, proactive, committed, dedicated, collaborative, inclusive, voluntary, mandatory}.
    \end{itemize}
    \item \textbf{Vérifier la construction syntaxique :} Noter les prépositions associées (\eng{participate \textbf{in}}, \eng{contribute \textbf{to}}, \eng{engage \textbf{in}}, \eng{volunteer \textbf{for/as}}, \eng{serve \textbf{on} a committee}).
    \item \textbf{Produire :} Réutiliser 5 à 7 mots de cette liste pour décrire une expérience personnelle (ex. : classe, association de quartier, stage).
\end{enumerate}
\end{methode}

\begin{exoresolu}[Compréhension et emploi lexical : Compte-rendu d'assemblée communautaire]
\textbf{Énoncé :} Lisez le texte ci-dessous et répondez aux questions.

\begin{textebox}[Minutes of the Buea Youth Association Monthly Meeting]
The monthly meeting of the Buea Youth Association (BYA) was held last Saturday at the Community Hall. The \textbf{chairperson}, Mr. Atanga, opened the session at 10:00 a.m. He praised the high \textbf{turnout} and thanked every \textbf{attendee} for their \textbf{commitment}. The main \textbf{agenda} item was the upcoming \textbf{clean-up campaign} scheduled for Labour Day. 

Mrs. Nkeng, the \textbf{treasurer}, presented the financial report. She highlighted the need for more \textbf{fundraising} initiatives to \textbf{sponsor} the event. Several members \textbf{volunteered} to \textbf{spearhead} specific sub-committees: logistics, communication, and refreshments. Mr. Fon, a new \textbf{stakeholder} from the local council, \textbf{pledged} financial support on behalf of the municipality. 

The \textbf{facilitator} reminded everyone that active \textbf{participation} is mandatory for all registered members. The meeting was \textbf{adjourned} at 12:30 p.m. after the \textbf{chairperson} \textbf{summarised} the \textbf{action points} and set the date for the next \textbf{follow-up meeting}.
\end{textebox}

\textbf{Questions :}
\begin{enumerate}
    \item Trouvez dans le texte les équivalents anglais des termes suivants : \textit{président de séance, affluence, participant, engagement, ordre du jour, campagne de nettoyage, trésorier, collecter des fonds, parrainer, se porter volontaire, diriger (un comité), partie prenante, promettre, animateur, points d'action, réunion de suivi}.
    \item Complétez les phrases suivantes avec la préposition correcte (\textit{in, to, for, on, as, by}) :
    \begin{itemize}
        \item All members must participate \_\_\_\_ the clean-up campaign.
        \item Mrs. Nkeng contributed \_\_\_\_ the financial report.
        \item Three members volunteered \_\_\_\_ the logistics committee.
        \item Mr. Fon serves \_\_\_\_ a stakeholder \_\_\_\_ the local council.
        \item The meeting was chaired \_\_\_\_ Mr. Atanga.
    \end{itemize}
    \item Réécrivez les phrases suivantes en remplaçant les parties soulignées par un mot du texte (forme adaptée) :
    \begin{itemize}
        \item The president \textbf{closed} the meeting at 12:30.
        \item We need more \textbf{money-raising} ideas.
        \item He \textbf{promised} to give money.
    \end{itemize}
\end{enumerate}

\tcblower

\textbf{Solution détaillée}

\textbf{1. Correspondances lexicales (Vocabulary Matching) :}
\begin{center}
\begin{tabularx}{\linewidth}{|X|X|}\hline
\rowcolor{popblueL} \textbf{Français} & \textbf{English (from text)} \\ \hline
Président de séance & \eng{chairperson} \\ \hline
Affluence & \eng{turnout} \\ \hline
Participant & \eng{attendee} \\ \hline
Engagement & \eng{commitment} \\ \hline
Ordre du jour & \eng{agenda} \\ \hline
Campagne de nettoyage & \eng{clean-up campaign} \\ \hline
Trésorier & \eng{treasurer} \\ \hline
Collecter des fonds (nom) & \eng{fundraising} \\ \hline
Parrainer (financièrement) & \eng{sponsor} \\ \hline
Se porter volontaire & \eng{volunteered} \\ \hline
Diriger / Mener (une initiative) & \eng{spearhead} \\ \hline
Partie prenante & \eng{stakeholder} \\ \hline
Promettre (solennellement) & \eng{pledged} \\ \hline
Animateur / Facilitateur & \eng{facilitator} \\ \hline
Points d'action & \eng{action points} \\ \hline
Réunion de suivi & \eng{follow-up meeting} \\ \hline
\end{tabularx}
\end{center}
\emph{Note :} \eng{Chairperson} est le terme neutre et moderne pour \eng{chairman/chairwoman}. \eng{Turnout} désigne le nombre de personnes présentes (affluence). \eng{Spearhead} implique de diriger l'avant-garde d'un projet.

\textbf{2. Prépositions (Grammar in Context) :}
\begin{itemize}
    \item participate \textbf{in} (participer \textbf{à} une activité) $\rightarrow$ \eng{participate in the clean-up campaign}.
    \item contribute \textbf{to} (contribuer \textbf{à} un rapport, une discussion) $\rightarrow$ \eng{contributed to the financial report}.
    \item volunteer \textbf{for} (se porter volontaire \textbf{pour} une tâche/un comité) $\rightarrow$ \eng{volunteered for the logistics committee}. (On peut aussi dire \eng{volunteered to join}).
    \item serve \textbf{as} a stakeholder \textbf{from} (or \textbf{representing}) the local council. \eng{Serve as} = agir en tant que. \eng{From} indique l'origine/appartenance.
    \item chaired \textbf{by} (voix passive : présidé \textbf{par}) $\rightarrow$ \eng{was chaired by Mr. Atanga}.
\end{itemize}

\textbf{3. Reformulation lexicale (Paraphrasing) :}
\begin{itemize}
    \item The president \textbf{adjourned} the meeting at 12:30. (\eng{To adjourn} = lever la séance, terme formel pour les réunions).
    \item We need more \textbf{fundraising} ideas. (\eng{Fundraising} est le nom composé standard ; \eng{money-raising} n'existe pas).
    \item He \textbf{pledged} to give money. (\eng{To pledge} = s'engager solennellement, promettre formellement, souvent pour un don).
\end{itemize}

\textbf{Conclusion :} Ce texte modèle concentre le vocabulaire institutionnel et associatif essentiel. La maîtrise des collocations verbales (\eng{participate in, volunteer for, spearhead, pledge to}) est cruciale pour l'expression écrite (rapports, comptes-rendus) et orale (simulation de réunion).
\end{exoresolu}

\begin{methode}[Maîtriser le vocabulaire de la gestion du temps (planning, deadlines, priorités)]
\textbf{Objectif :} Utiliser avec précision le lexique de l'organisation temporelle dans un contexte professionnel ou scolaire (emploi du temps, échéances, priorisation).
\begin{enumerate}
    \item \textbf{Distinguer les notions temporelles :}
    \begin{itemize}
        \item \textit{Moment précis :} \eng{deadline, due date, cutoff, timestamp, appointment, slot}.
        \item \textit{Durée/Processus :} \eng{timeline, schedule, timetable, agenda, planner, timesheet, duration, span, lead time}.
        \item \textit{Action sur le temps :} \eng{to schedule, to timetable, to plan, to prioritize, to allocate (time), to budget (time), to meet a deadline, to beat a deadline, to miss a deadline, to extend a deadline, to postpone, to reschedule, to fast-track}.
    \end{itemize}
    \item \textbf{Maîtriser les adjectifs de fréquence et d'urgence :} \eng{urgent, pressing, critical, imminent, overdue, pending, timely, punctual, prompt, chronic (late), last-minute}.
    \item \textbf{Structures utiles :}
    \begin{itemize}
        \item \eng{The deadline \textbf{is} Friday / \textbf{falls on} Friday.}
        \item \eng{We \textbf{have} a tight schedule / \textbf{are running} behind schedule / \textbf{are} ahead of schedule.}
        \item \eng{Please \textbf{submit} your report \textbf{by} Monday / \textbf{before} 5 p.m. / \textbf{no later than} Friday.}
        \item \eng{\textbf{Prioritise} tasks \textbf{according to} urgency / \textbf{using} the Eisenhower Matrix.}
    \end{itemize}
    \item \textbf{Entraînement :} Convertir un emploi du temps français en anglais (ex. : \textit{« Je dois rendre le devoir pour mardi »} $\rightarrow$ \eng{The assignment is due Tuesday / I have to submit the assignment by Tuesday.}).
\end{enumerate}
\end{methode}

\begin{exoresolu}[Grammaire et lexique : Gestion de projet au lycée]
\textbf{Énoncé :} Le texte suivant contient 10 erreurs lexicales ou grammaticales (mots mal choisis, prépositions incorrectes, faux amis). Identifiez-les, corrigez-les et justifiez chaque correction.

\begin{textebox}[Student's Project Log - Week 3]
\textbf{Monday:} The \textbf{deadline} for the science project is \textbf{in} Friday. I am \textbf{very late} because I started \textbf{since} last week only. I need to \textbf{program} my week carefully. I have a \textbf{charge} of work to do. I will \textbf{make} a planning for each day.
\textbf{Tuesday:} I \textbf{passed} three hours \textbf{on} research. My \textbf{time} is running out. The teacher said the \textbf{date limit} is fixed. I must \textbf{respect} the deadline.
\textbf{Wednesday:} I \textbf{advance} well. I have \textbf{finished} the introduction. I will \textbf{anticipate} the conclusion tomorrow. I have a \textbf{meeting} with my supervisor \textbf{at} 14h. We will \textbf{discuss about} the methodology.
\textbf{Thursday:} Crisis! I \textbf{lost} my USB key. It is a \textbf{catastrophe}. I have to \textbf{recommence} everything. I am \textbf{in delay}. I sent an email to \textbf{demand} an extension of the \textbf{delai}.
\textbf{Friday:} The teacher \textbf{accepted} my request. The new \textbf{deadline} is \textbf{on} next Monday. I am \textbf{safe} now. I will \textbf{profit} the weekend to work hard.
\end{textebox}

\tcblower

\textbf{Solution détaillée : Analyse et Corrections}

Voici les 10 erreurs identifiées, classées par type, avec la correction et l'explication grammaticale/lexicale.

\begin{center}
\begin{tabularx}{\linewidth}{|c|X|X|X|}\hline
\rowcolor{popblueL} \textbf{N°} & \textbf{Erreur (Texte original)} & \textbf{Correction} & \textbf{Justification / Règle} \\ \hline
1 & \eng{deadline ... \textbf{in} Friday} & \eng{\textbf{on} Friday} & \textbf{Préposition de temps :} \eng{on} pour les jours de la semaine/dates précises. \eng{In} = mois, années, saisons, parties de la journée (\eng{in the morning}). \\ \hline
2 & \eng{I am \textbf{very late}} & \eng{I am \textbf{running late} / \textbf{behind schedule}} & \textbf{Faux ami / Calque :} \eng{Late} (adjectif) = en retard (arrivée). Pour un projet/avancement : \eng{behind schedule} ou \eng{running behind}. \eng{Very late} suggère l'heure d'arrivée, pas l'état d'avancement. \\ \hline
3 & \eng{started \textbf{since} last week} & \eng{started \textbf{only} last week / \textbf{have only started} last week} & \textbf{Grammaire :} \eng{Since} + Present Perfect (\eng{have started}). Avec \eng{started} (Past Simple), on utilise \eng{last week} seul ou \eng{only last week}. \\ \hline
4 & \eng{\textbf{program} my week} & \eng{\textbf{plan} / \textbf{schedule} / \textbf{organize} my week} & \textbf{Faux ami :} \eng{To program} = programmer (informatique, machine, événement culturel). Pour son temps personnel : \eng{to plan} ou \eng{to schedule}. \\ \hline
5 & \eng{a \textbf{charge} of work} & \eng{a \textbf{lot of} work / a \textbf{heavy workload} / \textbf{plenty of} work} & \textbf{Faux ami :} \eng{Charge} = charge (électrique, responsabilité, accusation). \eng{Workload} = charge de travail. \eng{A charge of work} n'existe pas. \\ \hline
6 & \eng{\textbf{make} a planning} & \eng{\textbf{draw up} / \textbf{make} a \textbf{schedule} / \textbf{timetable} / \textbf{plan}} & \textbf{Faux ami / Anglicisme :} \eng{Planning} (nom) existe en anglais mais signifie « planification » (processus) ou « planning » (document) en français. En anglais standard, on préfère \eng{schedule} (US), \eng{timetable} (UK école/transport), \eng{plan}. Verbe : \eng{draw up a plan}. \\ \hline
7 & \eng{\textbf{passed} three hours} & \eng{\textbf{spent} three hours} & \textbf{Faux ami majeur :} \eng{To pass time} = passer le temps (s'occuper). \eng{To spend time} = consacrer du temps à une activité. \\ \hline
8 & \eng{\textbf{date limit}} & \eng{\textbf{deadline} / \textbf{cut-off date}} & \textbf{Calque du français :} \eng{Date limit} n'est pas l'expression idiomatique. \eng{Deadline} est le mot unique standard. \\ \hline
9 & \eng{\textbf{discuss about}} & \eng{\textbf{discuss}} & \textbf{Grammaire verbes transitifs :} \eng{Discuss} est transitif direct (discuss something). Pas de préposition. (Mais : \eng{talk about}, \eng{speak about}). \\ \hline
10 & \eng{\textbf{demand} an extension ... \textbf{delai}} & \eng{\textbf{request} / \textbf{ask for} an extension ... \textbf{deadline}} & \textbf{Registre / Faux ami :} \eng{Demand} = exiger (très fort, impoli pour un prof). \eng{Request} / \eng{ask for} = demander poliment. \eng{Delai} n'existe pas en anglais ($\rightarrow$ \eng{deadline, extension, postponement}). \\ \hline
\end{tabularx}
\end{center}

\textbf{Autres corrections notables (bonus) :}
\begin{itemize}
    \item \eng{at 14h} $\rightarrow$ \eng{at 2 p.m.} (format 12h standard en anglais, ou \eng{14:00} style militaire/horaire ferroviaire sans 'h').
    \item \eng{I \textbf{advance} well} $\rightarrow$ \eng{I am \textbf{making good progress} / \textbf{getting on} well}. \eng{Advance} (verbe intransitif) = avancer physiquement (armée) ou progresser dans une carrière ; pour un travail : \eng{make progress}.
    \item \eng{I will \textbf{anticipate} the conclusion} $\rightarrow$ \eng{I will \textbf{write/draft} the conclusion tomorrow}. \eng{Anticipate} = prévoir, s'attendre à, devancer. Ne signifie pas « faire à l'avance ».
    \item \eng{I \textbf{lost} my USB key} $\rightarrow$ \eng{I \textbf{misplaced} / \textbf{lost} my \textbf{flash drive / USB stick}}. \eng{USB key} est du Franglais ; \eng{flash drive} ou \eng{USB stick} est standard.
    \item \eng{I am \textbf{in delay}} $\rightarrow$ \eng{I am \textbf{behind schedule} / \textbf{running late}}. \eng{In delay} n'existe pas.
    \item \eng{I will \textbf{profit} the weekend} $\rightarrow$ \eng{I will \textbf{make the most of} / \textbf{take advantage of} the weekend}. \eng{Profit} est un nom (bénéfice) ou verbe intransitif (\eng{profit from}). Verbe transitif : \eng{benefit from} ou \eng{make the most of}.
\end{itemize}

\textbf{Version corrigée (Extrait) :}
\eng{``Monday: The \textbf{deadline} for the science project is \textbf{on} Friday. I am \textbf{behind schedule} because I \textbf{only started} last week. I need to \textbf{plan} my week carefully. I have a \textbf{heavy workload}. I will \textbf{draw up a schedule} for each day.''}
\end{exoresolu}

\begin{methode}[Maîtriser les collocations et expressions idiomatiques du thème]
\textbf{Objectif :} Reconnaître et produire les associations de mots figées (collocations) et les idiomes fréquents pour « prendre part » et « gérer son temps » afin de gagner en naturel et en précision (niveau B2/C1).
\begin{enumerate}
    \item \textbf{Identifier le type de collocation :}
    \begin{itemize}
        \item \textit{Verbe + Nom :} \eng{hold a meeting, chair a meeting, attend a meeting, call a meeting, adjourn a meeting, set a deadline, meet a deadline, miss a deadline, extend a deadline, allocate resources, prioritise tasks, volunteer for a task}.
        \item \textit{Adjectif + Nom :} \eng{active participation, key stakeholder, tight schedule, pressing deadline, prior commitment, voluntary work, mandatory attendance, unanimous decision}.
        \item \textit{Nom + Nom (composés) :} \eng{time management, team building, decision-making, problem-solving, timekeeping, schedule conflict, action item, follow-up}.
    \end{itemize}
    \item \textbf{Mémoriser les idiomes « temps/participation » :}
    \begin{center}
    \begin{tabularx}{\linewidth}{|X|X|X|}\hline
    \rowcolor{popblueL} \textbf{Idiome} & \textbf{Sens (FR)} & \textbf{Exemple contextuel} \\ \hline
    \eng{To \textbf{step up}} & S'investir davantage, prendre ses responsabilités & \eng{We need someone to \textbf{step up} and lead the team.} \\ \hline
    \eng{To \textbf{pitch in}} & Mettre la main à la pâte, participer activement & \eng{Everyone \textbf{pitched in} to clean the classroom.} \\ \hline
    \eng{To \textbf{pull one's weight}} & Faire sa part du travail & \eng{He always \textbf{pulls his weight} in group projects.} \\ \hline
    \eng{\textbf{In the nick of time}} & Juste à temps, in extremis & \eng{We submitted the report \textbf{in the nick of time}.} \\ \hline
    \eng{\textbf{Against the clock}} & Contre la montre, en urgence & \eng{We are working \textbf{against the clock} to finish.} \\ \hline
    \eng{\textbf{Buy time}} & Gagner du temps & \eng{Asking for an extension \textbf{buys us some time}.} \\ \hline
    \eng{\textbf{On the back burner}} & Mis de côté, en attente & \eng{The project is \textbf{on the back burner} for now.} \\ \hline
    \eng{\textbf{Time is of the essence}} & Le temps presse (formel/juridique) & \eng{Time is of the essence; we must decide now.} \\ \hline
    \end{tabularx}
    \end{center}
    \item \textbf{Exercice de substitution :} Remplacer un verbe générique (\eng{do, make, give, put}) par la collocation forte dans une phrase.
    \item \textbf{Attention au registre :} \eng{Pitch in} (informel) vs \eng{Contribute actively} (formel) ; \eng{Step up} (positif) vs \eng{Take over} (peut être négatif).
\end{enumerate}
\end{methode}

\begin{exoresolu}[Collocations et idiomes : Choix multiples et reformulation]
\textbf{Énoncé :}
\begin{enumerate}
    \item \textbf{Choisissez la collocation correcte pour compléter chaque phrase.}
    \begin{enumerate}
        \item The chairperson decided to \_\_\_\_ the meeting until next week because key members were absent.
              \begin{itemize}
                  \item[a)] postpone \item[b)] cancel \item[c)] adjourn \item[d)] delay
              \end{itemize}
        \item As a \_\_\_\_ member of the committee, you are expected to attend every session.
              \begin{itemize}
                  \item[a)] active \item[b)] acting \item[c)] actual \item[d)] activated
              \end{itemize}
        \item We are running \_\_\_\_ schedule; we must work faster to \_\_\_\_ the deadline.
              \begin{itemize}
                  \item[a)] behind / meet \item[b)] late / catch \item[c)] out of / reach \item[d)] off / do
              \end{itemize}
        \item The manager asked us to \_\_\_\_ tasks based on urgency and importance.
              \begin{itemize}
                  \item[a)] prioritize \item[b)] prioritise \item[c)] priorize \item[d)] preorize
              \end{itemize}
        \item Don't put the project on the \_\_\_\_; we need results now.
              \begin{itemize}
                  \item[a)] back burner \item[b)] back seat \item[c)] shelf \item[d)] side
              \end{itemize}
    \end{enumerate}
    \item \textbf{Réécrivez les phrases suivantes en utilisant l'idiome ou la collocation appropriée (forme conjuguée si nécessaire).}
    \begin{enumerate}
        \item Everyone helped to organize the school fair. (\eng{pitch in})
        \item We finished the exam just before the time limit. (\eng{in the nick of time})
        \item She contributes her fair share to the team effort. (\eng{pull one's weight})
        \item We need to delay the decision to get more information. (\eng{buy time})
    \end{enumerate}
\end{enumerate}

\tcblower

\textbf{Solution détaillée}

\textbf{1. Choix multiples (QCM) :}

\begin{enumerate}
    \item \textbf{Réponse : c) adjourn.}
    \textbf{Justification :} \eng{Adjourn} est le terme technique parlementaire/institutionnel pour « suspendre une séance pour la reprendre plus tard ». \eng{Postpone} s'applique à un événement unique (match, examen) pas à une séance récurrente qu'on « lève ». \eng{Cancel} = annuler définitivement. \eng{Delay} (verbe) = retarder le début, pas suspendre en cours.
    \item \textbf{Réponse : a) active.}
    \textbf{Justification :} Collocation standard : \eng{active member / active participation / active role}. \eng{Acting} = intérimaire (\eng{acting chairperson}). \eng{Actual} = réel (contraire de prévu). \eng{Activated} = mis en marche (machine, compte).
    \item \textbf{Réponse : a) behind / meet.}
    \textbf{Justification :} Expression figée : \eng{running behind schedule} (ou \eng{ahead of schedule}). Verbe collocatif avec \eng{deadline} : \eng{meet a deadline} (respecter), \eng{miss a deadline} (rater), \eng{beat a deadline} (deviner). \eng{Catch a deadline} n'existe pas.
    \item \textbf{Réponse : a) prioritize / b) prioritise.}
    \textbf{Justification :} Les deux orthographes sont correctes : \eng{prioritize} (US / Oxford spelling) et \eng{prioritise} (UK standard). Au Cameroun (norme britannique), \textbf{prioritise} est l'orthographe de référence. Les options c) et d) sont des fautes d'orthographe.
    \item \textbf{Réponse : a) back burner.}
    \textbf{Justification :} Idiome fixe : \eng{put on the back burner} = mettre en veille, déprioriser. \eng{Back seat} = \eng{take a back seat} (rester en retrait). \eng{On the shelf} = mis au placard (projet abandonné). \eng{On the side} = accessoire.
\end{enumerate}

\textbf{2. Reformulation guidée (Production) :}

\begin{enumerate}
    \item \eng{Everyone \textbf{pitched in} to organize the school fair.}
    \emph{Note :} \eng{Pitch in} implique un effort collectif, volontaire et énergique. Prétérit : \eng{pitched in}.
    \item \eng{We finished the exam \textbf{in the nick of time}.}
    \emph{Note :} Locution adverbiale figée, invariable. Place en fin de phrase.
    \item \eng{She always \textbf{pulls her weight} in the team. / She \textbf{pulled her weight} during the project.}
    \emph{Note :} Possessif obligatoire (\eng{her/his/their weight}). Sens : assumer sa part de responsabilité.
    \item \eng{We need to \textbf{buy (some) time} to get more information.}
    \emph{Note :} \eng{Buy time} = gagner du temps délibérément par une action (demander report, faire diversion).
\end{enumerate}

\textbf{Conclusion :} La connaissance de ces blocs lexicaux (\eng{meet a deadline, active member, running behind schedule, pitch in, back burner}) distingue un niveau B1 (vocabulaire générique : \eng{do, make, late, help}) d'un niveau B2/Bac (vocabulaire précis, idiomatique, collocationnel).
\end{exoresolu}

\begin{methode}[Éviter les faux amis et les calques du français (Spécial Participation/Temps)]
\textbf{Objectif :} Identifier les pièges lexicaux majeurs pour les francophones et produire l'anglais correct.
\begin{enumerate}
    \item \textbf{Lister les paires « Faux Ami / Vrai Sens » critiques pour ce thème :}
    \begin{center}
    \begin{tabularx}{\linewidth}{|X|X|X|}\hline
    \rowcolor{poporangeL} \textbf{Mot FR (Piège)} & \textbf{Mot EN incorrect (Faux ami)} & \textbf{Mot EN correct} \\ \hline
    \textit{Assister (à une réunion)} & \eng{Assist} & \eng{Attend} (\eng{Assist} = aider) \\ \hline
    \textit{Participer} & \eng{Participate} (sans prép.) & \eng{Participate \textbf{in}} \\ \hline
    \textit{Programme (emploi du temps)} & \eng{Program} & \eng{Schedule / Timetable / Agenda} \\ \hline
    \textit{Planning (document)} & \eng{Planning} & \eng{Schedule / Timetable / Plan / Diary} \\ \hline
    \textit{Charge (de travail)} & \eng{Charge} & \eng{Workload / Burden / Load} \\ \hline
    \textit{Délai (date limite)} & \eng{Delay / Delai} & \eng{Deadline / Time limit / Cut-off date} \\ \hline
    \textit{Retard (avancement)} & \eng{Delay / Late} & \eng{Behind schedule / Overdue / Running late} \\ \hline
    \textit{Demander (poliment)} & \eng{Demand} & \eng{Request / Ask for / Apply for} \\ \hline
    \textit{Profiter (de l'occasion)} & \eng{Profit} & \eng{Take advantage of / Make the most of / Benefit from} \\ \hline
    \textit{Anticiper (faire avant)} & \eng{Anticipate} & \eng{Do in advance / Prepare beforehand / Plan ahead} \\ \hline
    \textit{Animer (une réunion)} & \eng{Animate} & \eng{Chair / Facilitate / Lead / Moderate} \\ \hline
    \textit{Ordre du jour} & \eng{Order of the day} & \eng{Agenda} \\ \hline
    \textit{Compte-rendu} & \eng{Report / Account} & \eng{Minutes (pluriel) / Report} \\ \hline
    \end{tabularx}
    \end{center}
    \item \textbf{Règle de vérification (Reflexe « Bac ») :} Avant d'écrire un mot qui ressemble au français, demandez-vous : « Est-ce que ce verbe est transitif direct ? A-t-il la même préposition ? Est-ce que le sens correspond exactement ? » $\rightarrow$ Si doute $\rightarrow$ Synonyme simple (\eng{help} pour \eng{assist}, \eng{go to} pour \eng{attend}, \eng{plan} pour \eng{program}).
    \item \textbf{Traquer les calques syntaxiques :}
    \begin{itemize}
        \item \eng{Discuss about} $\times$ $\rightarrow$ \eng{Discuss} $\checkmark$.
        \item \eng{Attend to a meeting} $\times$ ($\eng{attend to}$ = s'occuper de) $\rightarrow$ \eng{Attend a meeting} $\checkmark$.
        \item \eng{Missing of time} $\times$ $\rightarrow$ \eng{Lack of time / Shortage of time} $\checkmark$.
        \item \eng{Take a decision} $\times$ (Anglicisme FR) $\rightarrow$ \eng{Make a decision} $\checkmark$.
    \end{itemize}
\end{enumerate}
\end{methode}

\begin{exoresolu}[Correction d'erreurs types « Faux Amis / Calques »]
\textbf{Énoncé :} Traduisez les phrases suivantes en anglais en évitant les pièges du français. Pour chaque phrase, soulignez le mot-clé qui pose problème et notez la règle appliquée.

\begin{enumerate}
    \item Le directeur a \textbf{assisté} à la réunion des parents d'élèves.
    \item Nous devons \textbf{programmer} la date de l'examen oral.
    \item J'ai une \textbf{charge} de travail énorme cette semaine.
    \item Le \textbf{délai} pour le dépôt des dossiers est vendredi.
    \item Il a \textbf{demandé} une augmentation de salaire.
    \item Nous avons \textbf{discuté à propos} du budget.
    \item Elle a \textbf{animé} l'atelier avec brio.
    \item Je vais \textbf{profiter} de mes vacances pour réviser.
    \item Le projet est \textbf{en retard}.
    \item Prenons une \textbf{décision} maintenant.
\end{enumerate}

\tcblower

\textbf{Solution détaillée : Traductions commentées}

\begin{enumerate}
    \item \eng{The principal \textbf{attended} the PTA meeting.}
    \textbf{Piège :} \textbf{Assister} $\neq$ \eng{Assist}. \eng{Assist} = aider (\eng{The nurse assisted the doctor}). \eng{Attend} = être présent à un événement. \textbf{Règle :} Faux ami sémantique.
    \item \eng{We must \textbf{schedule} / \textbf{set a date for} the oral exam.}
    \textbf{Piège :} \textbf{Programmer} $\neq$ \eng{Program} (sauf informatique/TV). \eng{Schedule} (verbe) = fixer une date/heure. \textbf{Règle :} Faux ami / Collocation.
    \item \eng{I have a huge \textbf{workload} this week. / My \textbf{workload} is heavy.}
    \textbf{Piège :} \textbf{Charge} $\neq$ \eng{Charge}. \eng{Workload} = charge de travail. \eng{Charge} = accusation, frais, responsabilité, électricité. \textbf{Règle :} Faux ami / Mot composé.
    \item \eng{The \textbf{deadline} for submitting applications is Friday. / The \textbf{cut-off date} is Friday.}
    \textbf{Piège :} \textbf{Délai} $\neq$ \eng{Delay} (retard) ni \eng{Delai} (inexistant). \eng{Deadline} = date butoir. \textbf{Règle :} Faux ami majeur / Calque.
    \item \eng{He \textbf{requested} / \textbf{asked for} a pay rise.}
    \textbf{Piège :} \textbf{Demander} $\neq$ \eng{Demand} (exiger, ton autoritaire). \eng{Request} (formel) / \eng{Ask for} (standard). \textbf{Règle :} Registre / Nuance pragmatique.
    \item \eng{We \textbf{discussed} the budget.}
    \textbf{Piège :} \textbf{Discuter à propos de} $\neq$ \eng{Discuss about}. \eng{Discuss} est transitif direct. \textbf{Règle :} Calque syntaxique (préposition parasite).
    \item \eng{She \textbf{chaired} / \textbf{facilitated} the workshop brilliantly.}
    \textbf{Piège :} \textbf{Animer} $\neq$ \eng{Animate} (donner vie, dessiner des cartoons). \eng{Chair} (présider formel), \eng{Facilitate} (guider les échanges), \eng{Lead} (diriger), \eng{Moderate} (modérer un débat). \textbf{Règle :} Faux ami sémantique.
    \item \eng{I will \textbf{make the most of} my holidays to revise. / I will \textbf{take advantage of} the break to revise.}
    \textbf{Piège :} \textbf{Profiter} $\neq$ \eng{Profit} (verbe intransitif : \eng{profit from}). \eng{Make the most of} / \eng{Take advantage of} = transitifs directs. \textbf{Règle :} Faux ami grammatical (transitivité) + lexical.
    \item \eng{The project is \textbf{behind schedule}. / The project is \textbf{overdue}.}
    \textbf{Piège :} \textbf{En retard} $\neq$ \eng{In delay} / \eng{Late} (pour un projet). \eng{Behind schedule} (avancement), \eng{Overdue} (dépassé la date limite). \textbf{Règle :} Calque structurel / Collocation.
    \item \eng{Let's \textbf{make a decision} now. / We need to \textbf{decide} now.}
    \textbf{Piège :} \textbf{Prendre une décision} $\neq$ \eng{Take a decision} (calque franglais toléré mais \eng{Make a decision} est l'idiome standard UK/US). \textbf{Règle :} Collocation verbale (\eng{Make} vs \eng{Take}).
\end{enumerate}

\textbf{Synthèse des 10 pièges :} \cle{Attend, Schedule, Workload, Deadline, Request, Discuss (direct), Chair/Facilitate, Make the most of, Behind schedule, Make a decision}.
\end{exoresolu}

\begin{methode}[Exploiter les familles de mots, la prononciation et l'orthographe]
\textbf{Objectif :} Dériver les formes nominales, adjectivales, verbales et adverbiales ; maîtriser l'accentuation tonique (stress shift) et l'orthographe des dérivés.
\begin{enumerate}
    \item \textbf{Construire le tableau de dérivation (Word Formation) :}
    \begin{center}
    \begin{tabularx}{\linewidth}{|X|X|X|X|X|}\hline
    \rowcolor{popgreenL} \textbf{Verbe} & \textbf{Nom (Action/Agent)} & \textbf{Nom (Abstrait/Qualité)} & \textbf{Adjectif} & \textbf{Adverbe} \\ \hline
    \eng{participate} & \eng{participant, participation} & \eng{participation} & \eng{participatory} & \eng{participatorily} (rare) \\ \hline
    \eng{contribute} & \eng{contributor, contribution} & \eng{contribution} & \eng{contributory} & \eng{contributorily} (rare) \\ \hline
    \eng{collaborate} & \eng{collaborator, collaboration} & \eng{collaboration} & \eng{collaborative} & \eng{collaboratively} \\ \hline
    \eng{volunteer} & \eng{volunteer, volunteerism} & \eng{volunteering} & \eng{voluntary} & \eng{voluntarily} \\ \hline
    \eng{prioritise} & \eng{prioritisation} & \eng{priority} & \eng{priority / prioritised} & \eng{prioritarily} \\ \hline
    \eng{schedule} & \eng{schedule, scheduler} & \eng{scheduling} & \eng{scheduled / unscheduled} & --- \\ \hline
    \eng{manage} & \eng{manager, management} & \eng{management} & \eng{managerial / manageable} & \eng{managerially / manageably} \\ \hline
    \eng{organise} & \eng{organiser, organisation} & \eng{organisation} & \eng{organisational / organised} & \eng{organisationally} \\ \hline
    \eng{commit} & \eng{commitment, committee} & \eng{commitment} & \eng{committed} & \eng{committedly} \\ \hline
    \eng{engage} & \eng{engagement} & \eng{engagement} & \eng{engaging / engaged} & \eng{engagingly} \\ \hline
    \end{tabularx}
    \end{center}
    \emph{Note :} Les formes en \textit{-ise/-isation} (UK) sont la norme au Cameroun ; les formes en \textit{-ize/-ization} (US/Oxford) sont acceptées mais moins attendues.
    \item \textbf{Règle d'accentuation (Stress Shift) :} Le suffixe \textbf{-tion / -sion / -ity / -ive / -ial} attire l'accent tonique sur la syllabe qui le précède immédiatement.
    \begin{itemize}
        \item \eng{volun'teer} (verbe/nom) $\rightarrow$ \eng{volun'tary} (adj) $\rightarrow$ \eng{volun'tarily} (adv) \textbf{MAIS} \eng{volunteer'ism}.
        \item \eng{'manage} $\rightarrow$ \eng{manage'ment} $\rightarrow$ \eng{mana'gerial} (accent sur 'ge' avant -ial).
        \item \eng{or'ganise} $\rightarrow$ \eng{organi'sation} $\rightarrow$ \eng{organi'sational}.
        \item \eng{pri'oritise} $\rightarrow$ \eng{prioriti'sation} $\rightarrow$ \eng{pri'ority} (exception : accent fixe sur 1ère pour \eng{priority}).
    \end{itemize}
    \item \textbf{Orthographe : Doublement de consonne / \textit{e} muet :}
    \begin{itemize}
        \item \eng{Commit} $\rightarrow$ \eng{Committed, Committee, Commitment} (doublement \textit{t} : verbe monosyllabe accentué ou bisyllabe accentué sur dernier syllabe = CVC).
        \item \eng{Manage} $\rightarrow$ \eng{Manageable} (garde \textit{e} pour \textit{g} doux /d\eng{ʒ}/), \eng{Management} (perd \textit{e}).
        \item \eng{Schedule} $\rightarrow$ \eng{Scheduling} (garde \textit{e} pour \textit{c} doux /s/ devant \textit{i}).
    \end{itemize}
    \item \textbf{Prononciation clé (approx. fr.) :} \eng{Volunteer} (« vol-on-ti-o »), \eng{Schedule} UK (« che-dyoule ») / US (« ske-doule »), \eng{Prioritise} (« pra-io-ri-taïz »), \eng{Collaborative} (« co-la-bo-ra-tive »), \eng{Commitment} (« co-mit-ment »).
\end{enumerate}
\end{methode}

\begin{exoresolu}[Formation de mots (Word Formation) \& Prononciation]
\textbf{Énoncé :}
\begin{enumerate}
    \item \textbf{Complétez le tableau suivant en dérivant le mot donné en majuscules. Attention à l'orthographe et à la catégorie grammaticale.}
    \begin{center}
    \begin{tabularx}{\linewidth}{|c|X|X|}\hline
    \rowcolor{poppurpleL} \textbf{Racine} & \textbf{Phrase à compléter} & \textbf{Mot formé} \\ \hline
    \eng{ENGAGE} & The students showed great \_\_\_\_\_\_\_\_ in the community project. & \eng{engagement} \\ \hline
    \eng{PARTICIPATE} & \_\_\_\_\_\_\_\_ in the workshop is mandatory for all teachers. & \\ \hline
    \eng{VOLUNTEER} & She works on a \_\_\_\_\_\_\_\_ basis at the local clinic. & \\ \hline
    \eng{PRIORITY} & You must \_\_\_\_\_\_\_\_ your tasks to meet the deadline. & \\ \hline
    \eng{MANAGE} & Time \_\_\_\_\_\_\_\_ is a crucial skill for students. & \\ \hline
    \eng{ORGANISE} & The \_\_\_\_\_\_\_\_ of the event took three months. & \\ \hline
    \eng{COLLABORATE} & It was a highly \_\_\_\_\_\_\_\_ effort between the two schools. & \\ \hline
    \eng{COMMIT} & We need a firm \_\_\_\_\_\_\_\_ from all members. & \\ \hline
    \end{tabularx}
    \end{center}
    \item \textbf{Indiquez la syllabe accentuée (en MAJUSCULE) pour les mots suivants :}
    \eng{volunteer (n/v), voluntary, voluntarily, prioritise, prioritisation, priority, organise, organisation, organisational, commit, commitment, committee.}
    \item \textbf{Dictée / Orthographe :} Écrivez la forme correcte pour les mots sonores suivants (contexte : email professionnel).
    \begin{itemize}
        \item [son: /ˌmænɪdʒmənt/] \_\_\_\_\_\_\_\_
        \item [son: /ˌɒrɡənaɪˈzeɪʃn/] \_\_\_\_\_\_\_\_
        \item [son: /kəˈmɪtmənt/] \_\_\_\_\_\_\_\_
        \item [son: /ˌvɒlənˈteəri/] \_\_\_\_\_\_\_\_ (adjectif)
        \item [son: /ˌvɒlənˈteərəli/] \_\_\_\_\_\_\_\_ (adverbe)
    \end{itemize}
\end{enumerate}

\tcblower

\textbf{Solution détaillée}

\textbf{1. Word Formation (Formation de mots) :}

\begin{center}
\begin{tabularx}{\linewidth}{|c|X|X|X|}\hline
\rowcolor{poppurpleL} \textbf{Racine} & \textbf{Phrase} & \textbf{Mot formé} & \textbf{Analyse morphologique} \\ \hline
\eng{ENGAGE} & The students showed great \_\_\_\_ in the community project. & \textbf{engagement} & Nom abstrait : Verbe + \textbf{-ment}. Pas de changement orthographique. \\ \hline
\eng{PARTICIPATE} & \_\_\_\_ in the workshop is mandatory... & \textbf{Participation} & Nom d'action : Verbe + \textbf{-tion} (chute \textit{e} final). Sujet de phrase $\rightarrow$ Nom singulier. \\ \hline
\eng{VOLUNTEER} & She works on a \_\_\_\_ basis... & \textbf{voluntary} & Adjectif : \eng{Volunteer} $\rightarrow$ \textbf{Voluntary} (racine latine \textit{volunt-}). \eng{Volunteer basis} $\times$. \\ \hline
\eng{PRIORITY} & You must \_\_\_\_ your tasks... & \textbf{prioritise} (UK) / \textbf{prioritize} (US) & Verbe : Nom \eng{Priority} + \textbf{-ise/-ize}. Orthographe : \textit{y} $\rightarrow$ \textit{i} + \textit{se/ze}. \\ \hline
\eng{MANAGE} & Time \_\_\_\_ is a crucial skill... & \textbf{management} & Nom : Verbe + \textbf{-ment}. Chute \textit{e} muet (\eng{Manage} $\rightarrow$ \eng{Manag-ment}). \\ \hline
\eng{ORGANISE} & The \_\_\_\_ of the event took three months. & \textbf{organisation} (UK) / \textbf{organization} (US) & Nom d'action : Verbe + \textbf{-ation} (chute \textit{e}). \textit{s} UK / \textit{z} US. \\ \hline
\eng{COLLABORATE} & It was a highly \_\_\_\_ effort... & \textbf{collaborative} & Adjectif : Verbe + \textbf{-ive} (chute \textit{e}). \eng{Collaboratory} n'existe pas (c'est un lieu). \\ \hline
\eng{COMMIT} & We need a firm \_\_\_\_ from all members. & \textbf{commitment} & Nom : Verbe + \textbf{-ment}. \textbf{Doublement du \textit{t}} (CVC accentué final : com-\textbf{MIT}). \\ \hline
\end{tabularx}
\end{center}

\textbf{2. Accentuation tonique (Word Stress) :}
\begin{center}
\begin{tabularx}{\linewidth}{|X|X|}\hline
\rowcolor{poppurpleL} \textbf{Mot} & \textbf{Accentuation (Syllabe forte en MAJUSCULE)} \\ \hline
\eng{volunTEER} (n/v) & Dernière syllabe (verbe/nom agent) \\ \hline
\eng{VOLun-tary} (adj) & \textbf{1ère} syllabe (règle -ary) \\ \hline
\eng{vol-un-TAR-i-ly} (adv) & \textbf{3ème} syllabe (avant -ily) \\ \hline
\eng{pri-OR-i-tise} & \textbf{3ème} syllabe (avant -ise) \\ \hline
\eng{pri-or-i-ti-SA-tion} & \textbf{6ème} syllabe (avant -tion) \\ \hline
\eng{PRI-or-i-ty} & \textbf{1ère} syllabe (exception \eng{priority}) \\ \hline
\eng{OR-gan-ise} & \textbf{1ère} syllabe \\ \hline
\eng{or-gan-i-SA-tion} & \textbf{5ème} syllabe (avant -tion) \\ \hline
\eng{or-gan-i-sa-TION-al} & \textbf{6ème} syllabe (avant -al) \\ \hline
\eng{com-MIT} & \textbf{Dernière} syllabe \\ \hline
\eng{com-MIT-ment} & \textbf{2ème} syllabe (racine verbale conservée) \\ \hline
\eng{com-MIT-tee} & \textbf{2ème} syllabe (racine verbale conservée) \\ \hline
\end{tabularx}
\end{center}
\emph{Astuce :} Les suffixes \textbf{-tion, -sion, -ic, -ical, -ity, -ive, -ial, -ian} sont « \cle{stress-attracting} » : ils tirent l'accent sur la syllabe qui les précède immédiatement. Les suffixes \textbf{-ment, -ness, -less, -ful, -ly, -er, -or, -ism, -ist} sont « \cle{stress-neutral} » : l'accent reste sur la racine.

\textbf{3. Dictée / Orthographe :}
\begin{enumerate}
    \item \textbf{Management} (\eng{Manage} + \textit{ment} $\rightarrow$ chute \textit{e}).
    \item \textbf{Organisation} / \textbf{Organization} (\eng{Organise} + \textit{ation} $\rightarrow$ chute \textit{e} ; \textit{s} UK / \textit{z} US).
    \item \textbf{Commitment} (\eng{Commit} + \textit{ment} $\rightarrow$ \textbf{doublement \textit{t}}).
    \item \textbf{Voluntary} (Racine latine \textit{volunt-} + \textit{ary}. Pas de \textit{e} après \textit{t}). \emph{Prononciation :} « vol-on-te-ri ».
    \item \textbf{Voluntarily} (\eng{Voluntary} + \textit{ly} $\rightarrow$ \textit{y} $\rightarrow$ \textit{i} + \textit{ly}). \emph{Prononciation :} « vol-on-te-reu-li » (5 syllabes).
\end{enumerate}

\textbf{Conclusion :} La dérivation est un levier majeur pour le \textbf{Use of English} (Partie grammaire/vocabulaire du Bac). Maîtriser les suffixes productifs (\textit{-tion, -ment, -ive, -al, -ise, -ary}) et leurs effets sur l'orthographe (chute \textit{e}, doublement, \textit{y}$\rightarrow$\textit{i}) et l'accentuation garantit des points faciles.
\end{exoresolu}

\begin{methode}[Réemployer le lexique en contexte : Rédaction d'un compte-rendu (Minutes) / Email formel]
\textbf{Objectif :} Produire un écrit structuré (compte-rendu de réunion ou email de suivi) intégrant le vocabulaire de la participation, de la gestion du temps, les collocations et la grammaire (passif, discours indirect, connecteurs).
\begin{enumerate}
    \item \textbf{Analyser la situation de communication (TAPE) :}
    \begin{itemize}
        \item \textbf{T}opic : Suivi d'un projet communautaire / scolaire / professionnel.
        \item \textbf{A}udience : Supérieur (principal, chef de service), pairs (collègues, membres), institution (conseil municipal). $\rightarrow$ Registre \textbf{formel / semi-formel}.
        \item \textbf{P}urpose : Informer des décisions, assigner des tâches (\eng{action points}), fixer des échéances (\eng{deadlines}), demander validation.
        \item \textbf{E}format : \textbf{Minutes} (structure codifiée) ou \textbf{Formal Email} (structure : Objet, Salutation, Contexte, Points clés, Action required, Clôture).
    \end{itemize}
    \item \textbf{Structure type \textit{Minutes} (Compte-rendu) :}
    \begin{enumerate}
        \item \textbf{Header :} Name of body, Date, Time, Venue, Attendees / Apologies (absents excusés), Absent (non excusés).
        \item \textbf{Opening :} \eng{The meeting was called to order by [Chair] at [Time].}
        \item \textbf{Minutes of previous meeting :} \eng{Adopted / Amended.}
        \item \textbf{Matters Arising / Agenda Items :} Numérotés. Utiliser le \textbf{Passif} (\eng{It was decided that...}, \eng{The report was presented by...}) ou \textbf{Reported Speech} (\eng{Mr X reported that...}).
        \item \textbf{Action Points (Crucial) :} Tableau : \eng{Action Item | Responsible Person | Deadline | Status}.
        \item \textbf{AOB (Any Other Business) :} Questions diverses.
        \item \textbf{Closing :} \eng{The meeting was adjourned at [Time]. Next meeting: [Date].}
        \item \textbf{Signature :} \eng{Prepared by [Secretary] / Approved by [Chair].}
    \end{enumerate}
    \item \textbf{Lexique obligatoire à insérer (Checklist) :} \eng{Agenda, Attendees, Apologies, Minutes, Matters arising, Action points, Deadline, Follow-up, Stakeholder, Chairperson, Facilitator, Quorum, Adjourned, Carried (vote), Unanimous.}
    \item \textbf{Relecture (Checklist BAC) :} Temps verbaux (Passé simple/Prétérit pour faits, Passif pour décisions), Prépositions (\eng{on} date, \eng{by} deadline, \eng{in} month), Faux amis (\eng{attend} vs \eng{assist}), Orthographe (\eng{commitment, schedule}).
\end{enumerate}
\end{methode}

\begin{exoresolu}[Production écrite : Rédaction de Minutes (Compte-rendu)]
\textbf{Énoncé :} Vous êtes le secrétaire du \textit{« Green Campus Club »} de votre lycée à Yaoundé. Rédigez le \textbf{compte-rendu (Minutes)} de la réunion tenue le samedi 12 octobre 2024, basée sur les notes brutes ci-dessous. Respectez la structure formelle. (180-220 mots).

\textbf{Notes brutes :}
\begin{itemize}
    \item Présents : Président (Amadou), Secrétaire (vous), Trésorière (Aïcha), 15 membres. Absents excusés : Vice-président (Paul - malade). Absent non excusé : Coordinateur logistique (Jean).
    \item Président ouvre séance 9h. Approuve PV réunion précédente.
    \item Point 1 : Campagne plantation arbres « Green October » prévue 26 oct. Budget 50 000 FCFA. Trésorière dit fonds insuffisants. Décision : organiser vente gâteaux vendredi 18 oct pour \textbf{lever des fonds} (fundraise). Aïcha \textbf{dirige} (spearhead) comité financement.
    \item Point 2 : Partenariat avec Mairie Yaoundé VII. M. le Maire \textbf{a promis} (pledged) 20 plants. Président \textbf{remercie}.
    \item Point 3 : Problème gestion temps : réunions trop longues. Décision : limiter à 1h30 max. \textbf{Ordre du jour} (agenda) à envoyer 48h avant.
    \item Fin séance 10h45. Prochaine réunion : 2 nov.
\end{itemize}

\tcblower

\textbf{Solution détaillée : Modèle de Minutes (Formal Style)}

\begin{textebox}[Minutes of the Green Campus Club Meeting]
\textbf{GREEN CAMPUS CLUB -- LYCÉE DE YAOUNDÉ}

\textbf{Minutes of the Ordinary General Meeting held on Saturday, 12th October 2024}

\textbf{Venue:} School Library, Room 101 \hfill \textbf{Time:} 09:00 -- 10:45

\textbf{1. ATTENDANCE}
\textbf{Chairperson:} Amadou Bello \\
\textbf{Secretary:} [Votre Nom] \\
\textbf{Treasurer:} Aïcha Diko \\
\textbf{Members Present:} 15 registered members \\
\textbf{Apologies for Absence:} Paul Biya (Vice-Chairperson) -- medical leave. \\
\textbf{Absent without Apology:} Jean Mbarga (Logistics Coordinator).

\textbf{2. OPENING}
The meeting was \textbf{called to order} by the Chairperson, Amadou Bello, at 09:00. A \textbf{quorum} was declared present. The Chairperson welcomed all \textbf{attendees} and thanked them for their \textbf{commitment}.

\textbf{3. MINUTES OF PREVIOUS MEETING}
The minutes of the meeting held on 14th September 2024 were \textbf{read, adopted unanimously}, and \textbf{signed} by the Chairperson and Secretary. \textbf{Matters arising:} None.

\textbf{4. AGENDA ITEMS}

\textbf{4.1 ``Green October'' Tree Planting Campaign (26th October)}
The Treasurer, Aïcha Diko, \textbf{presented the financial report}. She highlighted a budget shortfall for the 50,000 FCFA target. 
\textbf{Decision:} A \textbf{fundraising} bake sale shall be organized on Friday, 18th October. 
\textbf{Action Point:} Aïcha Diko to \textbf{spearhead} the Fundraising Sub-committee (Logistics, Communication, Procurement). \textbf{Deadline:} 17th October for final preparation.

\textbf{4.2 Partnership with Yaoundé VII Council}
The Chairperson \textbf{reported} that the Mayor has \textbf{pledged} 20 seedlings. The Club \textbf{expressed its gratitude}. 
\textbf{Action Point:} Chairperson to \textbf{follow up} with Council liaison officer by 20th October for delivery logistics.

\textbf{4.3 Time Management \& Meeting Efficiency}
Members \textbf{raised concerns} regarding excessive meeting duration. 
\textbf{Decision:} Future meetings are \textbf{capped at 90 minutes maximum}. The \textbf{agenda} and supporting documents must be \textbf{circulated} at least 48 hours \textbf{in advance}.

\textbf{5. ANY OTHER BUSINESS (AOB)}
None raised.

\textbf{6. CLOSURE}
The meeting was \textbf{adjourned} at 10:45 after the Chairperson \textbf{summarised the action points}. 
\textbf{Next Meeting:} Saturday, 2nd November 2024, 09:00, School Library.

\vspace{0.5cm}
\textbf{Prepared by:} \hrulefill \hspace{2cm} \textbf{Approved by:} \hrulefill \\
{}[Votre Nom], Secretary \hfill Amadou Bello, Chairperson \\
Date: 12/10/2024 \hfill Date: 12/10/2024
\end{textebox}

\textbf{Analyse linguistique du modèle (Pourquoi c'est une copie « Bac ») :}
\begin{itemize}
    \item \textbf{Structure :} Respect strict du format standard (Header, Attendance, Opening, Minutes Prev, Agenda Items numérotés, Action Points distincts, AOB, Closure, Next Meeting, Signatures).
    \item \textbf{Grammaire :} Usage massif du \textbf{Passif} (\eng{was called to order, were read, adopted, signed, presented, shall be organized, are capped, must be circulated, was adjourned}) $\rightarrow$ Impersonnalité, objectivité, focus sur l'action/décision.
    \item \textbf{Vocabulaire cible (Coché) :} \eng{Attendees, Apologies, Quorum, Unanimously, Matters arising, Financial report, Shortfall, Fundraising, Spearhead, Pledged, Follow up, Circulated, In advance, Capped at, Adjourned, Action points}.
    \item \textbf{Collocations :} \eng{Call to order, Declare a quorum, Raise concerns, Express gratitude, Cap at, Circulate agenda}.
    \item \textbf{Zéro faute :} Pas de \eng{assist} pour \eng{attend}, pas de \eng{program} pour \eng{schedule}, pas de \eng{delay} pour \eng{deadline}, pas de \eng{discuss about}.
    \item \textbf{Registre :} Formel, concis, précis. Pas de « I think », « We decided » (sauf reported speech), mais « It was decided », « The meeting agreed ».
\end{itemize}
\end{exoresolu}

\begin{attention}[Les 5 erreurs qui coûtent des points au Bac (Leçon 23)]
\begin{enumerate}
    \item \textbf{« J'ai assisté à la réunion » $\rightarrow$ \eng{I \textbf{assisted} the meeting.} \textcolor{red}{FAUX.}} \\
    \eng{Assist} = Aider (\eng{assist someone}). Le verbe pour « être présent » est \textbf{\eng{Attend}} (\eng{Attend a meeting / class / conference}). \textit{Astuce :} \eng{Assistant} = Assistant (adjoint), pas « participant ».
    
    \item \textbf{« Le délai est vendredi » $\rightarrow$ \eng{The \textbf{delay} is Friday.} \textcolor{red}{FAUX.}} \\
    \eng{Delay} = Retard (noun). « Délai / Date limite » = \textbf{\eng{Deadline}} ou \textbf{\eng{Cut-off date}} ou \textbf{\eng{Time limit}}. \eng{The deadline falls on Friday.}
    
    \item \textbf{« J'ai une charge de travail énorme » $\rightarrow$ \eng{I have a huge \textbf{charge} of work.} \textcolor{red}{FAUX.}} \\
    \eng{Charge} = Accusation, responsabilité, frais, électricité. « Charge de travail » = \textbf{\eng{Workload}} (nom composé) ou \textbf{\eng{Heavy workload}}.
    
    \item \textbf{« Nous avons discuté à propos du budget » $\rightarrow$ \eng{We \textbf{discussed about} the budget.} \textcolor{red}{FAUX.}} \\
    \eng{Discuss} est \textbf{transitif direct} (pas de préposition). \eng{We \textbf{discussed} the budget.} (Mais : \eng{We \textbf{talked about} / \textbf{spoke about} the budget}).
    
    \item \textbf{« Le projet est en retard » $\rightarrow$ \eng{The project is \textbf{in delay} / \textbf{late}.} \textcolor{red}{FAUX.}} \\
    \eng{In delay} n'existe pas. \eng{Late} s'applique aux personnes/heures. Pour un projet/tâche : \textbf{\eng{Behind schedule}} ou \textbf{\eng{Overdue}} (si date limite dépassée). \eng{We are running behind schedule.}
\end{enumerate}
\textbf{Autres pièges « Bonus » :} \eng{Program} (pour \eng{Schedule/Plan}), \eng{Profit} (verbe transitif pour \eng{Make the most of}), \eng{Animate} (pour \eng{Chair/Facilitate}), \eng{Take a decision} (calque pour \eng{Make a decision}), \eng{Participate} sans \textbf{\eng{in}}.
\end{attention}

\newpage

\coursec{Exercices}

\courssub{Je vérifie mes connaissances}

\begin{exercice}[QCM : Le bon mot pour le bon contexte]
Choisissez la réponse correcte (a, b ou c) pour compléter chaque phrase.
\begin{enumerate}
\item \eng{The \cle{deadline} for submitting the application form is next Monday.}
\begin{enumerate}
\item \textbf{a)} delay \item \textbf{b)} deadline \item \textbf{c)} diary
\end{enumerate}
\item \eng{Our team managed to finish the project \cle{in the nick of time}.}
\begin{enumerate}
\item \textbf{a)} on time \item \textbf{b)} in the nick of time \item \textbf{c)} at times
\end{enumerate}
\item \eng{The meeting has been \cle{postponed} until further notice.}
\begin{enumerate}
\item \textbf{a)} cancelled \item \textbf{b)} postponed \item \textbf{c)} advanced
\end{enumerate}
\item \eng{Every member must \cle{pull their weight} if we want to succeed.}
\begin{enumerate}
\item \textbf{a)} pull their weight \item \textbf{b)} pull their leg \item \textbf{c)} pull through
\end{enumerate}
\item \eng{She is very \cle{punctual}; she always arrives five minutes early.}
\begin{enumerate}
\item \textbf{a)} punctual \item \textbf{b)} punctuality \item \textbf{c)} punctually
\end{enumerate}
\end{enumerate}
\end{exercice}

\begin{aretenir}[Mots-clés : Participation et Gestion du temps]
\begin{center}
\begin{tabularx}{\linewidth}{|X|X|X|}
\hline
\rowcolor{popblueL} \textbf{English} & \textbf{Français} & \textbf{Collocation / Exemple} \\
\hline
\cle{deadline} & date limite, échéance & \eng{meet a deadline} (respecter une échéance), \eng{miss a deadline} \\
\cle{schedule} / \cle{timetable} & emploi du temps, planning & \eng{a tight schedule} (planning serré), \eng{draw up a schedule} \\
\cle{volunteer} (n./v.) & bénévole / être volontaire & \eng{volunteer for a task} (se porter volontaire) \\
\cle{coordinator} & coordonnateur(trice) & \eng{The coordinator briefed the team.} \\
\cle{participant} & participant(e) & \cle{participation} (n.), \cle{participate} (v.) \\
\cle{punctual} / \cle{punctuality} & ponctuel(le) / ponctualité & \eng{arrive on time} / \eng{be punctual} \\
\cle{postpone} & reporter & \eng{postpone to/until} + date \\
\cle{commitment} & engagement, dévouement & \eng{show commitment} \\
\hline
\end{tabularx}
\end{center}
\end{aretenir}

\begin{exercice}[Vrai ou Faux justifié : « A Busy Week at GHS Limbe »]
Lisez le texte ci-dessous, puis dites si les affirmations sont vraies (V) ou fausses (F). Justifiez chaque réponse par une phrase ou un extrait du texte.
\begin{textebox}[A Busy Week at GHS Limbe]
Last week, the Students' Council of Government High School Limbe organised a “Clean Campus Campaign”. The \cle{schedule} was tight: activities ran \cle{around the clock} from Monday to Friday. On Monday, the \cle{coordinator} briefed the \cle{volunteers} on safety rules. Tuesday was dedicated to collecting plastic waste around the playground. Unfortunately, the heavy rain \cle{forced} us to \cle{postpone} the tree-planting session to Thursday. By Friday, the school grounds looked spotless. The principal praised the students for their \cle{commitment} and \cle{punctuality}. “You have shown that teamwork makes the dream work,” he said.

\medskip
\textbf{Glossaire :} \textit{schedule (n.) : emploi du temps ; around the clock : jour et nuit / sans arrêt ; coordinator (n.) : coordonnateur ; volunteer (n.) : bénévole ; to postpone (v.) : reporter ; commitment (n.) : engagement ; punctuality (n.) : ponctualité.}
\end{textebox}
\begin{enumerate}
\item The campaign lasted for two weeks only.
\item The tree-planting session took place on Tuesday as planned.
\item The principal congratulated the students on their attitude.
\item Volunteers received safety instructions on Monday.
\end{enumerate}
\end{exercice}

\begin{attention}[Faux amis fréquents : \textit{Assist} vs \textit{Attend}, \textit{Actually} vs \textit{Currently}, \textit{Delay} vs \textit{Deadline}, \textit{Profit} vs \textit{Benefit/Take advantage}, \textit{Responsible of} vs \textit{Responsible for}]
Ne traduisez pas mot à mot depuis le français !
\begin{center}
\begin{tabularx}{\linewidth}{|X|X|X|}
\hline
\rowcolor{poporangeL} \textbf{Français} & \textbf{Anglais correct} & \textbf{Piège (Faux ami / Calque)} \\
\hline
Assister à (être présent) & \cle{attend} & \textit{Assist} = aider (help) \\
Actuellement & \cle{currently} / \cle{at present} & \textit{Actually} = en fait, en réalité \\
Date limite, échéance & \cle{deadline} & \textit{Delay} = retard \\
Tirer profit de, bénéficier de & \cle{benefit from} / \cle{take advantage of} & \textit{Profit} (n.) = bénéfice ; (v.) = rapporter du profit \\
Être responsable de & \cle{be responsible for} & \textit{Responsible of} n'existe pas (préposition \textbf{for}) \\
\hline
\end{tabularx}
\end{center}
\end{attention}

\begin{exercice}[Faux amis et calques du français : Correction]
Les phrases suivantes contiennent des erreurs dues à l'interférence du français. Corrigez-les et expliquez l'erreur (faux ami ou calque).
\begin{enumerate}
\item I will \textbf{assist} to the meeting tomorrow.
\item \textbf{Actually}, the shops are closed on Sundays. (Meaning: \textit{Actuellement})
\item The \textbf{delay} for the homework is Friday.
\item We need to \textbf{profit} from this opportunity to learn.
\item He is \textbf{responsible} of the cleaning team.
\end{enumerate}
\end{exercice}

\begin{aretenir}[Formation de mots : Familles lexicales du thème]
\begin{center}
\begin{tabular}{|l|l|l|l|}
\hline
\rowcolor{popgreenL} \textbf{Verbe} & \textbf{Nom (chose/action)} & \textbf{Nom (agent)} & \textbf{Adjectif / Adverbe} \\
\hline
\cle{participate} & \cle{participation} & \cle{participant} & \cle{participatory} / \cle{actively} \\
\cle{organise} (UK) & \cle{organisation} & \cle{organiser} & \cle{organised} / \cle{organisedly} \\
\cle{punctuate} & \cle{punctuality} & -- & \cle{punctual} / \cle{punctually} \\
\cle{volunteer} & \cle{volunteering} & \cle{volunteer} & \cle{voluntary} / \cle{voluntarily} \\
\cle{commit} & \cle{commitment} & -- & \cle{committed} \\
\cle{manage} & \cle{management} & \cle{manager} & \cle{manageable} \\
\hline
\end{tabular}
\end{center}
\emph{Note :} L'orthographe britannique (\textit{organise, organisation}) est la norme au Cameroun.
\end{aretenir}

\begin{exercice}[Formation de mots et grammaire contextuelle]
Complétez les phrases avec la forme correcte du mot entre parenthèses (nom, adjectif, adverbe ou verbe conjugué).
\begin{enumerate}
\item The \_\_\_\_\_\_\_\_\_\_ (participate) of every student is compulsory for the sports day.
\item She is a very \_\_\_\_\_\_\_\_\_\_ (punctual) person; she is never late.
\item The \_\_\_\_\_\_\_\_\_\_ (organize) of the event did a great job.
\item If we \_\_\_\_\_\_\_\_\_\_ (not / manage) our time well, we will miss the \_\_\_\_\_\_\_\_\_\_ (deadline).
\item All \_\_\_\_\_\_\_\_\_\_ (participate) must register before Wednesday.
\end{enumerate}
\end{exercice}

\courssub{Je m'entraîne}

\begin{exercice}[Traduction : Du français vers l'anglais]
Traduisez les phrases suivantes en anglais en réemployant le vocabulaire de la leçon.
\begin{enumerate}
\item La date limite pour s'inscrire est fixée au 15 mai.
\item Les bénévoles ont travaillé jour et nuit pour préparer la salle.
\item Nous avons dû reporter le match à cause de la pluie.
\item Chaque élève doit participer aux tâches ménagères (pull his/her weight).
\item Le coordonnateur a remercié tout le monde pour leur engagement.
\end{enumerate}
\end{exercice}

\begin{aretenir}[Structures utiles pour la transformation]
\begin{itemize}
\item \cle{Due to} / \cle{Owing to} / \cle{Because of} + nom / gérondif (cause) $\neq$ \textit{Because} + sujet + verbe.
\item \cle{The deadline for} + gérondif / nom \textbf{is} + date.
\item \cle{Work around the clock} = travailler jour et nuit / sans relâche.
\item \cle{Never} + adjectif positif (late) = \textit{Always} + adjectif négatif (on time / punctual).
\item \textit{Reported speech} : “Can you help me...?” $\rightarrow$ She asked me \textbf{to help} her (infinitif). \textit{This/These} $\rightarrow$ \textit{That/Those}.
\end{itemize}
\end{aretenir}

\begin{exercice}[Transformation de phrases]
Réécrivez les phrases en commençant par les mots indiqués, sans en changer le sens.
\begin{enumerate}
\item The meeting was postponed because of the strike. \\
\textit{Due to \_\_\_\_\_\_\_\_\_\_}
\item You must submit your report before Friday. \\
\textit{The deadline \_\_\_\_\_\_\_\_\_\_}
\item They worked all night to finish on time. \\
\textit{They worked around \_\_\_\_\_\_\_\_\_\_}
\item He always arrives on time. \\
\textit{He is never \_\_\_\_\_\_\_\_\_\_}
\item “Can you help me carry these boxes?” she asked. \\
\textit{She asked me to \_\_\_\_\_\_\_\_\_\_}
\end{enumerate}
\end{exercice}

\begin{exercice}[Texte à trous : Banque de mots]
Complétez le texte ci-dessous en choisissant les mots appropriés dans la banque. Attention : deux mots sont en trop.
\begin{center}
\begin{tabular}{|c|c|c|c|c|c|c|c|c|c|}
\hline
\textbf{volunteers} & \textbf{schedule} & \textbf{postponed} & \textbf{take part} & \textbf{on time} & \textbf{lend a hand} & \textbf{deadline} & \textbf{coordinator} & \textbf{cancelled} & \textbf{punctual} \\
\hline
\end{tabular}
\end{center}
\begin{textebox}[School Fundraising Day]
The \_\_\_\_\_\_\_\_\_\_ for the annual fundraising day was published last week. The main \_\_\_\_\_\_\_\_\_\_ reminded everyone that the \_\_\_\_\_\_\_\_\_\_ for registration is Tuesday. Over fifty \_\_\_\_\_\_\_\_\_\_ have already signed up to \_\_\_\_\_\_\_\_\_\_ in the various activities. Unfortunately, the football tournament was \_\_\_\_\_\_\_\_\_\_ due to a waterlogged pitch. We hope the new date will allow everyone to arrive \_\_\_\_\_\_\_\_\_\_. Please \_\_\_\_\_\_\_\_\_\_ wherever you can!
\end{textebox}
\end{exercice}

\begin{exercice}[Appariement : Expressions idiomatiques et gestion du temps]
Associez chaque expression idiomatique (1--8) à sa signification française (A--H).
\begin{enumerate}
\item In the nick of time
\item Around the clock
\item Pull one's weight
\item Beat the clock
\item Against the clock
\item Time is money
\item A race against time
\item On the dot
\end{enumerate}
\begin{itemize}
\item[A)] Travailler sans relâche / 24h sur 24
\item[B)] Faire sa part du travail
\item[C)] Juste à temps / au dernier moment
\item[D)] Précisément à l'heure dite
\item[E)] Le temps, c'est de l'argent
\item[F)] Faire quelque chose le plus vite possible avant une échéance
\item[G)] Une course contre la montre
\item[H)] Essayer de finir avant l'heure limite
\end{itemize}
\end{exercice}

\courssub{Je me prépare au Bac}

\begin{exercice}[Compréhension et langue : « The Buea Youth Entrepreneurship Summit »]
Lisez attentivement le texte suivant et répondez aux questions.
\begin{textebox}[The Buea Youth Entrepreneurship Summit]
From the 10th to the 12th of August, the University of Buea hosted the third edition of the Youth Entrepreneurship Summit. The \cle{theme} was “Innovation and Time Management”. Over 300 young \cle{participants} from the Southwest and Northwest regions gathered to \cle{share ideas} and \cle{network}.

The \cle{programme} was intense. On Day 1, the \cle{keynote speaker}, Mrs. Ekema, stressed the importance of \cle{punctuality} in business. “If you cannot manage your time, you cannot manage a business,” she declared. Workshops ran \cle{simultaneously} in three halls. Day 2 focused on practical skills: writing a business plan, \cle{pitching} to investors, and using digital tools for \cle{scheduling}. A highlight was the “Elevator Pitch Contest” where contestants had exactly 60 seconds to convince the jury.

On the final day, the \cle{award ceremony} took place. The winner, a 22-year-old student from Bamenda, received a grant of 500\,000 FCFA. In his speech, he thanked the \cle{organisers} and urged his peers to \cle{seize opportunities} and \cle{meet deadlines}. The summit ended with a commitment to create a \cle{mentorship} network for young startups.

\medskip
\textbf{Glossaire :} \textit{keynote speaker (n.) : conférencier principal ; to pitch (v.) : présenter un projet de façon concise ; scheduling (n.) : planification ; grant (n.) : subvention ; to urge (v.) : exhorter ; mentorship (n.) : mentorat / tutorat.}
\end{textebox}

\textbf{A. Comprehension Questions} (Answer in your own words as far as possible)
\begin{enumerate}
\item Where and when did the summit take place?
\item What was the main message of the keynote speaker?
\item Describe the “Elevator Pitch Contest”.
\item What did the winner receive and what did he advise others to do?
\end{enumerate}

\textbf{B. Vocabulary in Context}
Find words or expressions in the text that mean:
\begin{enumerate}
\item People who take part in an event (para 1)
\item At the same time (para 2)
\item To take advantage of chances (para 3)
\item Guidance from an experienced person (para 3)
\end{enumerate}

\textbf{C. Grammar / Word Formation}
\begin{enumerate}
\item Rewrite in the passive voice: “The organisers \textbf{hosted} the summit at the University of Buea.”
\item Give the noun form of: \textbf{punctual} ; \textbf{participate} ; \textbf{organise}.
\item Complete with the right preposition: The workshop focused \_\_\_\_\_\_ digital tools. / He insisted \_\_\_\_\_\_ the importance of time.
\end{enumerate}
\end{exercice}

\begin{exercice}[Traduction : Français $\rightarrow$ Anglais]
Traduisez le paragraphe suivant en anglais. (Environ 80 mots)
« L'organisation d'un événement scolaire demande une bonne planification. Le coordinateur doit établir un emploi du temps précis et fixer une date limite pour les inscriptions. Les bénévoles doivent arriver à l'heure et travailler sans relâche pour que tout soit prêt. Si une activité est reportée, il faut en informer tout le monde rapidement. C'est seulement ainsi que l'événement sera une réussite. »
\end{exercice}

\begin{exercice}[Expression écrite : Sujets de type Bac]
\textbf{Choisissez UN des deux sujets ci-dessous. Écrivez environ 100--120 mots.}

\textbf{Sujet 1 (Article de magazine scolaire) :}
You are the president of the Environment Club. Write an article for your school magazine inviting students to take part in a “Clean-Up Campaign” next Saturday. Mention the date, the deadline for signing up, the programme (gathering, distribution of tools, cleaning zones, refreshments) and the benefits for the community.

\textbf{Sujet 2 (Lettre formelle / Email) :}
You are the Class Prefect of Upper Sixth A. Write a formal email to the Principal requesting permission to organise a “Career Day” for final year students. Explain the objectives, propose a date and a draft schedule, and ask for the school's support (venue, invited professionals, time off classes).

\medskip
\textbf{Activité orale guidée (Jeu de rôle) :}
En binôme. \textbf{Élève A :} Vous êtes le coordonnateur de la « Journée Culturelle » de l'école. Vous accueillez un volontaire. Demandez ses disponibilités, expliquez le planning, donnez la date limite d'inscription, insistez sur la ponctualité. \textbf{Élève B :} Vous êtes un élève volontaire. Posez des questions sur les tâches, proposez votre aide (\textit{lend a hand}), confirmez votre engagement. \textit{Utilisez au moins 5 expressions de la leçon (deadline, schedule, volunteer, on time, take part, postpone, pull one's weight, in the nick of time...).}
\end{exercice}

\begin{savaistu}[Culture \& Contexte : Le bénévolat au Cameroun et dans le monde anglophone]
Au Cameroun, l'esprit de \cle{volunteering} (bénévolat) est fort : \textit{“Njangi”} (tontines solidaires), nettoyage des quartiers le samedi matin (\textit{“Clean-up Saturday”} imposé par certains maires), brigades de jeunesse. Dans le monde anglophone, le \cle{volunteering} est souvent valorisé sur le CV (\textit{resume}) pour l'université ou l'emploi (\textit{“Community service”}). Des organisations comme \textit{VSO} (Voluntary Service Overseas) ou le \textit{Peace Corps} (USA) envoient des volontaires. L'expression \eng{“Many hands make light work”} (Plusieurs mains rendent le travail léger / L'union fait la force) résume cet esprit.
\end{savaistu}

\newpage

\coursec{Corrigés des exercices}

\begin{corrige}[Exercice 1 — QCM : Le bon mot pour le bon contexte]
\begin{enumerate}
\item \textbf{b)} \cle{deadline} (date limite, échéance). \emph{Delay} = retard ; \emph{diary} = agenda / journal intime. La phrase signifie : \eng{The deadline for submitting the application form is next Monday.} (« La date limite pour soumettre le formulaire de candidature est lundi prochain. »)
\item \textbf{b)} \cle{in the nick of time} (juste à temps, au dernier moment). \emph{On time} = à l'heure (ponctualité normale) ; \emph{at times} = parfois. Traduction : \eng{Our team managed to finish the project in the nick of time.} (« Notre équipe a réussi à finir le projet juste à temps. »)
\item \textbf{b)} \cle{postponed} (reporté). \emph{Cancelled} = annulé ; \emph{advanced} = avancé (dans le temps). Traduction : \eng{The meeting has been postponed until further notice.} (« La réunion a été reportée jusqu'à nouvel ordre. »)
\item \textbf{a)} \cle{pull their weight} (faire sa part du travail). \emph{Pull someone's leg} = se moquer de quelqu'un ; \emph{pull through} = s'en sortir (maladie, difficulté). Traduction : \eng{Every member must pull their weight if we want to succeed.} (« Chaque membre doit faire sa part si nous voulons réussir. »)
\item \textbf{a)} \cle{punctual} (adjectif : ponctuel). \emph{Punctuality} = nom (ponctualité) ; \emph{punctually} = adverbe (ponctuellement). Après le verbe \emph{is}, il faut un adjectif attribut. Traduction : \eng{She is very punctual; she always arrives five minutes early.} (« Elle est très ponctuelle ; elle arrive toujours cinq minutes en avance. »)
\end{enumerate}
\textbf{Point de vigilance :} distinguer toujours la \emph{nature grammaticale} du mot attendu (nom, adjectif, adverbe) avant de choisir.
\end{corrige}

\begin{corrige}[Exercice 2 — Vrai ou Faux justifié]
\begin{enumerate}
\item \textbf{Faux.} Justification : \eng{“activities ran around the clock from Monday to Friday”} — la campagne a duré une semaine (lundi à vendredi), pas deux.
\item \textbf{Faux.} Justification : \eng{“the heavy rain forced us to postpone the tree-planting session to Thursday”} — la séance de plantation d'arbres n'a pas eu lieu mardi ; elle a été reportée à jeudi.
\item \textbf{Vrai.} Justification : \eng{“The principal praised the students for their commitment and punctuality”} — le proviseur a félicité les élèves pour leur attitude (engagement et ponctualité).
\item \textbf{Vrai.} Justification : \eng{“On Monday, the coordinator briefed the volunteers on safety rules”} — les bénévoles ont bien reçu des consignes de sécurité le lundi.
\end{enumerate}
\textbf{Méthode :} pour un Vrai/Faux justifié, citez toujours l'extrait exact du texte entre guillemets anglais, puis reformulez en français si demandé.
\end{corrige}

\begin{corrige}[Exercice 3 — Faux amis et calques : Correction]
\begin{enumerate}
\item \textbf{Correction :} \eng{I will \cle{attend} the meeting tomorrow.} (« J'assisterai à la réunion demain. ») \quad \emph{Erreur :} faux ami. \emph{Assist} = aider (\emph{help}) ; \emph{attend} = assister à, être présent à. Notez aussi : \emph{attend} se construit \textbf{sans préposition} (jamais \emph{attend to} pour « assister à »).
\item \textbf{Correction :} \eng{\cle{Currently} (ou \cle{At present}), the shops are closed on Sundays.} (« Actuellement, les magasins sont fermés le dimanche. ») \quad \emph{Erreur :} faux ami. \emph{Actually} = en fait, en réalité ; \emph{currently / at present} = actuellement.
\item \textbf{Correction :} \eng{The \cle{deadline} for the homework is Friday.} (« La date limite pour les devoirs est vendredi. ») \quad \emph{Erreur :} faux ami. \emph{Delay} = retard ; \emph{deadline} = date limite, échéance.
\item \textbf{Correction :} \eng{We need to \cle{benefit from} (ou \cle{take advantage of}) this opportunity to learn.} (« Nous devons profiter de cette occasion pour apprendre. ») \quad \emph{Erreur :} calque du français « profiter de ». Le verbe \emph{to profit} signifie « tirer un profit financier » ; pour « bénéficier de », on dit \emph{benefit from} ou \emph{take advantage of}.
\item \textbf{Correction :} \eng{He is \cle{responsible for} the cleaning team.} (« Il est responsable de l'équipe de nettoyage. ») \quad \emph{Erreur :} calque syntaxique. L'adjectif \emph{responsible} exige la préposition \textbf{for} (jamais \emph{of}).
\end{enumerate}
\end{corrige}

\begin{corrige}[Exercice 4 — Formation de mots et grammaire contextuelle]
\begin{enumerate}
\item \cle{participation} — nom (action), placé après l'article \emph{The} et sujet de \emph{is}. \eng{The participation of every student is compulsory for the sports day.} (« La participation de chaque élève est obligatoire pour la journée sportive. »)
\item \cle{punctual} — adjectif, après \emph{very}, qualifiant \emph{person}. \eng{She is a very punctual person; she is never late.}
\item \cle{organiser} — nom d'agent, sujet de \emph{did} ; orthographe britannique en \textbf{-iser} (norme au Cameroun). \eng{The organiser of the event did a great job.}
\item \cle{do not manage} / \cle{don't manage} — présent simple dans la subordonnée introduite par \emph{if} (conditionnel de type 1 : \emph{if + present} $\rightarrow$ \emph{will} + base verbale) ; \cle{deadline} — nom, après l'article \emph{the}. \eng{If we don't manage our time well, we will miss the deadline.}
\item \cle{participants} — nom d'agent au pluriel, après \emph{All}, sujet de \emph{must register}. \eng{All participants must register before Wednesday.}
\end{enumerate}
\textbf{Point de vigilance :} dans une conditionnelle avec \emph{if}, on n'emploie jamais \emph{will} dans la proposition introduite par \emph{if} (jamais \emph{if we will not manage}).
\end{corrige}

\begin{corrige}[Exercice 5 — Traduction : Du français vers l'anglais]
\begin{enumerate}
\item \eng{The \cle{deadline} for \cle{registration} (ou \cle{signing up}) is set for May 15th.} — \emph{La date limite pour} se rend par \emph{the deadline for} + nom ou gérondif ; la date se place sans article : \emph{May 15th}.
\item \eng{The \cle{volunteers} worked \cle{around the clock} to prepare the hall (ou \cle{venue}).} — \emph{Jour et nuit} se traduit idiomatiquement par \emph{around the clock}, jamais par un calque du type \emph{day and night} dans ce contexte soutenu (bien que \emph{day and night} soit acceptable).
\item \eng{We \cle{had to postpone} the match \cle{due to} (ou \cle{because of}) the rain.} — \emph{Avoir dû} = \emph{had to} (passé de \emph{have to}) ; \emph{à cause de} + nom = \emph{due to / because of}.
\item \eng{Every student must \cle{pull their weight} with the household chores.} — L'expression idiomatique \emph{pull one's weight} remplace avantageusement le calque \emph{participate to the household chores} ; notez \emph{their} (singulier neutre) après \emph{every student}.
\item \eng{The \cle{coordinator} thanked everyone for their \cle{commitment}.} — \emph{Remercier qqn pour qch} = \emph{thank someone for} + nom/gérondif.
\end{enumerate}
\end{corrige}

\begin{corrige}[Exercice 6 — Transformation de phrases]
\begin{enumerate}
\item \eng{Due to \cle{the strike, the meeting was postponed}.} — \emph{Due to} + nominal reprend exactement la cause exprimée par \emph{because of} ; l'ordre des propositions est inversé.
\item \eng{The deadline \cle{for submitting your report is Friday}.} (ou \eng{The deadline for your report is Friday.}) — Structure : \emph{the deadline for} + gérondif (\emph{submitting}) ou nom ; \emph{before Friday} devient simplement \emph{is Friday}.
\item \eng{They worked around \cle{the clock}.} — \emph{All night} (toute la nuit) est reformulé par l'idiome \emph{around the clock} (sans arrêt).
\item \eng{He is never \cle{late}.} — Équivalence : \emph{always on time} = \emph{never late} (négation + antonyme).
\item \eng{She asked me to \cle{help her carry those boxes}.} — Discours rapporté : demande $\rightarrow$ \emph{ask + objet + to}-infinitif ; \emph{me} reste \emph{me} (ici l'interlocuteur est le rapporteur) ; \emph{these} $\rightarrow$ \emph{those} ; \emph{can} disparaît devant l'infinitif.
\end{enumerate}
\textbf{Point de vigilance :} ne jamais écrire \emph{She asked me to helping} ni conserver \emph{these} dans le discours rapporté.
\end{corrige}

\begin{corrige}[Exercice 7 — Texte à trous : Banque de mots]
Les deux mots en trop sont \textbf{cancelled} et \textbf{punctual}.
\begin{enumerate}
\item \cle{schedule} — \eng{The schedule for the annual fundraising day was published last week.} (le planning a été publié)
\item \cle{coordinator} — nom d'agent, précédé de l'adjectif \emph{main} : \eng{The main coordinator reminded everyone...}
\item \cle{deadline} — \eng{...that the deadline for registration is Tuesday.} (structure \emph{deadline for} + nom)
\item \cle{volunteers} — \eng{Over fifty volunteers have already signed up...} (pluriel après \emph{fifty})
\item \cle{take part} — après \emph{to} (infinitif de but) : \eng{...signed up to take part in the various activities.}
\item \cle{postponed} — passif (\emph{was} + participe passé) : \eng{...the football tournament was postponed due to a waterlogged pitch.} (\emph{cancelled} est impossible : on espère une \emph{new date}, donc le tournoi est reporté, pas annulé.)
\item \cle{on time} — \eng{...will allow everyone to arrive on time.} (\emph{punctual} est un adjectif ; il faudrait un adverbe ou une locution pour modifier \emph{arrive}.)
\item \cle{lend a hand} — impératif : \eng{Please lend a hand wherever you can!} (donnez un coup de main !)
\end{enumerate}
\end{corrige}

\begin{corrige}[Exercice 8 — Appariement : Expressions idiomatiques]
\begin{center}
\begin{tabularx}{\linewidth}{|X|X|X|}
\hline
\rowcolor{popblueL} \textbf{Expression} & \textbf{Réponse} & \textbf{Signification} \\
\hline
1. In the nick of time & \textbf{C} & Juste à temps / au dernier moment \\
2. Around the clock & \textbf{A} & Travailler sans relâche / 24h sur 24 \\
3. Pull one's weight & \textbf{B} & Faire sa part du travail \\
4. Beat the clock & \textbf{F} & Faire quelque chose le plus vite possible avant une échéance \\
5. Against the clock & \textbf{H} & Essayer de finir avant l'heure limite \\
6. Time is money & \textbf{E} & Le temps, c'est de l'argent \\
7. A race against time & \textbf{G} & Une course contre la montre \\
8. On the dot & \textbf{D} & Précisément à l'heure dite \\
\hline
\end{tabularx}
\end{center}
\textbf{Exemples :} \eng{The meeting starts at 8 o'clock on the dot.} (« La réunion commence à 8 heures précises. ») ; \eng{We are working against the clock to finish before Friday.} (« Nous travaillons contre la montre pour finir avant vendredi. »)
\end{corrige}

\begin{corrige}[Exercice 9 — Compréhension et langue : type Bac]
\textbf{Barème indicatif (sur 20) :} A. Compréhension : 8 pts (2 pts par question) ; B. Vocabulaire : 4 pts (1 pt par mot) ; C. Grammaire / Formation de mots : 8 pts (passif 2 pts ; noms 3 pts ; prépositions 3 pts).

\textbf{A. Comprehension Questions}
\begin{enumerate}
\item \eng{The summit took place at the University of Buea, from the 10th to the 12th of August.} (« Le sommet a eu lieu à l'Université de Buea, du 10 au 12 août. »)
\item \eng{Her main message was that punctuality and time management are essential in business: if you cannot manage your time, you cannot manage a business.} (« Son message principal : la ponctualité et la gestion du temps sont essentielles dans les affaires. »)
\item \eng{It was a contest in which each contestant had exactly sixty seconds to present a project and convince the jury.} (« Un concours où chaque candidat disposait d'exactement 60 secondes pour présenter son projet et convaincre le jury. »)
\item \eng{He received a grant of 500,000 FCFA and urged his peers to seize opportunities and meet deadlines.} (« Il a reçu une subvention de 500\,000 FCFA et a exhorté ses camarades à saisir les opportunités et à respecter les échéances. »)
\end{enumerate}

\textbf{B. Vocabulary in Context}
\begin{enumerate}
\item \cle{participants} (paragraphe 1)
\item \cle{simultaneously} (paragraphe 2)
\item \cle{seize opportunities} (paragraphe 3)
\item \cle{mentorship} (paragraphe 3)
\end{enumerate}

\textbf{C. Grammar / Word Formation}
\begin{enumerate}
\item \eng{The summit \cle{was hosted} at the University of Buea \cle{by the organisers}.} — Passif au prétérit : \emph{was} + participe passé \emph{hosted} ; l'agent est introduit par \emph{by}.
\item \emph{punctual} $\rightarrow$ \cle{punctuality} ; \emph{participate} $\rightarrow$ \cle{participation} ; \emph{organise} $\rightarrow$ \cle{organisation} (orthographe britannique).
\item \eng{The workshop focused \cle{on} digital tools.} / \eng{He insisted \cle{on} the importance of time.} — \emph{Focus on} et \emph{insist on} : préposition \textbf{on} obligatoire dans les deux cas.
\end{enumerate}
\textbf{Points de vigilance :} ne pas écrire \emph{insist about} ni \emph{focus in} ; au passif, ne pas oublier l'auxiliaire \emph{was} (jamais \emph{the summit hosted by...} seul).
\end{corrige}

\begin{corrige}[Exercice 10 — Traduction : Français $\rightarrow$ Anglais]
\textbf{Barème indicatif (sur 10) :} respect du sens 4 pts ; lexique de la leçon 3 pts ; grammaire (structures, inversion finale) 2 pts ; orthographe et ponctuation 1 pt.

\textbf{Traduction modèle :}
\eng{“Organising a school event requires good planning. The coordinator must draw up a precise schedule and set a deadline for registration. Volunteers must arrive on time and work around the clock so that everything is ready. If an activity is postponed, everyone must be informed quickly. Only in this way will the event be a success.”}

\textbf{Notes de traduction :}
\begin{itemize}
\item \emph{L'organisation d'un événement} $\rightarrow$ \cle{Organising a school event} (gérondif sujet, plus naturel que \emph{The organisation of...}).
\item \emph{Établir un emploi du temps} $\rightarrow$ \cle{draw up a schedule} (collocation correcte ; éviter \emph{do a schedule}).
\item \emph{Fixer une date limite} $\rightarrow$ \cle{set a deadline}.
\item \emph{Travailler sans relâche} $\rightarrow$ \cle{work around the clock}.
\item \emph{Est reportée} $\rightarrow$ passif : \cle{is postponed}.
\item \emph{C'est seulement ainsi que...} $\rightarrow$ \cle{Only in this way will the event be...} : inversion sujet--auxiliaire obligatoire après \emph{Only} + complément circonstanciel en tête de phrase.
\end{itemize}
\textbf{Points de vigilance :} ne pas traduire \emph{demande} par \emph{demands} (faux ami : \emph{demand} = exiger) ; préférer \emph{requires}. Ne pas écrire \emph{in the hour} pour « à l'heure » : dire \emph{on time}.
\end{corrige}

\begin{corrige}[Exercice 11 — Expression écrite : Sujets de type Bac]
\textbf{Barème indicatif (sur 20) :} respect de la consigne et du format (article ou email) 4 pts ; richesse et exactitude du lexique de la leçon 6 pts ; correction grammaticale 5 pts ; cohérence et organisation 3 pts ; orthographe et ponctuation 2 pts.

\textbf{Modèle Sujet 1 (Article de magazine scolaire) :}
\eng{“JOIN US FOR A CLEANER CAMPUS THIS SATURDAY!}
\eng{By the President of the Environment Club}
\eng{Next Saturday, 25th May, the Environment Club is launching its “Clean-Up Campaign”. We invite every student to \textbf{take part} in making our school shine. The \textbf{deadline} for signing up is Wednesday 22nd May, at the club office.}
\eng{The \textbf{schedule} is simple: gathering at 7:30 am sharp at the main gate — \textbf{be on time}! — then distribution of gloves and bags, and cleaning of the assigned zones: playground, gardens and back streets. Refreshments will be provided at 10:30 am.}
\eng{This is our chance to \textbf{pull our weight} for the community. A clean school means a healthier learning environment for everyone. Do not wait until the last minute — register now and show your \textbf{commitment}!”}

\textbf{Modèle Sujet 2 (Email formel) :}
\eng{Subject: Request for Permission to Organise a Career Day — Upper Sixth A}
\eng{Dear Principal,}
\eng{I am writing on behalf of the Upper Sixth A class to request your authorisation to organise a “Career Day” for final-year students. The main objective is to guide students in their professional orientation by allowing them to meet professionals from various sectors, such as medicine, engineering, finance and entrepreneurship.}
\eng{We propose Saturday, 15th June, from 9:00 am to 1:00 pm. A draft \textbf{schedule} includes: 9:00 am — opening ceremony; 9:30 am — panel discussions; 11:00 am — question-and-answer sessions; 12:30 pm — closing remarks.}
\eng{We would be grateful for the school's support: the use of the main hall as a \textbf{venue}, assistance in contacting invited professionals, and permission for students to be excused from classes that morning. The \textbf{deadline} for confirming the speakers is 1st June.}
\eng{We remain at your disposal for any further information.}
\eng{Yours faithfully,}
\eng{[Name], Class Prefect, Upper Sixth A}

\textbf{Points de vigilance :}
\begin{itemize}
\item Article : titre accrocheur, signature, paragraphes courts, ton motivant ; employer l'impératif (\emph{register now!}).
\item Email formel : objet (\emph{Subject}), formule d'appel (\emph{Dear Principal}), formule de politesse finale (\emph{Yours faithfully}), pas de contractions (\emph{we would}, pas \emph{we'd}).
\item Réemployer au moins cinq mots-cibles : \cle{deadline}, \cle{schedule}, \cle{take part}, \cle{on time}, \cle{commitment}, \cle{volunteers}, \cle{venue}.
\item Éviter les calques : « prendre part » = \emph{take part in} (jamais \emph{take part to}) ; « à l'heure » = \emph{on time} (jamais \emph{in time} pour la ponctualité simple).
\end{itemize}

\textbf{Activité orale guidée — Pistes de réalisation (jeu de rôle) :}
\begin{itemize}
\item \textbf{Élève A (coordinator) :} \eng{“Thank you for volunteering! Can you tell me when you are available? Here is the schedule for the Cultural Day. The deadline for registration is Friday. Please be on time — we start at 8 o'clock sharp.”}
\item \textbf{Élève B (volunteer) :} \eng{“I am free on Saturday morning. What tasks can I take part in? I would be happy to lend a hand with the decorations. You can count on my commitment — I will pull my weight!”}
\end{itemize}
Critères d'évaluation orale : prononciation claire (noter \emph{schedule} : « ché-dioul » en britannique, « ské-dioul » en américain), usage d'au moins cinq expressions de la leçon, fluidité et interaction naturelle.
\end{corrige}

\newpage

\coursec{Fiche bilan}

\coursec{Lesson 23 -- Vocabulary: Words and expressions related to taking part and managing time in community/workplace/school activities}

\begin{textebox}[Staff Meeting: Preparing the Annual Prize-Giving Day]
Good morning everyone. I \cle{call this meeting to order} at 10:00 a.m. Thank you for \cle{attending} on such short notice. As you know, the Prize-Giving Day is a major \cle{milestone} in our school calendar. The \cle{deadline} for the final programme is next Friday. We are currently \cle{behind schedule} on the printing of certificates.

Mr. Fouda, as \cle{secretary}, please read the \cle{minutes} of the last meeting.
[Minutes read] \dots \cle{Adopted}.

Item 2 on the \cle{agenda}: \cle{Allocation} of tasks. We need to \cle{prioritize} the venue decoration and the guest list. I propose we \cle{allocate} two teachers per committee. Mrs. Mendo, would you \cle{chair} the logistics committee? Mr. Bello, can you \cle{facilitate} the sponsorship drive?

We need a \cle{quorum} of 10 for the vote on the budget. Do we have a quorum? Yes.
Let's not \cle{delay} the decision on the date. I suggest we \cle{postpone} the cultural rehearsals to next week to \cle{buy time} for the academic awards verification.

Any \cle{Other Business}? No? Then I declare the meeting \cle{adjourned} at 11:30. Next meeting: Monday, same \cle{timeslot}. Thank you.
\end{textebox}

\begin{popfigure}
\centering
\begin{tikzpicture}[
    mindmap,
    concept color=popblue,
    level 1/.append style={level distance=4.5cm, sibling angle=360/7},
    level 2/.append style={level distance=3cm},
    every node/.style={font=\small, align=center},
    branch1/.style={concept color=popblue, grow=90},
    branch2/.style={concept color=poporange, grow=140},
    branch3/.style={concept color=popgreen, grow=190},
    branch4/.style={concept color=poppurple, grow=240},
    branch5/.style={concept color=poppink, grow=290},
    branch6/.style={concept color=popgold, grow=340},
    branch7/.style={concept color=popdark, grow=40}
]
\node[concept, text=white, scale=1.2] {Taking Part \&\ Managing Time}
    child[branch1] { node[concept, text=white] {Participation\ Verbs}
        child { node[rectangle, rounded corners, fill=popblueL, draw=popblue, inner sep=3pt] {attend, join, take part} }
        child { node[rectangle, rounded corners, fill=popblueL, draw=popblue, inner sep=3pt] {contribute, volunteer, engage} }
        child { node[rectangle, rounded corners, fill=popblueL, draw=popblue, inner sep=3pt] {chair, facilitate, second} }
    }
    child[branch2] { node[concept, text=white] {Time Management\ Verbs}
        child { node[rectangle, rounded corners, fill=poporangeL, draw=poporange, inner sep=3pt] {schedule, plan, prioritize} }
        child { node[rectangle, rounded corners, fill=poporangeL, draw=poporange, inner sep=3pt] {meet deadlines, allocate, delay} }
        child { node[rectangle, rounded corners, fill=poporangeL, draw=poporange, inner sep=3pt] {postpone, reschedule, streamline} }
    }
    child[branch3] { node[concept, text=white] {Key Nouns\&\ Roles}
        child { node[rectangle, rounded corners, fill=popgreenL, draw=popgreen, inner sep=3pt] {agenda, minutes, quorum} }
        child { node[rectangle, rounded corners, fill=popgreenL, draw=popgreen, inner sep=3pt] {chairperson, secretary, stakeholder} }
        child { node[rectangle, rounded corners, fill=popgreenL, draw=popgreen, inner sep=3pt] {deadline, timeline, milestone} }
    }
    child[branch4] { node[concept, text=white] {Collocations\&\ Idioms}
        child { node[rectangle, rounded corners, fill=poppurpleL, draw=poppurple, inner sep=3pt] {take the floor, run a meeting} }
        child { node[rectangle, rounded corners, fill=poppurpleL, draw=poppurple, inner sep=3pt] {time is money, behind schedule} }
        child { node[rectangle, rounded corners, fill=poppurpleL, draw=poppurple, inner sep=3pt] {on the agenda, call to order} }
    }
    child[branch5] { node[concept, text=white] {False Friends\&\ Calques}
        child { node[rectangle, rounded corners, fill=poppinkL, draw=poppink, inner sep=3pt] {assist $\neq$ assister (help)} }
        child { node[rectangle, rounded corners, fill=poppinkL, draw=poppink, inner sep=3pt] {attend $\neq$ attendre (be present)} }
        child { node[rectangle, rounded corners, fill=poppinkL, draw=poppink, inner sep=3pt] {preside $\neq$ présider (chair)} }
    }
    child[branch6] { node[concept, text=white] {Word Families}
        child { node[rectangle, rounded corners, fill=popgoldL, draw=popgold, inner sep=3pt] {participate $\to$ participation} }
        child { node[rectangle, rounded corners, fill=popgoldL, draw=popgold, inner sep=3pt] {organize $\to$ organization} }
        child { node[rectangle, rounded corners, fill=popgoldL, draw=popgold, inner sep=3pt] {manage $\to$ management} }
    }
    child[branch7] { node[concept, text=white] {Contexts:\ School, Work, Community}
        child { node[rectangle, rounded corners, fill=popdark!20, draw=popdark, inner sep=3pt] {PTA meetings, staff meetings} }
        child { node[rectangle, rounded corners, fill=popdark!20, draw=popdark, inner sep=3pt] {community service, volunteering} }
        child { node[rectangle, rounded corners, fill=popdark!20, draw=popdark, inner sep=3pt] {project planning, shift roster} }
    };
\end{tikzpicture}
\legende{Carte mentale : Vocabulaire de la participation et de la gestion du temps (Leçon 23)}
\end{popfigure}

\begin{definition}[Lexique clé : Verbes d'action]
\begin{center}
\begin{tabularx}{\linewidth}{|X|X|X|}\hline
\rowcolor{popblueL} \textbf{Verbe (EN)} & \textbf{Traduction (FR)} & \textbf{Contexte / Collocation} \\ \hline
\eng{attend} & assister à, être présent à & \eng{attend a meeting / school regularly} \\ \hline
\eng{chair} / \eng{preside over} & présider, diriger (une réunion) & \eng{chair the session, preside over the AGM} \\ \hline
\eng{convene} & convoquer, réunir & \eng{convene an emergency meeting} \\ \hline
\eng{facilitate} & faciliter, animer & \eng{facilitate a workshop / discussion} \\ \hline
\eng{second a motion} & appuyer une motion & \eng{I second the proposal.} \\ \hline
\eng{allocate} & allouer, affecter (temps, ressources) & \eng{allocate time for questions} \\ \hline
\eng{prioritize} & prioriser, classer par ordre d'importance & \eng{prioritize urgent tasks} \\ \hline
\eng{meet a deadline} & respecter un délai & \eng{We must meet the Friday deadline.} \\ \hline
\eng{postpone} / \eng{defer} & reporter, ajourner & \eng{postpone the meeting to Monday} \\ \hline
\eng{streamline} & rationaliser, simplifier & \eng{streamline the workflow} \\ \hline
\eng{volunteer} & se porter volontaire & \eng{volunteer for a task} \\ \hline
\eng{contribute} & contribuer, apporter sa pierre & \eng{contribute to the discussion} \\ \hline
\eng{engage} & s'engager, impliquer & \eng{engage stakeholders} \\ \hline
\eng{schedule} & planifier, programmer & \eng{schedule a meeting for 3 p.m.} \\ \hline
\eng{minute} (v) & rédiger le procès-verbal & \eng{The secretary minuted the decisions.} \\ \hline
\end{tabularx}
\end{center}
\end{definition}

\begin{definition}[Lexique clé : Noms, rôles et unités de temps]
\begin{center}
\begin{tabularx}{\linewidth}{|X|X|X|}\hline
\rowcolor{poporangeL} \textbf{Terme (EN)} & \textbf{Traduction (FR)} & \textbf{Note / Exemple} \\ \hline
\eng{agenda} & ordre du jour & \eng{Item 3 on the agenda} \\ \hline
\eng{minutes} (pl.) & procès-verbal, compte rendu & \eng{read the minutes of the last meeting} \\ \hline
\eng{quorum} & quorum (nombre min. requis) & \eng{We have a quorum.} \\ \hline
\eng{chairperson} / \eng{chair} & président(e) de séance & \eng{address the chair} \\ \hline
\eng{secretary} & secrétaire & \eng{take the minutes} \\ \hline
\eng{stakeholder} & partie prenante & \eng{key stakeholders} \\ \hline
\eng{deadline} & date limite, échéance & \eng{tight deadline, extend the deadline} \\ \hline
\eng{milestone} & étape clé, jalon & \eng{reach a milestone} \\ \hline
\eng{timeslot} / \eng{slot} & créneau horaire & \eng{book a 30-minute slot} \\ \hline
\eng{roster} / \eng{rota} & planning, tableau de service & \eng{the duty roster for the week} \\ \hline
\eng{timeline} & chronogramme, calendrier & \eng{project timeline} \\ \hline
\eng{action point} / \eng{action item} & point d'action, décision à exécuter & \eng{Action point: Mr. Bello to contact sponsors} \\ \hline
\end{tabularx}
\end{center}
\end{definition}

\begin{definition}[Collocations \& Expressions idiomatiques (à connaître par cœur)]
\begin{center}
\begin{tabularx}{\linewidth}{|X|X|}\hline
\rowcolor{poppurpleL} \textbf{Expression anglaise} & \textbf{Équivalent français / Sens} \\ \hline
\eng{Take the floor} & Prendre la parole (en réunion) \\ \hline
\eng{Call the meeting to order} & Ouvrir la séance \\ \hline
\eng{Table a motion} & Mettre une motion aux voix (UK) / Ajourner (US) \\ \hline
\eng{Run behind schedule} & Avoir du retard sur l'horaire \\ \hline
\eng{Press for time} / \eng{Short on time} & Être pressé par le temps \\ \hline
\eng{Time is money} & Le temps, c'est de l'argent \\ \hline
\eng{At the eleventh hour} & Au dernier moment \\ \hline
\eng{Buy time} & Gagner du temps \\ \hline
\eng{On the agenda} & À l'ordre du jour \\ \hline
\eng{Action points} / \eng{Action items} & Points d'action, décisions à exécuter \\ \hline
\eng{Stick to the agenda} & Respecter l'ordre du jour \\ \hline
\eng{Run out of time} & Manquer de temps \\ \hline
\end{tabularx}
\end{center}
\end{definition}

\begin{attention}[Faux amis \& Calques fréquents (Pièges pour francophones)]
\begin{center}
\begin{tabularx}{\linewidth}{|X|X|X|}\hline
\rowcolor{poppinkL} \textbf{Mot anglais} & \textbf{Piège (FR)} & \textbf{Vrai sens / Correction} \\ \hline
\eng{Assist} & \eng{Assister à} (être présent) & \textbf{Help / Support.} \eng{Assist a colleague.} Pour « assister à » : \eng{attend}. \\ \hline
\eng{Attend} & \eng{Attendre} & \textbf{Être présent à.} \eng{Attend a meeting.} Pour « attendre » : \eng{wait for}. \\ \hline
\eng{Preside} & \eng{Présider} (semble OK mais registre) & Souvent remplacé par \textbf{\eng{chair}} (plus courant). \eng{Preside} = très formel. \\ \hline
\eng{Intervene} & \eng{Intervenir} (venir réparer) & \textbf{Prendre la parole / S'interposer.} \eng{Intervene in a debate.} \\ \hline
\eng{Report} & \eng{Rapporter} (apporter) & \textbf{Faire un compte rendu / Signaler.} \eng{Report the minutes.} \\ \hline
\eng{Delay} (n/v) & \eng{Délai} (n) & \textbf{Retard.} \eng{Delay} = retard. Pour « délai » : \eng{deadline} ou \eng{time limit}. \\ \hline
\eng{Agenda} & \eng{Agenda} (carnet) & \textbf{Ordre du jour.} Carnet = \eng{diary} (UK) / \eng{planner} (US). \\ \hline
\eng{Minute} (n) & \eng{Minute} (temps) & \textbf{Procès-verbal (souvent pl. \eng{minutes}).} Temps = \eng{minute} (60 sec). \\ \hline
\end{tabularx}
\end{center}
\end{attention}

\begin{propriete}[Familles de mots (Dérivation) -- Orthographe \& Prononciation]
\begin{center}
\begin{tabular}{|c|c|c|c|}\hline
\rowcolor{popgreenL} \textbf{Verbe} & \textbf{Nom (action/état)} & \textbf{Nom (agent)} & \textbf{Adjectif} \\ \hline
\eng{participate} & \eng{participation} & \eng{participant} & \eng{participatory} \\ \hline
\eng{organize} & \eng{organization} & \eng{organizer} & \eng{organized} \\ \hline
\eng{manage} & \eng{management} & \eng{manager} & \eng{managerial} \\ \hline
\eng{contribute} & \eng{contribution} & \eng{contributor} & \eng{contributory} \\ \hline
\eng{schedule} & \eng{schedule} / \eng{scheduling} & \eng{scheduler} & \eng{scheduled} \\ \hline
\eng{volunteer} & \eng{volunteering} & \eng{volunteer} & \eng{voluntary} \\ \hline
\end{tabular}
\end{center}

\textbf{Prononciation repères :} \eng{schedule} $\to$ « chèd-youl » (UK) / « skèd-joul » (US) ; \eng{agenda} $\to$ « a-djèn-da » ; \eng{quorum} $\to$ « kwo-ram » ; \eng{minutes} $\to$ « mi-nits » (pas « mi-nutes ») ; \eng{prioritize} $\to$ « pra-io-ri-taïz » ; \eng{contributory} $\to$ « kon-tri-byou-t'ri ».
\end{propriete}

\begin{methode}[Réviser le lexique efficacement]
\begin{enumerate}
\item \textbf{Cartes mentales par contexte :} Crée 3 cartes (École, Travail, Communauté) avec verbes/noms/expressions spécifiques.
\item \textbf{Tableau de collocations :} Note \textbf{Verbe + Nom} (ex. \eng{chair a meeting}, \eng{take minutes}, \eng{meet a deadline}). Apprends les blocs, pas les mots isolés.
\item \textbf{Chasse aux faux amis :} Liste les 8 faux amis du tableau ci-dessus ; écris une phrase correcte pour chacun.
\item \textbf{Simulation orale :} Enregistre-toi en train de « présider » une mini-réunion de 2 min (ordre du jour, tour de parole, clôture).
\end{enumerate}
\end{methode}

\begin{methode}[Rédiger un compte-rendu (Minutes) -- Structure type]
\begin{enumerate}
\item \textbf{Header :} Organisation, Date, Time, Venue, Attendees / Apologies (absents excusés).
\item \textbf{Opening :} \eng{The meeting was called to order at [time] by [Chair].}
\item \textbf{Minutes of previous meeting :} \eng{Adopted / Amended.}
\item \textbf{Matters arising / Agenda items :} Sujet $\to$ Discussion $\to$ \textbf{Decision / Action point} (Who? What? When?).
\item \textbf{AOB (Any Other Business) :} Questions diverses.
\item \textbf{Closing :} \eng{The meeting adjourned at [time]. Next meeting: [date].}
\item \textbf{Signature :} \eng{Chairperson / Secretary.}
\end{enumerate}
\end{methode}

\begin{aretenir}[Checklist avant le Bac]
$\square$ Je connais par cœur les \textbf{15 verbes clés} (\eng{attend, chair, convene, facilitate, second, allocate, prioritize, meet a deadline, postpone, streamline, volunteer, contribute, engage, schedule, minute}).\par
$\square$ Je sais traduire et employer les \textbf{12 noms essentiels} (\eng{agenda, minutes, quorum, chairperson, secretary, stakeholder, deadline, milestone, timeslot, roster, timeline, action point}).\par
$\square$ J'utilise correctement \textbf{10 collocations/idiomes} (\eng{take the floor, call to order, behind schedule, press for time, at the eleventh hour, buy time, on the agenda, action points, table a motion, stick to the agenda}).\par
$\square$ J'évite les \textbf{5 faux amis majeurs} : \eng{assist} (aider), \eng{attend} (être présent), \eng{delay} (retard), \eng{agenda} (ordre du jour), \eng{minutes} (PV).\par
$\square$ Je forme les \textbf{familles de mots} sans faute d'orthographe (\eng{participate/participation}, \eng{organize/organization}, \eng{manage/management}, \eng{contribute/contribution/contributory}).\par
$\square$ Je prononce correctement \eng{schedule, agenda, quorum, minutes, prioritize, contributory}.\par
$\square$ Je sais structurer un \textbf{compte-rendu formel (minutes)} avec les rubriques standards.\par
$\square$ Je sais \textbf{présider une réunion} (ouvrir, donner la parole, voter, clore) en anglais oral.\par
$\square$ Je gère le vocabulaire de la \textbf{planification} (\eng{roster, timeline, allocate, slot}) pour un stage ou job d'été.\par
$\square$ Je distingue \textbf{UK / US} : \eng{diary/planner, timetable/schedule, postponing/tabling a motion}.\par
$\square$ Je suis capable de \textbf{répondre à un sujet d'expression écrite} (lettre formelle, article, discours) en intégrant ce lexique naturellement.
\end{aretenir}

\findecours

\end{document}
